% ============================================================
%  H. Høffding, Religionsfilosofi
%  TRANSLATION (English)
%
%  EDITION.  This translation follows the THIRD EDITION, whose
%  preface Høffding signed "Carlsberg d. 28. Februar 1924".
%  The book was first published 1901 (2nd ed. 1906); the
%  catalog entry is filed under 1901 for that reason, but the
%  text translated here is the 1924 revision, which is
%  Høffding's own final version.
%
%  In the Fortale til tredie Udgave Høffding records what he
%  changed:
%    - the epistemological section (ch. II) was SHORTENED,
%      "kun det strengt nødvendige" retained, because he had
%      since treated the matter in Den menneskelige Tanke
%      (1910), Totalitet (1917), Relation (1921), Analogi (1923);
%    - the two groupings of religious-psychological types were
%      MERGED ("en Sammenarbejdelse af de to ... Grupperinger").
%      Accordingly the 1901 subsection "Hovedgrupper af
%      personlige Forskelligheder" (III.E.b) is ABSENT from the
%      3rd ed. and has been dropped from this translation;
%    - a number of notes were DELETED as outdated or
%      superfluous, and new ones added.
%
%  SOURCE.  Translated from the Lindhardt og Ringhof e-book
%  (ISBN 9788726277555, 2020), which sets the 3rd edition.
%  Per the repo rule, the Danish e-book text is NOT transcribed,
%  committed, or published: the epub is gitignored and only this
%  English translation goes on the site.  The partial
%  transcription.tex in this directory is keyed to the 1901
%  text and is no longer the source of truth for the translation;
%  it is retained but unlinked from the catalog.
%
%  Cf. B. E. Meyer's 1906 translation (The Philosophy of
%  Religion, London: Macmillan), which renders the FIRST edition
%  by way of the German, and so differs from this text both in
%  substance and in what it contains.
% ============================================================
\documentclass[12pt,a4paper]{book}
\usepackage[utf8]{inputenc}
\usepackage[T1]{fontenc}
\usepackage{libertinus}
\usepackage{libertinust1math}
\usepackage{textalpha}
\usepackage[english]{babel}
\usepackage{enumitem}
\usepackage[protrusion=true,expansion=true]{microtype}
\usepackage{setspace}
\usepackage[margin=1.2in]{geometry}
\usepackage{fancyhdr}
\setlength{\headheight}{14.5pt}
\addtolength{\topmargin}{-2.5pt}
\usepackage{hyperref}

\hypersetup{
  pdftitle={Philosophy of Religion},
  pdfauthor={Harald Høffding},
  hidelinks
}

\pagestyle{fancy}
\fancyhf{}
\fancyhead[LE,RO]{\thepage}
\fancyhead[RE]{\textit{Philosophy of Religion}}
\fancyhead[LO]{\textit{\leftmark}}

\setstretch{1.3}

\title{\textbf{Philosophy of Religion}\\[6pt]
  \large\textit{Religionsfilosofi}}
\author{Harald Høffding\\[4pt]
  \small Translated from the Danish}
\date{Copenhagen and Christiania:\\
  Gyldendalske Boghandel, Nordisk Forlag\\[6pt]
  \small First edition 1901; second edition 1906.\\
  \small \textit{Translated from the third edition of 1924.}}

\begin{document}
\maketitle
\thispagestyle{empty}

\vspace*{\fill}
\begin{center}
\textit{Alte dubitat, qui altius credit.}\\[6pt]
(He doubts more deeply who believes more deeply.)
\end{center}
\vspace*{\fill}
\newpage

\tableofcontents
\newpage

%% ============================================================
%%  PREFACE TO THE THIRD EDITION
%% ============================================================
\chapter*{Preface to the Third Edition}
\addcontentsline{toc}{chapter}{Preface to the Third Edition}
\label{ch:fortale3}
\markboth{Preface to the Third Edition}{}

\noindent When I worked out my \emph{Philosophy of Religion} in the
years 1896--1901, I had not yet given any fuller account of
epistemology. Nor had I arrived at one when the second edition was due
to appear (1906). Later I set out my epistemological position in ``The
Human Thought'' (1910) and in a series of investigations bearing on the
subject of the present book (on Totality 1917, on Relation 1921, and on
Analogy 1923). It has thereby been possible, in this third edition, to
shorten the epistemological section; on several points I could now
simply refer to the works named, which would otherwise have required
detailed justification. Only the strictly necessary epistemological
content has been retained. In the religious-psychological section too I
was able to refer to works that threw light on the topics (``The Great
Humour'' 1916; ``Experiences and Interpretations'' 1918; ``Leading
Thoughts in the Nineteenth Century'' 1920; together with the essay
``Foreword and Postscript to my Philosophy of Religion,'' printed in
``Minor Works''). I have besides carried out a consolidation of the two
groupings of religious-psychological types, determined by different
points of view, between which I earlier thought it possible to
distinguish. Of the notes, a number have been omitted as outdated or
superfluous. Where it was necessary, I have added new notes.

The book's fundamental thought has remained the same. Its significance
has confirmed itself to me over the years by what I have experienced or
learned. Often it is precisely objections that have led to fresh
corroboration. The book's fundamental thought is not a key that unlocks
everywhere, but it leads one to pose fruitful questions and to enter
upon investigations that one would otherwise be inclined to think
superfluous. It warns against dogmatism of every kind, and it points at
the same time to a deep-lying human need that can hardly die out so long
as the life of the spirit is not being extinguished. It does not solve
the great problem it treats, but it can lead to ever-renewed reflection.
And I hope that the book will continue to work in the direction here
indicated.

\bigskip

\begin{flushright}
Carlsberg, 28 February 1924.\\
\textsc{Harald Høffding}.
\end{flushright}

\newpage

%% ============================================================
%%  I. PROBLEM AND PROCEDURE
%% ============================================================
\chapter{Problem and Procedure}
\label{ch:problem}

% The e-book sets this epigraph without word spaces and with corrupted
% breathings/accents ("Οὕϰἑστιλὐεινἀγνοοῦντατò δεσμόν"), plainly a defect of
% the 2020 digitisation rather than a reading of the print. Word division and
% diacritics restored here; the epub's form is recorded above. Aristotle,
% Metaphysics III.1, 995a29.
\begin{flushright}
Οὐκ ἔστι λύειν ἀγνοοῦντα τὸ δεσμόν.\\
\textsc{Aristotle}.\\[4pt]
\textit{(One cannot loose a knot without knowing where the binding lies).}
\end{flushright}

\bigskip

\noindent\textbf{1.} By \emph{philosophy of religion} one could mean
either a thinking that springs from religion and builds on religion, or a
thinking that takes religion as its object. It is in this latter sense
that the term shall be taken here. The thinking that springs from and
builds on religion itself belongs to the subject matter of the philosophy
of religion and should not be called philosophy of religion.

By religion we think first of all of a state of mind in which feeling
and need, fear and hope, enthusiasm and devotion play a larger role than
reflection and inquiry, and in which intuition and image hold more sway
than investigation and reasoning. To be sure, again and again in the
religious life there will arise an involuntary need to give an account
of one's condition and of its content---of the worth of one's motives
and the validity of one's thoughts. But the religious thinking that thus
arises does not pose the religious problem in a decisive way. Its
problems arise within religion, but religion itself does not become a
problem. No other starting point is yet recognized than religion
itself; in any case other starting points are so subordinate in relation
to religion that they do not acquire decisive significance. Of this kind
is religious thinking in those times in which religion is everything on
the spiritual plane. These are the classical times of religion: partly
the great times of origins, in which religion with all the power of
originality draws all strength and interest to itself; partly the great
times of organization, in which all existing culture is shaped or bent
into obedience under the highest religious ideas. In such classical
times there is a great unity, or at least a great harmony, in the domain
of spiritual life; the division of labour that sets in at a certain stage
of cultural development has not yet made itself felt in this domain.
% 3rd-ed. addition: the "Arbejdsdeling" clause is not in the 1901 text.
Christianity had such a classical time in the age of
the primitive Christian community, and later in the great organizational
period at the height of the Middle Ages. A religious problem forms only
when other sides of spiritual life---science and art, morality and
social life---begin to free themselves and to appear with the claim of
having independent value. They then judge religion from their points of
view and by their standards, whereas religion in its classical times
either simply ignored them---as in the times of origins---or---as in the
times of organization---valued them from its own point of view and by
its own standard, so that the question became whether these two
different evaluations could be brought into mutual harmony.---There must
be definite historical conditions present before the religious problem
can be posed in a decisive way. It is, if one will, only in unhappy
times that a religious problem exists. For this problem is the
expression of a spiritual disharmony: the different elements of
spiritual life no longer work together in so unified a way as before;
they lead in different directions, among which one must perhaps
ultimately choose. But then both the necessity and the possibility of an
investigation are also present.

That religion is made into a problem will be offensive to many. But
thought must have the right to investigate everything, once it has
awakened, and only thought itself can mark out its boundaries. Who else
should mark out its boundaries? He who perceives no problem has of
course no reason to think; but neither has he any reason to prevent
others from thinking. He who fears losing his spiritual support must
keep away. No one wishes to rob the poor man of his only lamb---but then
the poor man must not unnecessarily come dragging it along the highway
and demand that traffic be stopped on his account. Experience shows,
moreover, that it is rather the oxen than the lambs that loudly announce
themselves with their indignation in season and out of season
(especially out of season). It is not so much the genuinely poor in
spirit as the stiff and blustering dogmatists who take offense when free
inquiry claims the right to stir on the religious plane as on any other.

The investigation I wish to undertake here addresses itself neither to
those who are finished nor to those who are fearful. The finished are
found in every camp, not least among the so-called ``freethinkers''---a
class of people who, like the ``Worms'' in Linnaeus's system, can be
characterized only by negative predicates, because the class is supposed
to encompass so many different forms. The finished are in possession of
a definite negative or positive solution to the religious problem and
have thereby lost the sense for further thinking about it. The fearful
dare not think about it.---Rather, the investigation addresses itself to
those who are \emph{becoming}. ``One who is becoming will always be
grateful''---no matter in which direction his becoming goes.

Philosophy is, as is well known, richer in ideas, perspectives, and
discussions than in definitive results, and the philosopher would often
be in a sorry state if he did not believe that the constant striving
after truth belongs to the highest spiritual values, which it is the
task of all true religion to uphold. Even if we could learn nothing else
from a philosophy-of-religion investigation, we could perhaps learn
something about the causes of the strife that rages over religious
questions, and about this strife's significance for the development of
spiritual life. And should the religious problem prove insoluble, we
might perhaps learn why it is insoluble.

\medskip

\noindent\textbf{2.} The philosophy of religion should not proceed from
a finished philosophical system. In a certain sense all thinking must be
systematic, for ``system'' literally means ``that which stands
together,'' and it is an inexorable demand placed on our thoughts that
they must be able to ``stand together'' with one another. There is an
inner coherence and consistency which in every domain is the ultimate
mark of truth, and which must therefore also be required in the domain
of the philosophy of religion. But the philosophy-of-religion enterprise
will be most fruitful where religion and its phenomena are not placed in
relation to some completed philosophical conception or other, but are
illuminated by a philosophizing that is itself in the process of
investigating whether a completion is possible at all. The task is to
discuss religion's relation to spiritual life. Religion is itself a kind
of or form of spiritual life, and it is judged rightly only when seen in
its relation to other kinds or forms. It is no alien standard that is
thereby applied to religion. For there is a service from which no
spiritual power---even if it appears under the greatest names---can
exempt itself, namely to serve the deep and rich development of
spiritual life. The philosophy-of-religion task will therefore above all
be to find the means for investigating whether religion, under the
conditions of our spiritual culture, can continue to render such
service. In solving, or better: in treating this task, the philosophy of
religion cannot come equipped with other presuppositions than those
which every other science brings to the treatment of its subject. Very
often it will surely be compelled to declare that there are questions
that cannot be answered, because the conditions in the spiritual domain
are more complex and more diverse than positive and negative dogmatists
suppose. Yet it is not without significance that the problems are
discussed. The conscious insight into the impossibility of a solution
(\textit{docta ignorantia}) naturally invites one to investigate whether
it is then so absolutely necessary that a solution be found. In
addition, one's eye will be sharpened for the particular conditions
under which our spiritual culture lives, and above all for the
significance of the division of labour in the spiritual domain, through
which the religious problem becomes a special form of the general
problem of culture.
% REVISED IN THE 3RD ED. The 1901 text continued instead: "...conditions
% which would probably not continue to persist, just as they even now do
% not exert equal influence on all personalities." The 1924 text drops
% that clause and substitutes the division-of-labour point, matching the
% new clause added above in the "classical times" paragraph.

A philosophy-of-religion investigation can suitably be arranged so as to
consist of an epistemological, a psychological, and an ethical part.

In its classical times, religion satisfies all of humankind's spiritual
needs, including its drive for knowledge. Religious ideas then contain
for human beings the explanation of existence in its entirety and in its
individual parts. But when an independent science has developed, another
kind of understanding or explanation appears alongside the religious
one, and the question becomes whether these two kinds are compatible.
Many hold that the relation is this, that the new, scientific mode of
explanation increasingly replaces the religious explanation where
particulars are concerned, but that it cannot be applied to existence
considered as a whole---to what one customarily calls the ``first'' or
the ``last'' questions. Then one must discuss whether religion is able
to solve the riddles that science cannot solve. If it should then turn
out that religious conceptions, whether with respect to the whole or to
the particular, cannot help with anything that could rightly be called
understanding or explanation, then one must raise the question what
significance these conceptions can still have.

This last question leads from the epistemological philosophy of religion
to the psychological. For if religious conceptions no longer have
cognitive value, the value they might possess must consist in being
expressions of sides of the life of the soul other than the
intellectual. An attempt must then be made to describe the religious life
of the soul, especially the relation of conceptions to religious
experience and religious feeling. Through this it will become apparent
that religion in its innermost nature has to do not with the
understanding of existence but with the \emph{evaluation} of it, and
that religious conceptions express the relation in which factual
existence, as we know it, stands to that which for us gives life its
highest value. The core of religion is---according to the hypothesis to
which we are led---a faith in the persistence of value in existence.
This faith appears in great, vivid images in the popular religions,
especially in the highest among them. Even in those who stand outside
every popular religion, such a faith may stir, even though it does not
here take definite shape.

The transition from psychological to ethical philosophy of religion is
naturally motivated by the question: what ethical value does it have
that one believes in the persistence of value? Ethics concerns the
production of values, and the possibility cannot be dismissed that the
work of producing values leaves no time and no strength for living in
the thought of their persistence. The ethical philosophy of religion
forms a parallel to the epistemological. In its classical times,
religion directly and indirectly determines the evaluation of all
actions and all conditions of life. The religious problem arises on this
point when an independent ethics has developed that seeks, by its own
path, the grounding and motivation for the evaluation of human actions.
And on this point the problem is posed in the sharpest and most pressing
way, for the ultimate standard for judging religion must belong to the
ethical domain, where the final reckoning of the spiritual account takes
place.---

Within the framework here indicated, we shall be able to find room for
all essential aspects of the religious problem. I wish to treat this
problem not only with the intellectual interest that is aroused by a
great and comprehensive subject, but also in the mood that comes from
the consciousness of standing before a spiritual form of life in which
the human race through millennia has deposited some of its deepest and
most inward experiences.

\medskip

\noindent\textbf{3.} Before we enter upon the individual parts of the
philosophy of religion, it will be useful to give a somewhat more
detailed characterization of the content of the religious problem.

In the Protestant world it is fairly generally acknowledged that one must
distinguish between knowledge and faith. The Catholic Church does not
acknowledge this distinction in its full extent and therefore does not
acknowledge any religious problem in the intellectual domain. It
maintains that God's existence can be proved by the paths that Thomas
Aquinas, supported by Aristotle, pointed out, and it has even declared
it a matter of faith to accept that these proofs suffice. The tradition
from the great time of the Middle Ages, when ``one did not doubt and was
therefore nearer to the truth'' (to use a sentence from one of Thomas
Aquinas's modern biographers), is continually upheld, and it is not
acknowledged that a disharmony has really entered into spiritual life.
The distinction between knowledge and faith concedes that such a
disharmony is present---that the unity of spiritual life is broken. Is
this distinction itself then only a station on the way, in that
development will lead to an elimination of all religion from spiritual
life, because only that can endure which can enter into a harmonious
relation with the other elements of life?---Whether one now assumes that
it is only the classical time of religion that is past, or that its time
is altogether past, the question cannot be set aside whether spiritual
value is lost in the process of dissolution that has set in, or whether
a replacement can be found for the spiritual unity and self-containment
that was lost. This consideration alone shows that the question of the
persistence of value is an essential aspect of the religious problem,
whatever solution one inclines toward. It is not least from this side
that philosophy in the last century and a half (from the time of
Rousseau, Lessing, and Kant onward) has treated the religious problem.
Both Romantics and Positivists have been attentive to this aspect of
it.---It must of course be remembered that even if one could show that
no spiritual energy was lost in the dissolution of religion, it does not
follow that the energy would under the new conditions receive as
valuable an application as in the classical times of religion. In the
spiritual domain no less than in the material, the conservation of
energy is not without further ado the same as the persistence of value.
Energy can persist without value persisting, while the reverse cannot
well be imagined, since energy in and of itself will be a condition for
the persistence and development of what is valuable. Many freethinkers
proceed without further ado from the assumption that human life will
become more powerful and richer when religion falls away---an assumption
that is not self-evident and that rests on the presupposition that
psychic equivalents are present---equivalents both in terms of value and
in terms of energy. The task will then be to demonstrate these
equivalents and thereby establish the persistence of value. But it is a
great question and an essential aspect of the religious problem whether
such equivalents can be demonstrated.---Moreover, the problem of the
persistence of value will also arise within the Church, for religion and
the religious life undergo constant changes; and the question becomes
whether the values from the classical times of religion persist through
these changes---e.g., whether primitive Christianity, both in terms of
value and in terms of energy, is preserved in the Christianity of our
day. In its sharpest form this problem has appeared in the assertion:
``the Christianity of the New Testament no longer exists.''---

But the religious problem concerns not only the persistence of value
within the human world. It raises not only a psychological but also a
cosmological question, namely about the relation between that which
stands as the highest value for human beings and the whole of existence.
Is there any connection between values and the laws and forces of
existence? Are these laws and forces perhaps even ultimately determined
by the highest values? Or can we perhaps not ascribe any validity to the
concept of value beyond the domain of human life?

The hypothesis that will later (in the psychological philosophy of
religion) be sought to be grounded holds that the proposition of the
persistence of value is the actual religious axiom, which is expressed
in different ways in the different religions and at the different
religious standpoints. The religious problem thus concerns the question
whether this axiom can really be taken to hold. And at the same time
this axiom, insofar as it expresses the fundamental thought of all
religion, can be used as a standard for the consistency and significance
of the individual religion or the individual religious standpoint.
Finally---as already noted---the relation between ethics and religion
can, with the help of this axiom, be expressed in a very simple way: how
does faith in the persistence of value relate to the work of finding,
producing, and preserving values?---

If the proposition of the persistence of value were to be discussed in
its full scope, the discussion would become interminable, and clear
starting points for conducting it would soon be lacking. I therefore
emphasize at once the limits within which my investigation will move,
and which must be kept in mind throughout all that follows.---When I
attempt to show that the said proposition is the fundamental thought of
all religion, this attempt will not have failed because it turns out
that no religion expresses the proposition with clarity and consistency.
It will be sufficient if an explicit need or tendency to hold fast to
this fundamental thought can be demonstrated, so that the standard for
the evaluation of religions in their mutual relations becomes the degree
to which they are able to articulate and carry through the
proposition.---When the persistence of value is spoken of, this could be
understood to mean that value does not vanish, but that there is
constantly value in existence, whether it decreases or increases or
alternates between decrease and increase. But I understand the
expression \emph{the persistence of value} in analogy with the
expression \emph{the conservation of energy}, such that the proposition
asserts a constant persistence of value through all changes of form. It
will perhaps be necessary here to distinguish between potential and
actual value, just as one distinguishes between potential and actual
energy, and the one distinction is just as clear and just as justified
as the other.---But the persistence of value through all changes of form
and despite the difference between possibility and actuality is
nonetheless only a minimum. It may be maintained that the mere
persistence of value is not satisfying and actually contains a
contradiction, since value can persist only by growing, since mere
repetition dulls, while on the other hand change is valuable only when
it brings enhancement. That an increase of value does not conflict with
the conservation of energy needs no further demonstration; the increase
of value presupposes only a new and more perfect application of energy,
not an increase of energy as a whole. It is for the sake of simplicity
that I limit the investigation to the persistence of value. And it will
also turn out that religious consciousness has at bottom limited itself
to this as well. The persistence of value will present us with
difficulties and tasks enough, even if we do not enter into everything
that lies in the unavoidable consequence that value can persist only if
it increases. Yet we must not forget that this consequence is
present.---Finally, when the persistence of value is under discussion,
one might ask: which value and how much value? I enter into these
questions only insofar as the difference between the different religions
and religious standpoints essentially depends on the different answers
they give to them. It is clear that the proposition of the persistence
of value can be accepted whatever answers one gives. The hypothesis that
the persistence of value is the fundamental religious proposition is not
shaken by the fact that the different religions and religious
standpoints diverge greatly from one another with respect to the
character and extent of the valuable in whose persistence they believe.

Value is the property of something that it brings about a real
satisfaction or can become the means to such. Value can therefore be
immediate or mediate. Immediate value we seek to hold fast when it is
given, and to produce when it is not given. We thus make it our end.
That we make something an end presupposes, conversely, that we have an
experience or a supposition of its immediate value. Mediate value
belongs to that which serves the attainment of an end---that is, of an
immediate value. Mediate value need not be potential value. For what
serves as a means for attaining something of immediate value itself
acquires value in our eyes, even though its value stands as derived.
There are all possible transitions between mediate and immediate value,
and through the shifting of motives the former passes over into the
latter, so that what at first had value only as a means now acquires
value as an end. Potential value, on the other hand, often means only
the possibility that something will acquire value as a means---that is,
the possibility of mediate value.---Which values are experienced and
therefore acknowledged at the different standpoints will depend on the
different character of those standpoints. It is a being's nature that
determines its need, and it is its need that determines what acquires
value for it. The religious axiom thus shows the necessity of a
religion's character being determined by the nature and need of the
human beings who profess it. For one cannot seriously believe in the
persistence of values that one does not know at least approximately from
one's own experience.---

% REVISED IN THE 3RD ED. — the largest change in chapter I.
% The 1901/1906 text closed the chapter with a passage on Kant and the
% fundamental problem: "It is to Kant's philosophy that we owe the
% independence of the problem of value from the problem of knowledge. He
% taught us to distinguish between evaluation and explanation. [...] The
% fundamental philosophical problem concerns the continuity of existence in
% its relation to the special and individual forms of existence. And if the
% hypothesis that the proposition of the persistence of value is the
% religious axiom turns out to be well-founded, then the religious problem
% concerns the continuity of existence from a distinctive point of view and
% thus falls under the fundamental problem as a special form of it."
% The 1924 text deletes all of that and substitutes the passage below on
% value and wholeness, with cross-references to Den menneskelige Tanke
% (1910) and Totalitet som Kategori (1917) — books that did not exist when
% the first edition was written. The "special form" formula survives, but
% relocated and recast: it now appears earlier in the chapter, where the
% religious problem is called a special form of the general problem of
% culture. The full 1901 wording is kept in transcription.tex.

Every application of the concept of value presupposes a relation to a
whole whose persistence and development require definite conditions.
Such wholes are organisms, personalities, societies, and finally the
human race. In an individual whole endowed with consciousness (a
personality), feelings of pleasure and displeasure, in all their various
degrees and forms, appear as symptoms of a need to maintain or to call
forth the conditions of the whole's persistence. Here, then, value is
known by immediate experience. But by way of analogy the concept of
value---and with it the concept of wholeness, which makes rational
discussion of values possible---can be applied to all wholes.---A fuller
demonstration than can be given here of the significant connection
between the concepts of value and wholeness I have attempted in
\emph{Den menneskelige Tanke} (pp. 238--241; 346--349) and in
\emph{Totalitet som Kategori}, ch.~5.

\begin{center}---\end{center}


%% ============================================================
%%  II. EPISTEMOLOGICAL PHILOSOPHY OF RELIGION
%% ============================================================
\chapter{Epistemological Philosophy of Religion}
\label{ch:epistemology}

\section{Understanding}
\label{sec:understanding}

\begin{flushright}
\textit{Ihm ziemt's die Welt im Innern zu bewegen.}\\[2pt]
\textsc{Goethe}.\\[2pt]
\textit{(It befits him to move the world from within.)}
\end{flushright}

\bigskip

\noindent\textbf{4.} Against the supposition I have put forward, that
the core of religion is faith in the persistence of value, it might seem
to tell that theoretical motives have after all played, and still for
many play, a considerable role in the development of religious ideas. I
have moreover already remarked that since religion in its classical
times is everything for the human being, it in such times also satisfies
its drive for knowledge---that is, provides it with means for and forms
of understanding and explanation of existence. In the classical times of
religion there exists no independent drive for knowledge, any more than
special means and forms for satisfying such a drive. The distinction
between knowledge and faith, between understanding and evaluation, comes
forward only in the critical times---that is, in the times when religion
has become a problem. This distinction began to make its way only after
a number of violent collisions between religion and science. Only
reluctantly has one begun, from the religious side, to learn that it is
not religion's business to give us a scientific understanding of the
world. What all theologians now repeat, that the Bible is not to teach
us natural science, no one would hear when Bruno, Galileo, and Spinoza
said it. The persecution of heretics ran its course, and only afterwards
did one see the rightness of what the heretics had said. In other
domains the same conflict repeats itself; in our days it is the
humanities whose independence is at stake.

Every great religion appears in history with a world-view, in that it
gathers up within itself, in a distinctive manner, the whole circle of
ideas of its time. Christianity, for example, stood in the first
centuries as attractive to many by its intellectual content as well.
\textsc{Tatian} thus went over to Christianity because this ``philosophy
of the barbarians,'' as he called it, seemed to him to explain what he
had not before understood: ``the coming-to-be of the world.'' The old
apologists stood generally at the standpoint that Christianity contained
the best solution to the riddles of thought. When Plato and other Greek
philosophers seemed to hint at similar solutions, this was to be
explained, according to Justin Martyr, either by their having been
influenced by the Israelite prophets, or else by the fact that ``the
Word'' (Logos), before it became flesh, had worked illuminatingly within
human beings. In comparison with the earlier religions, Christianity
had---as is generally the case with the higher religions in relation to
the lower---a certain rational character, and above all a simplicity, a
plain grandeur, which could not but work attractively on many who knew
only the earlier religions with their motley mythologies. In so far as
the world-view given in the biblical writings was not detailed enough,
the Church Fathers and the Schoolmen supplemented it by the aid of Greek
science.

Just as Christianity does for the western portion of humanity, so
Buddhism stands for the eastern portion not merely as a great religion,
but as a great world-view, whose ideas have a certain rational character
and cannot be understood without acquaintance with the antecedent Indian
thought.

When the opposition between religious and scientific understanding of
existence is spoken of, the weight is as a rule laid on the fact that
science---especially in the domains of astronomy, geology, and
biology---has led to results that conflict with religion's transmitted
teachings. A great part of the work of the modern apologists accordingly
goes to showing the compatibility of the transmitted religious ideas
with those scientific results that really stand firm. Towards modern
natural science one applies, with greater or lesser success---and with
greater or lesser taste---the same process of accommodation that the
apologists of antiquity applied towards Greek science. But for the
illumination of the religious problem as a whole, this entire discussion
for and against is of slight interest, except perhaps in so far as it
furnishes characteristic examples of the accommodation which the life of
religious ideas---like everything living---undergoes, or strives to
undergo, in the face of new conditions of life. What is the main matter,
on the contrary, is that it is two wholly different kinds or types of
``understanding'' that here stand over against one another. They are two
ways of thinking and of explaining events---two ways so different that
once one has begun with the one, one will not be able to find any point
at which it is justified and possible to go over to the other. The whole
mental condition is so heterogeneous that the word ``understanding'' is
in reality used here in two wholly different senses. Let me take a
single example. A sudden storm and tidal flood, which brings about the
destruction of many fishermen at sea, is explained in natural-scientific
fashion by the states of the air and the sea in the preceding days. But
in the funeral oration over the drowned, the ``understanding'' of what
has happened is found in this, that God wished to give the survivors a
sign, that they should be converted.\textsuperscript{1} At what point in
the series of natural causes can one now suppose that this divine
intervention has taken place? The natural causes form a connected
series, from which no link can be taken out. A choice must be made
between the two ``understandings.'' The one kind of understanding aims,
so far as possible, to form connected series of links, each of which can
be pointed out in experience. The more the state of the air and the sea
at the moment of the catastrophe has been brought into connected
sequence with their states in the preceding days, the better---according
to the one type of understanding---one is said to understand the
catastrophe. According to the other type, it is understood by the value
it acquires through its influence on the minds of those who suffer under
it, and this value is regarded as intended---as the purpose of a divine
intervention.

The scientific type of understanding has developed in a clear and
fruitful manner especially in the course of the last three centuries. It
does not exclude a religious evaluation, but it excludes a confusion of
this evaluation with explanation proper. It is the task of philosophical
epistemology to investigate the presuppositions and the forms of thought
that lie at the foundation of scientific knowledge, and it will be of
philosophy-of-religion interest to bring out these presuppositions and
forms of thought and to compare them with the foundation of religious
``understanding.'' Through this we shall be able to obtain information
as to how far religion will still be able to have significance for the
satisfaction of our drive for knowledge.

\subsection{Causal Explanation}
\label{ssec:causal}

\noindent\textbf{5.} When one sets up the proposition that every event
has its cause---the so-called causal principle---one thereby establishes
the principle that existence is intelligible. This proposition can never
become more than a hypothesis, since we can never understand everything,
never point out causes for all events. But it contains the
presupposition with which we involuntarily meet events that awaken our
attention; at any rate it is the presupposition that brings us to look
about for earlier events when an event has occurred.

We can here disregard entirely the question of the causal principle's
justification and origin, since the mutually divergent conceptions here
offered need not acquire any philosophy-of-religion significance.
Whether the acknowledgement of this proposition is due to habit or to an
inner necessity in our nature, the need to find causes stirs more or
less in everyone, and the religious conception bears witness to it as
well as the scientific. Indeed, religious ideas have a peculiar
epistemological interest precisely in this, that they testify how deeply
the need to find causes lies in human nature, since it expresses itself
even there where no independent scientific interest has yet arisen. The
opposition between religion and science comes forward only when regard
is had to the \emph{character} of the causes one acknowledges as right,
real causes (\textit{veræ causæ}). \emph{Both religious and scientific
understanding are a leading back of the unknown to the known; the
difference between them rests on what this known is.}

Scientific understanding requires that the cause of an event be
demonstrated in events that themselves establish themselves in
experience just as well as the event to be explained. The task here
becomes to interpret the series of events by finding as many individual
links in this series, and as close a connection between them, as
possible. The aim is to interpret nature by nature itself, just as one
interprets a passage in a book by comparing it with other passages in
the same book. Religious understanding does not make the demand that the
cause should stand in continuous connection with the effect and itself
be demonstrated as an experience just as well as the event to be
explained. It does not investigate its causes in the critical manner
that science does. It does not explain nature by nature itself, but by
something that is different from or stands outside nature, that
``pushes in from without.'' Thus, at any rate, religious understanding
most frequently appears in history.
% 3RD-ED. CUT. The 1901 text continued here: "; whether it is necessary
% that it appear in this form, and therefore in marked opposition to
% scientific understanding, will be shown later." The 1924 text drops the
% forward reference and begins a new sentence.
In contrast to this religious understanding, the special form that
empirical science---in both the mental and the material domains---gives
to the causal principle can be designated as \emph{the principle of
natural causes.} If it cannot be applied, we understand nothing, when it
is the scientific type of understanding and the corresponding
intellectual habit that prevail in us.

Only the principle of natural causes satisfies the need for a special
explanation of the particular, determinate event. Galileo thus showed
that in the debate between the different contending astronomical
hypotheses it could be of no use to appeal to God's intervention, since
for God's omnipotence it is just as easy to let the sun move round the
earth as to let the earth move round the sun. This cause therefore
cannot help us to a choice between the hypotheses. It is a key that
unlocks everywhere; but it is not such a master key that science seeks;
science's interest is to investigate the peculiar character of the
different locks, and it therefore demands a key of its own for every
lock. If a master key were sufficient, the difference of the locks would
become meaningless---and would contain an unsolved problem!---To appeal
to God's will gives no scientific understanding, but is---scientifically
regarded---an admission that one does not understand the event.
% 3RD-ED. CUT. A passage of the 1901 text beginning "Fremdeles. I
% Tvivlstilfælde…" stood here and is absent from the 1924 text.

The deepest justification for the principle of natural causes I find in
a need for continuity, which lies in the nature of our consciousness and
which already shows itself in the general need to find causes. Our
consciousness strives involuntarily to gather up and hold together its
elements and to preserve unity with itself, even when it stands before a
new content. This new content it then strives to shape, alter, and
arrange, so that the new shows itself to be a new form of something old.
Self-recognition then becomes possible in spite of the shifting of
content, and the unity of the personal life mirrors itself in the
continuity thereby brought about. The work of knowledge is here closely
bound up with the work on the development of character. We speak of a
character---in the pronounced sense---only where unity and continuity
prevail in the whole striving, and where there is not again and again a
leaping over to something new. In analogy with this, knowledge seeks to
find as much unity and continuity as possible, to exhibit the new as a
transformation or continuation of the old; and only where this succeeds
is a full appropriation and understanding won. The principle of natural
causes, which leads one to build series as long and as continuous as
possible, is therefore no accident, but is closely bound up with the
very nature of personality. It is therefore incorrect when personality
and science are set in sharp opposition to one another. Scientific work
is a work of personality. And the more this work succeeds, the more does
causal explanation approach grounding and recognition. The differences
in existence are reduced to the least possible when we survey, at a
single glance, ``under the form of eternity,'' great regions of
existence.
% 3rd ed. reads "under Evighedens Form"; the 1901 text had "under
% Evighedens Synspunkt" (under the point of view of eternity).

Religious understanding too---in the mythical and dogmatic forms in
which it most frequently appears in history---leads the unknown back to
the known. But this known is here as a rule the will of one or more
personal beings. The religious consciousness supposes itself to know
what the divine will aims at, and recognizes in events the expression of
this will. Its interest is not to find an inner connection among events
themselves, but to see all events (in the whole of nature, or in some
single domain) as the expression of one and the same power. It places
itself at the centre, or supposes itself to be at the centre, and from
there it surveys the whole circle of existence.

The difficulty for the religious conception here becomes this, that the
radius with which this circle is described is not known. From one and
the same centre many different circles can be described. Therefore every
``central'' consideration (as it has moreover been attempted not only on
religious presuppositions, but also in speculative philosophical form)
will be powerless with respect to the understanding of the special
events. A theological or speculative epistemology will therefore not be
able to be carried through. There do not even exist serious attempts at
one; but we must hold to the epistemology that empirical and critical
philosophy has developed on the basis of the history of empirical
science.

\medskip

\noindent\textbf{6.} When the scientific type of understanding has begun
to prevail, but is excluded at certain selected points where religious
understanding still claims dominion, there arises the bastard concept
called wonder, \emph{miracle.} The miracle presupposes that both types
make themselves felt, but that there is a leap over from the one type to
the other. Where either a certain degree of scientific understanding or
of knowledge is not present, the miracle can only mean a deviation from
the accustomed. The natural man wonders when the hitherto unseen or
unheard-of happens. But the miracle proper presupposes a certain
knowledge of natural laws---a continuity that is broken---a rule from
which an exception is made. Miracles must therefore occur according to
the law of parsimony; if they come so frequently that they become the
rule, they form a continuous connection, and the concept of miracle is
abolished.

If one does not set up the principle of natural causes in a wholly
dogmatic manner, one cannot maintain in advance the impossibility of
miracles. But the decisive thing is here again the way the problem is
posed and the intellectual habit. For him who asks and sets up his
problems according to the scientific type, for him a miracle can never
come to be. He may come to stand before something he cannot explain,
since no cause (in the sense in which he speaks of cause) can be found,
or at any rate has hitherto been found. But he can never discover that
there is no natural cause. The religious consciousness itself
cannot---without ascribing omniscience to itself---maintain that an
event \emph{cannot} be explained according to natural laws. To know that
something falls outside the laws of nature, one must know them all and
know all possible special applications of them. Only by a revelation
could one come to know that in the particular case a miracle is present.
It is not the miracle that confirms the revelation, but the revelation
that decides what is a miracle. And since a revelation is itself a
miracle, one can only by a miracle come to know whether a miracle has
happened.---It is consistent when the Catholic Church, which decides
what is revelation, also demands to decide what is a miracle and what is
not. Nothing may be taken for a miracle unless it is sanctioned as such
by ``apostolic and ordained authority.''\textsuperscript{2}

In the case of miracles of which one is not oneself an eyewitness, the
question will particularly be whether the report is so securely
transmitted, and whether the eyewitness has been so keen an observer,
that one will sooner be able to assume a deviation from the laws of
nature than an error in transmission and observation. A report that does
not derive from an observer who has sufficient acquaintance with the
natural conditions and natural laws in question has no scientific
significance.

\medskip

\noindent\textbf{7.} The preceding consideration rested on the
scientific concept of cause, according to which an event is understood
by being seen in its inner connection with other events, as a
continuation of them or as a transformation into new form of their
content. The presupposition of all such causal explanation is a
principle of unity that makes connection possible. We reach this
principle by analysis of the very concept of causal connection or
reciprocal action. It is the bond that holds existence together and
makes it a whole, more than a chaos of elements---even if it is a whole
we cannot survey.

The principle of unity at which we have thus arrived recalls the concept
of God. But the customary line of thought often reaches its concept of
God by another road: not by going back to the presupposition of the
causal series, to the principle that makes such a series possible at
all, but by continuing the causal series backwards to a ``first cause.''
By virtue of this line of thought one will say, in the face of the
preceding development: ``Religion gladly leaves it to science---a few
supernatural events excepted---to find the finite causes and to spin out
the causal series as far as possible. The causes that science can find
are nevertheless always subordinate, derived causes (\textit{causæ
secundæ}). What religion maintains is that existence becomes meaningless
if we do not at last find the origin of the whole causal series in a
first cause, an absolute cause, in which thought finds rest---where all
further questioning falls away---and from which the full, clear light
streams out over everything.''

This line of thought stems from \textsc{Aristotle} and is taken up by
\textsc{Thomas Aquinas}, on whose authority it is still espoused by
Catholic theology,\textsuperscript{3} as well as by many Protestant
theologians. It is the so-called cosmological proof of God's
existence---the most important of all ``proofs''---that we here stand
before.

That the series of causes must be limited is inferred from the
consideration that otherwise no complete explanation of any link
whatever in the series of events would have been reached. ``If there is
no first,'' says \textsc{Aristotle}, ``there is no cause at all.'' It is
thus the condition of existence's intelligibility that one here supposes
oneself to be relying on. It is a rationalistic tendency that lies at
the foundation of this theological line of thought. Existence must be
comprehensible!---that is the underlying postulate. Aristotle does not,
however, assume that the first cause began to work at a definite point
of time, so that ``the world'' has a beginning. And Thomas also concedes
that it is only by the way of faith that one comes to assume that the
world has a beginning (\textit{mundum incepisse, sola fide tenetur}).

\textsc{Kant} has subjected this proof to a criticism that leads him to
the result that it rests on ``a false self-satisfaction of human
reason.'' In the main his criticism stands in force to this very day. In
what follows I make use in part of his objections, adding to them some
other misgivings.

Even if we could build a connected causal series, at what point should
we then break off our search for finite causes in order to hang the
whole chain on a ``first cause''---as Zeus in Homer threatens to hang the
whole world on a peak of Mount Olympus? What mark can we have that we
are at the ``next to first'' cause, so that we must now set about making
the decisive leap?---And why should the causal series not be able to be
infinite? That an infinite series is unthinkable is a dogmatic
assertion. There can be series which, by virtue of the law that holds
for the relation between their links, can always be continued. Such
series we think---not by running through all the links, which is of
course impossible---but by thinking the law of their connection. Such a
series is the causal series. It is built by virtue of the causal
principle, which says that everything that happens has its cause. This
proposition must---if it is accepted at all---also be accepted as
holding for every single link in the series, for the ``first'' one has
gone back to as well as for all the others. The assumption of an
absolutely first cause would thus conflict with the causal
principle---although it was precisely to be justified by existence's
absolute intelligibility, and so by the causal principle!---

Popular theology will especially hold to the form of the cosmological
proof that runs thus: ``Since everything has its cause, the world too
must have its cause.'' One operates here with the concept ``world''
without noticing that it contains the same or similar difficulties as
those that just showed themselves in the concept of God. If by ``world''
one understands the sum total of all that exists, or at any rate of all
``finite'' existing things, then it is a concept that can never be
completed. The given is never completed; new experiences are always
coming in, which make new conceptual determinations necessary. We
therefore never come, by the road of scientific inquiry, to a point at
which there could be a question of a cause of ``the whole world.'' The
given constitutes a whole, which we analyse in order to find its laws
and its elements; but every given whole points over to other wholes; we
seek to work all given wholes together into a still higher whole; but
however far this work be continued, the same task will always present
itself anew. There is here a series of like kind with the causal
series---a series that by virtue of its own law cannot be completed. The
concept ``world'' is therefore properly a false concept, though one
operates with it so briskly, plays with it as with a ball.

\subsection{The World of Space}
\label{ssec:space}

\noindent\textbf{8.} When the religious consciousness historically shows a
tendency to limit the series of causes, this is in part to be explained by
the need for intuitive vividness. The religious consciousness wants a vivid
image of its object, and such an image presupposes limitation. With this in
turn hangs together its tendency to refer the divine to definite places in
space. Spatial ideas are our clearest ideas, and what we wish to represent
clearly we therefore seek to shape into spatial images or schemata. With
this there naturally combines the idea that everything that exists must
belong at a definite place in the world, must be localized. What is not at
some definite place seems, to the childlike conception, which is decidedly
realistic, not to be at all. Just as, with respect to the explanation of the
world, it would be the simplest thing if we could come upon the deity by
following the series of causes link by link, so also, with respect to the
conception of the world, it would be the simplest thing if the human being
needed only to lift his eyes upward to find the deity, or at any rate could
behold its dwelling.

Antiquity's conception of the world met this need. For it the world was
limited. The earth stood at the world's centre, and the vault of heaven,
which was not thought to be extraordinarily far off, formed the world's
outermost boundary. The bright sky has, especially for the Indo-Germanic
peoples, been regarded as deity or as the seat of the deity. The stem from
which many of the Indo-Germanic languages form their word for ``God''
properly denotes sky or heavenly. Thus the Indian \textit{devas}, the Persian
\textit{deva}, the Greek Ζεύς (Δίος), the Latin \textit{deus}, and so on.
% The e-book sets the Greek with Latin letterforms mixed in ("Zεύς (Δίoς)");
% corrected here, as with the Aristotle epigraph in ch. I. A digitisation
% defect, not a print reading.
Herodotus says of the Persians that they call the sky Zeus, and with this
agrees the fact that the natural background of the god Ahura Mazda (Ormuzd)
is assumed by the more recent history of religion to be the sky. The Greeks
supposed, indeed, in the oldest time, that the gods dwelt more particularly
on Mount Olympus, but later thought of the sky or ``the aether'' as their
dwelling. Aristotle says: ``All human beings assume that there are gods, and
all, both barbarians and Hellenes, assign to the deity the uppermost place.''
The opposition between heaven and earth thus became quite naturally, for
antiquity, one with the opposition between the divine and the human, between
the eternal and the changeable, between the perfect and the imperfect. The
expressions ``the higher'' and ``the lower'' had originally a literal
meaning; the symbolic meaning is here, as everywhere, the later one.
Aristotle's incorporation of antiquity's astronomical conception into his
system became of decisive importance. Since both this conception and his
whole system seemed to be compatible with the biblical circle of ideas, and
capable of serving for the further working-out of it, they were taken up by
ecclesiastical theology and through it became dominant right into the modern
age.

Among the Jewish people too we meet with the great religious significance of
the opposition between heaven and earth. In the opinion of some, Jahveh was
originally a thunder-god; at any rate he reveals himself with special majesty
in the heavenly phenomena. Heaven is his seat, and the earth his footstool.
The New Testament conception stands here entirely on the ground of the Old.
There is sometimes talk of several heavens (2 Corinthians XII, 2; Ephesians
IV, 10), and God's dwelling is then the heavenly holy of holies, which is
above all the heavens. Accounts such as that of Jesus's visible ascension and
of Paul's rapture to the third heaven testify that the ancient picture of the
world is presupposed. As these accounts stand, they are plainly to be taken
literally, though the more recent theology (even the orthodox) is on this
point inclined to apply the rationalistic explanation.

There was both vividness and security in the narrow frame within which the
picture of the world thus presented itself. Without being disturbed and
disquieted by infinite distances, the religious consciousness could give
itself wholly to the great experiences of life and regard the all-decisive
process of redemption---resolved on in heaven and accomplished on
earth---as that about which the whole outer life of the world also turned.
There was as it were a ladder between heaven and earth; the two worlds, the
world of the highest value and the world of contending values, stood in
vivid reciprocal action. Poets like Homer and Dante are testimony to what
foothold and what support the limited picture of the world afforded the human
imagination.

But within the religious consciousness itself there shows itself, at a
somewhat higher stage of development, an opposite tendency, which at the cost
of vividness aims at abolishing the localization of the content of religious
ideas. It came to appear an imperfection that the deity should really have
its dwelling at a definite place, so that it must either move through space
in order to come to the human being, or the human being must travel long
roads to meet his God. God must be omnipresent and all-embracing if he is not
to be limited and imperfect, as a human being is. However short the distance
between heaven and earth was thought to be, it nevertheless became too long
for the religious need. The deity had to stand in a far more intimate relation
to the human being than localization at definite places made possible.

In the biblical writings and among the ecclesiastical theologians we find,
then, a double current. Vividness is held fast; but where it is the
inwardness of feeling, rather than the imagination's need for intuition, that
leads the word, there the significance of the spatial relation is abolished.
God dwells in heaven---but he is truly not far from us, for in him we are and
live and move. Jesus ascends to heaven and shall come thence again---but he
lives in believers and is with them all their days.

The religious conception has here, in its classical times, a
\emph{Both---And}, which later is gradually replaced by an
\emph{Either---Or}.

\medskip

\noindent\textbf{9.} What made it possible for the theologians of earlier
times to hold fast to this Both---And was the bold use they made of
allegorical or symbolic interpretation. An allegorical interpretation did
not, according to their conception, exclude the literal meaning. The relation
between the two remained in a certain indeterminacy, which not seldom could
lead to violent collisions. ``There was lacking,'' says \textsc{Adolf
Harnack} in his history of dogma, ``a definite rule for how far the letter of
Holy Scripture should hold. Has God human shape, has he eyes, has he a voice,
does Paradise lie on the earth, do the dead rise with all their limbs, with
hair too, and so on?---to all these questions and a hundred like them no
certain answer was given, and therefore there was strife between the adherents
of one and the same confession. Allegorize all had to, but all had on the
other hand also to take certain passages of Scripture literally. But what a
difference there was between an Epiphanius and a Gregory of Nyssa, and how
many nuances there were between the crude anthropomorphists and the
spiritualists! These last for the most part feared the arguments of the
former, often also their fists, and they had to be anxious about their own
orthodoxy.''\textsuperscript{4}

\textsc{Augustine} is certainly the one who has contributed most to
displacing the merely literal conception. His Platonic studies have here had
decisive influence and have---in conjunction with his profound and strongly
agitated inner life---led him to a clearer view of the difference between the
spiritual and the material, and thereby led him to utter thoughts in which he
stands as the precursor of Descartes and Kant. What in this connection
especially interests us is his rejection of all spatial determinations as
invalid for the divine. Where he speaks of his own earlier conception, he
says: ``I did not know that God is a spirit that has no limbs with length and
breadth, any more than bodily mass (\textit{moles}). The bodily
(\textit{moles}) is less in a part than in the whole, and if it is infinite,
it is less in a limited part than in the infinite, and it is not everywhere
as a whole (\textit{tota ubiqve}), as God is.'' ``God is present in his
wholeness everywhere, and is yet nowhere in place (\textit{ubiqve totus, et
nusqvam locorum})!'' ``God dwells further within my being than my own inmost
self, and he is more exalted than the highest I can reach (\textit{interior
intimo meo, superior summo meo}).'' ``God stands above my soul, but not in
the same way as heaven stands above the earth.''\textsuperscript{5}

When thoughts such as these were appropriated and carried through, they had
to lead to strife with the old picture of the world and its childlike
opposition between heaven and earth. The Schoolmen of the Middle Ages held
fast to the immaterial concept of God, independent of spatial determinations,
and from them it passed over to the Mystics, who besides were influenced by
Neoplatonic thoughts.---``Lord,'' Suso's female disciple asks him, ``where is
God?'' The answer runs: ``The masters say: God has no Where\ldots. God is a
circle whose centre is everywhere and whose circumference is
nowhere.''\textsuperscript{6} Thus, through inner contradiction, the validity
of the spatial relation is abolished.

Towards accounts such as that of the visible ascension one took up in the
Middle Ages this position: that one held fast that the events reported had
taken place as they were described, but that they had a spiritual meaning
which was the essential thing. \textsc{Hugh of St.\ Victor} explains that by
``the highest'' is meant the inmost in us, and that ascending to God means
that we sink ourselves into ourselves and there find something that is
``higher'' than ourselves. But he will not therefore deny that heaven and hell
are definite places in the world. His contemporary \textsc{Abelard} maintains
(in his remarkable ``Dialogue between a Philosopher, a Jew, and a Christian'')
that what is said of God in bodily form is to be understood---not, as the
common man understands it, bodily and literally---but mystically and
allegorically. By heaven's exaltation above the earth is meant rather the
exalted character (\textit{sublimitas}) of the life to come than the situation
of the material heaven, and that Jesus sits at God's right hand denotes not a
spatial position (\textit{localis positio}), but equality of dignity
(\textit{æqvalis dignitas}). Jesus's ascension has indeed taken place bodily
as it is reported, but it denotes nevertheless ``a nobler kind of ascent''
(\textit{melior ascensus}), namely that which takes place in the souls of
believers. Unlearned and simple people would not understand what is not
presented in vivid form.\textsuperscript{7}

The characteristic Both---And here comes clearly forward. It had its great
advantages in a pedagogical respect, since one did not need to make any sharp
distinction between idea and reality, symbol and truth. With his bodily eyes
the human being could look up to heaven, the seat of power and grace, but the
same heaven he had in the inner states in which the divine made known its
authority and bestowed its consolation. The evangelical narratives could be
taken according to the letter and yet at the same time be interpreted
spiritually and edifyingly. The absolute differences of place in the world of
space appeared immediately as the expression of differences of value.

\medskip

\noindent\textbf{10.} There came, however, the time when one had to choose.
There arose a conflict between an idealistic conception, which inculcates the
purely spiritual significance of religious ideas, and a realistic or
materialistic conception, which held to vividness and literalness. Such a
conflict is known in many religions. In the \emph{Upanishads,} the idealistic
interpretation of the Vedic religion, it is said that Bráhman, the deity, is
imperishable, while ``name, place, time, and body'' perish, wherefore none of
these determinations holds for Bráhman. In Xenophanes's and Plato's criticism
of the Greek popular faith a similar tendency shows itself. Within
Mohammedanism it comes forward when, for example, the sensuous and vivid
depiction of the joys of Paradise is interpreted allegorically as a
presentation of spiritual enjoyments.\textsuperscript{8}

When it began to dawn on thought that there are no absolute determinations of
place, since every place is determined by its relation to other places, these
again by their relation to others, and so forth, the old picture of the world
was moving towards its dissolution. It went with the series of places as it
went with the series of causes; one could not get the series absolutely
completed, so that the picture could be bounded. And when now Copernicus
showed that the simplest picture of the world was obtained by thinking of the
earth as moving round the sun, so that it was no longer the world's centre,
there fell away the clear and secure frame within which the content of
religious ideas had hitherto been able to be localized. It was besides soon
shown---by Giordano Bruno---that it was inconsistent to hold fast to the
limitation of ``the world'' when one accepted the relativity of all
determinations of place and the Copernican astronomy. An infinite horizon
opened.\textsuperscript{9} The difference between heaven and earth fell away,
and an Either---Or had to replace the earlier Both---And: only the idealistic
meaning of the word ``heaven'' was now tenable, if this word was to express a
religious idea.

To this there now came also another movement of thought. Just as one had
hitherto found spiritual differences of value expressed in a vivid way by
spatial distances, so one had generally made no sharp difference between the
spiritual and the material. With the beginning of natural science this
difference came forward more clearly, and it came---with Kepler and
Descartes---to stand as a chief mark of the material that it is extended in
space. It had now to become a special problem, not merely how the spiritual
and the material stand related to one another, but also whether the spiritual
world is as comprehensive as the material world, according to the new picture
of the world, showed itself to be. The security with which one had hitherto
been able to find spiritual things expressed in material forms was quite
particularly shaken by this. Formerly one had been able literally to see
``spirits''; after the Cartesian reform of psychology this was no longer
possible. The spiritual world came to stand in opposition to the world of
space, of extension, and of vividness, however one might think of the inner
relation between them. Everything vivid could now at last, that is to say,
from a religious point of view, acquire only symbolic significance.

But it was easier for philosophy than for theology to draw this consequence.
With thinkers like Bruno and Spinoza the religious ideas underwent an
inwardization and an infinitization that corresponded to the new conception
of the world and carried further what Augustine and the Mystics had already
hinted at. After long resistance and misgiving, theology began to strike into
the same road. And often we now see the assertion put forward that the
idealistic meaning of such expressions as ``heaven'' (or ``hell,'' and so on)
is the only valid one, and the only one that the Bible and the Church have
ever approved. \textsc{Cardinal Newman} declares, for example, that the
Council of Trent does not teach that Purgatory causes sensuous pain, though he
concedes that this was the Latin Church's tradition, and that he himself has
seen in the streets of Naples paintings that represented souls burning in
Purgatory. But in an ecclesiastically sanctioned handbook on the doctrine of
Purgatory the question is raised: ``Is the flame of Purgatory really a bodily
flame?'' And to this it is said: ``Cardinal Bellarmine answers that the general
opinion of theologians is that the flame of Purgatory is a true and real
flame, of the same kind as our earthly fire.'' And it is added that the great
majority of Church Fathers, Schoolmen, and theologians teach that Purgatory,
like hell, is found at the earth's centre, and that this doctrine is built on
Holy Scripture and has always been maintained by the
Church.\textsuperscript{10}---Here idealism and realism stand sharply over
against one another. One may by particular distinctions seek to veil the
breach in spiritual development that let an Either---Or replace a
Both---And (just as this had in its time replaced the immediate, childlike
localization); but history will always bring it forth anew, and it is of the
very greatest importance for the right understanding of the religious problem
that it not be forgotten.

By the Copernican and the Cartesian reform there have been drawn more definite
boundaries than could be drawn before, between thought and intuition, reality
and symbolism. The philosophy-of-religion significance of this is easy to see.
Of quite special significance is the expansion of the world, the infinite
horizons that the material world has gained, and in relation to which the
spiritual world now seems to be so limited. How will one---the question arose
naturally---now be able to hold fast to the conviction that spiritual values
are the highest thing in existence, that about which everything turns? Now the
outer world no longer turns about the home of the only spiritual life we know
from experience!---There arose an opposition between psychology and cosmology
which the past did not know. Great spirits like \textsc{Jacob Böhme} and
\textsc{Blaise Pascal} felt keenly the problem that lay in this, and which in
religion's classical times could not be posed.---It is not, however, only for
positive religion that a great, perhaps insoluble problem is here posed. For
every view of life the same problem must present itself in one form or
another. Naturally there are problems at all only for those who have not from
indolence or anxiety given up thinking. To the problem itself we shall later
return.

\subsection{The Course of Time}
\label{ssec:time}

\noindent\textbf{11.} The need for intuitive vividness makes itself felt
with respect to time as well, and here too it leads to a shortening and a
closing-off of the ideas. In this respect, however, we meet with one of the
most significant differences that the history of religion anywhere exhibits,
in that the idea of a completed course of time has no significance for so
religious a people as the Indians, whereas in the faith of Zarathustra and in
Christianity it is of quite decisive significance. For the Indians, time was
only resultless movement, change and unrest, from which it was a matter of
freeing oneself. Historical development therefore acquired no value; the
abolition of temporal existence by the attainment of ``Nirvana'' was the
highest thing. A similar trait is found in Greek thought; Plato's doctrine of
ideas especially is a classical example of it: true reality is the
unchangeable, that which always is what it is; the changeable phenomena of
the world of experience, caught up in becoming, are at last only appearance.
In the Zarathustra faith of the Iranians, on the other hand, the development
of the world stands as a historical process which through a comparatively
short span of time (12,000 years) reaches its conclusion, when the great
struggle between good and evil---in nature, in human life, and in the world
of spirits---has been carried to its end. Here the idea of a world-history in
the literal sense of the word appears. All human work and all human conduct
of life thereby come to stand in a mighty, cosmic illumination. The course of
time, surveyable by the imagination, exercises a concentrating influence on
the minds of human beings. There stands a goal before thought, towards which
all striving is directed, and the human being stands as the fellow-worker of
the gods. How and where this faith in the significance of history first arose
is uncertain. It is not certain that it is of Indo-Germanic origin; possibly
it is due to influences from the Semitic or pre-Semitic
side.\textsuperscript{11} It is likewise uncertain how far this faith came
from the Persians to post-exilic Judaism and through this to Christianity.
But the Zarathustra faith stands at any rate as the first great popular
religion in which the idea of life as historical development towards a goal
determines the whole conception. The courage of human beings is strengthened
in the struggle of life by the great visions of the future, by the thought of
a kingdom that is to come and that is to be consummated through a judgment of
the world.

In the New Testament this conception prevails in a pronounced way especially
in the three first Gospels, in the Pauline epistles, and in the Revelation of
John. In the Gospel of John the weight is laid with great emphasis on the
point that eternal life is not merely a future but an already present life.
He who believes has already passed over from death to life. The judgment, the
decision, is first and foremost an inner one; the outer, future decision
comes rather to stand as a concluding consequence. With respect to this
idealism, which already in the present sees the ideal realized, there is
indeed only a difference of degree between the various New Testament
writings; but this difference of degree is of great significance; it grounds
two different types of doctrine, and it is a psychological and
literary-historical riddle that both the writing (``the Revelation of John'')
in which the consummation of ``the kingdom'' is most depicted as future, and
the one (``the Gospel of John'') in which its ideal presence is most
profoundly maintained, could in tradition be ascribed to one and the same
author. The opposition comes forward in the New Testament all the more
strongly because the conclusion and the judgment were expected within the
then living generation, whereby the tension between present and future had to
be sharpened, and the living, forward-directed expectation came to form a
stronger contrast to the inner possession.

In the Zarathustra faith the idealistic conception does not---presumably on
account of the sober and practical character of the Iranians---come forward
so clearly. But in the New Testament there appears on this point---as with
respect to space---a characteristic Both---And. And here there is greater
possibility that a Both---And can be maintained. For there is in itself no
contradiction in the ideal's having taken root and so far being already
present, and its being successively unfolded in its whole extent. The
question---a question we shall later come back to---is whether such a longer
successive unfolding is foreseen and taught in the New Testament. It is,
according to the New Testament conception, only a short time that is to
elapse before the conclusion, and it is generally a comparatively short time
that lies between the creation and the judgment. The whole course of time
lies between these two fixed points; in the interval the great religious
events find their place, and the whole time rounds itself off into a natural
whole for the imagination. There is a beginning, and there is an end, and
there is no dizzying distance between them.---But when the two fixed points
of support are taken away---when it is asked what was before the
``beginning'' and what will be after the ``end''---vividness is abolished,
and there easily sets in a dizziness, just as when it dawns on consciousness
that space can have no absolute boundaries.

\medskip

\noindent\textbf{12.} It follows from the very concept of time that every
moment must lie between two other moments. Thus neither a first nor a last
moment can be thought. Just as the causal principle must find application to
every link in the causal series, and the spatial relation hold for every
place in space, so must the temporal relation too hold for every link in the
temporal series. Every course of time that we represent to ourselves---the
one from ``the creation'' to ``the judgment'' as well---can then only be a
wave in time's great, unsurveyable sea. A completed course of time means a
period that is filled out by the realization of a certain result, a certain
goal. It is the content that makes time be divided into periods. But the
temporal relation spreads itself, by virtue of its own inner law, out beyond
the ``beginning'' and the ``end''; and if our thoughts can give no content,
and therefore no division into periods, beyond the points of support that
have hitherto stood fast for us, then we end at any rate with a question, a
problem, just as with the series of causes and of determinations of place.

That the religious consciousness finds an imperfection in the temporal
relation is easily intelligible. The misfortune of development in time is
always more or less this, that the one period of life comes to stand as a
mere means for the other. Means and end are separated, and life is divided
between joyless labour and idle enjoyment. Just as little as any individual
human being ought to be made a mere means for other human beings, just as
little ought single moments within a human being's life to be mere means for
other moments, past and present mere means for the future. This can be
avoided if the work and the development themselves acquire immediate value
and can thereby themselves become the end or a part of the end. The child is
then not merely a future grown human being; childhood becomes an independent
period of life with its own peculiar tasks and its own distinctive value. And
so must every period of life, every part of the course of time, be conceived.
In this way it becomes possible ``in the midst of time to live in eternity,''
without falling into mystical contemplation. What is external in the
differences of time thereby falls away. ``Eternity'' then stands not as time's
continuation or as a distant time, but as the expression of the continuity of
the valuable through the changes of the times. It then becomes possible to
look time's incompletability and infinity in the face without becoming dizzy
and without receiving the impression of the restless and the purposeless. As
with causal explanation and with the world of space, the only possible
solution of the difficulties in the case of the temporal relation lies in the
direction of an inwardizing conception. It is the inner law and the valuable
content, not the external differences, on which the weight must at last be
laid. And here is something that can subsist independently of the limited
frame to which the religious consciousness, in its mythical and dogmatic
forms, is inclined to cling.

\begin{center}---\end{center}

\section{Concluding Thoughts}
\label{sec:concluding}

\begin{flushright}
\textit{By speculative philosophy one obtains not so much answers to the}\\
\textit{questions that common human understanding poses, as one comes}\\
\textit{to recognize that the questions themselves, being founded on}\\
\textit{concepts that have no truth, must fall away.}\\[2pt]
\textsc{Poul Møller}.
\end{flushright}

\bigskip

\noindent\textbf{13.} The final presupposition of the scientific
understanding of existence is, as I have sought to show, a principle of unity
that lies at the foundation of all connection in the course of time, in
spatial extension, and in the causal series. Can thought now give a closer
determination of, a concluding conceptual form for, this principle of unity?
This question is all the more important because in the reference to a
principle of unity there might seem to be a possibility of a connection
between the religious and the scientific conception of existence.

We are all compelled to bring our knowledge to a close. There are points at
which we in fact cannot get further. But this closing-off may be due to purely
factual circumstances. The one grows tired before the other, or the one has
not so many intellectual requirements as the other. The closing-off is here
different for different individuals, for different ages, different peoples. It
may be limitation in time and space that hinders knowledge from getting
sufficient material. It is not that kind of closing-off I here ask about; I
ask rather about a closing-off that does not rest on special, individual, or
historical circumstances, but on the laws of thought itself.

The principle of unity is the necessary presupposition of existence's
intelligibility. And since the fact that existence is within certain limits
intelligible must be assumed to hang together with existence's own nature, we
have a right to speak of a force or a power which brings it about that
everything that exists, and everything that happens, stands in inner
connection, is held together in a relation of continuity. If one defines God
as the principle of existence's intelligibility, and therefore again as the
principle of existence's unity, then there seems here to be a possibility of a
concept of God harmoniously compatible with scientific knowledge. In the
ideas of God of the popular religions there enters---at any rate at the higher
stages---this determination as a more or less prominent element. But it must
be well noted that in the formation of religious concepts quite other
interests than the purely theoretical and intellectual are at work. Every
religious standpoint gathers up in its concept of God the highest valuable
thing it knows. Ethical and aesthetic interests, but besides and above all the
enthusiasm or feeling of dependence awakened during the struggle of life,
drive the ideas forward to the most inward and highest concentration, which
finds vent in the exclamation: God!---The religious problem arises only
through a spiritual division of labour, by which in particular the
intellectual interest is separated from the other interests and seeks its
satisfaction by the development of special paths and forms of thought.

It is then no wonder that scientific and religious formations of concepts do
not correspond to one another. The philosophical and the religious language do
not stand related to one another in such a way that one could without more ado
translate a thought from the one language into the other, or find, for a
concept formed in the one domain, an entirely corresponding concept in the
other.

Not only for the philosophical consciousness, but for the religious
consciousness too, the great problem of the relation between unity and
multiplicity presents itself. But while philosophical thought holds to unity
as the necessary presupposition of the law-governed connection between the
manifold elements, the religious consciousness in its mythical and dogmatic
forms has a tendency to place unity and multiplicity in an external relation
of opposition to one another, as though they were two different beings or
powers. It thus forms two different concluding concepts, and it calls them:
God and the world; the world is thought as a unity, in that all the
multiplicity is gathered into a whole. Towards this formation of concepts
philosophical thought stands in such a way that it finds one and the same
difficulty in both these concepts, and therefore cannot see anything but that
all that has been achieved is that the problem has been doubled; indeed there
arises thereby into the bargain a third problem, in the relation of the two
beings or powers to one another. If it is so, as I shall now try to show, this
becomes of great importance for religious discussion. For from it follows
that, because a thinker rejects the religious concept of God, one has no right
to presuppose that he merely believes in what in religious language is called
``the world.'' The concept ``world'' contains philosophical difficulties
similar to those the concept ``God'' contains, when it is to be absolutely
concluding. And often it is the case that thinkers who are regarded as
atheists reject the current concept of ``the world'' rather than the concept
of God. This holds of philosophers like Spinoza and
Fichte.\textsuperscript{12}---Polemical fanaticism naturally does not let
itself be corrected, however often this state of affairs be pointed out.

If God and the world are regarded as two different beings or powers, then they
must limit one another. The one ceases where the other begins. The clockmaker
and the clock he has made are two different things; in a similar external
relation the customary conception places God and the world. But then God
cannot be an infinite being. ``The world'' with its laws and forces is the
deity's barrier. It is properly a polytheistic conception that arises in this
way; ``the world''---as different from ``God''---properly becomes itself a
god, and a certain personification of the world then also appears
involuntarily, even within views that with great pomp proclaim themselves
monotheistic.

\medskip

\noindent\textbf{14.} The concept ``world'' is, as already remarked earlier,
not a legitimate formation of concepts.
% LACUNA IN THE E-BOOK. The 2020 setting runs "14. som allerede tidligere
% bemærket, ikke en berettiget Begrebsdannelse." — the subject has dropped out
% between the section number and "som", leaving a headless sentence. The sense
% is fixed by the back-reference ("as already remarked earlier") to §7's
% "Begrebet 'Verden' er derfor egentligt et falsk Begreb" and by the paragraph
% that follows, which is about that concept throughout; supplied accordingly.
% Check against the KB scan if a printed copy is ever to hand.
If this concept is to denote a whole, a law-bound totality, then it is an
ideal towards which inquiry steers, but which on account of the
inexhaustibility of experience it will never be able completely to satisfy.
Our thought is always spelling out what reality is. The mark by which a
distinction can be drawn between merely subjective idea and objective reality
is the law-governed connection that can be pointed out among the given
phenomena. But partly there are always arising new phenomena as links in
experience, and partly there show themselves to obtain, among the phenomena
already given, far more complicated and mysterious relations of connection
than we could at first suppose. We therefore never finish applying our
criterion of reality, and the concept of reality therefore becomes an ideal
concept.\textsuperscript{13} From the religious side one has as a rule not the
patience to be content with the progressive work into which this ideal leads
inquiry, namely to find ever deeper connection among ever more phenomena. One
breaks off the work and regards the structure as completed. The concept
``world'' is therefore the expression of a half-thought. At the lowest stage
of religious development this concept is not yet found. The idea of a cosmos,
an ordered whole of existing things, presupposes a certain degree of thought.
By a precipitate closing-off there has then been formed the pseudo-concept of
an absolute whole, and after it had been introduced into religious language it
has figured there ever since, as though it were a good and defensible concept.

But let us for a moment acknowledge this formation of concepts. The question
then becomes that of the relation between the two beings or powers. Everywhere
where there is a causal relation there must---as we have seen in what
precedes---be presupposed a principle that conditions the connection between
the links. If there is nothing that brings them into connection, the causal
relation is unintelligible. This must hold also for the causal relation
between God and the world. But then the bond, the principle of unity, which
makes their mutual relation possible, would itself be the real deity---or, if
one prefers, the real world would then consist of God + ``the world,'' that is,
of the whole that the two together form. We then come again into an
incompletable series---a sign that one has been at the last thought we can
really form, but has gone past it. This last valid thought is the principle of
unity, the principle of conformity to law in existence. If one does not heed
it, it avenges itself by reporting again---like the goblin who moves with
you---every time one supposes oneself to have made an advance. There is
something teasing about all limit-concepts, which comes out especially when
one wants in an external manner to tie a knot in the thread of thought instead
of holding to the inner connection that makes the thread a thread. Only if
there could be given us from the theological side a wholly different
epistemology from the one we must now work with would this teasing quality of
our limit-concepts possibly be able to fall away. As it is, it follows us right
up into the highest theological and metaphysical speculations. Plato already
has (in the dialogue \emph{Parmenides)} seen in this teasing quality a
difficulty for his doctrine of ideas: if it is the ideas that lie at the
foundation of sensible things, then there arises a new question with respect
to the relation between the ideas and the things; this relation too must have
its idea, and so forth. With Spinoza and Lotze the problem comes forward in
the relation between substance and the individual beings (\textit{modi}). I
cannot see that theology, by its doctrine of God and the world, helps us
beyond this difficulty of principle, which epistemologically is an indirect
proof that we stand at a limit-concept. In the principle of unity we have a
thought that is concluding \emph{for us}, but at the same time a thought that
demands ever new applications in incompletable series and series of series.
One is deceived if one would find the concluding thing as a link in a series,
instead of in the principle of series-formation itself.

The possibility cannot be dismissed that the incompletability of experience
and of knowledge hangs together with the fact that existence itself is not
finished, but is in constant development. One easily forgets that experience
and knowledge themselves also belong to existence. When we speak of knowing
existence, it properly means that the whole of existence is to be expressed by
a single part of it or a single side of it. So long as this task is not
solved, existence itself as a whole is not finished. And conversely, if
existence itself is not finished, our knowledge cannot be finished either.
Besides: to know the relation between the part or side of existence we call
knowledge, and the rest of existence or existence as a whole, there would be
needed a new knowledge---and so forth. Seen from this side too, the problem
presents itself ever anew.

It is on the whole a fundamental law for all our concepts that they express
relations, and that every concept presupposes a relation to other concepts. It
is the inner law of all movement of thought. From this law, which may be
called the law of relation, it follows that no concept can be formed of
anything that does not stand in relation to something else. All movement of
thought consists in pointing out likeness or difference, in finding ground or
consequence, in deriving cause or effect. But likeness and difference, ground
and consequence, cause and effect find application only where different links
can be placed together in mutual relation.\textsuperscript{14}

Every formation of a concept is a limitation, which is again abolished when
thought reflects on the larger connection out of which it has taken that which
is gathered up in the concept---and without which the concept's content could
not be determined. The law of relation itself testifies to thought's unity, in
that different terms of a relation are gathered up in one concept. But it
testifies also to thought's constant limitation, since besides the terms that
are gathered up in the single concept there are always presupposed other
terms, whose closer determination becomes a new task. It is the nobility of
our thought that it can see its own limit and at the limit always
hear---by virtue of its own law---an \textit{Excelsior!} Every concept of God,
Fichte has said, is a concept of an idol. But it is precisely a spark of
divinity in human thought that, in the face of all pious attempts to formulate
the divine, it can recognize that they are idolatry.

\begin{center}---\end{center}

\section{Thought and Image}
\label{sec:thought-image}

\begin{flushright}
\textit{Und deines Geistes höchster Feuerflug}\\
\textit{Hat schon am Gleichniss, hat am Bild genug.}\\[2pt]
\textsc{Goethe}.
\end{flushright}

\bigskip

\noindent\textbf{15.} We do not cease to form ideas, even when we have
reached the limit of all knowledge, so that clear and contradiction-free
concepts can no longer be formed. The religious need in particular will
constantly lead to such formation of ideas. It will appear on closer
investigation that the arising of all such ideas is due to
\emph{analogy}.

The word analogy properly denotes the same as proportion, but one is
accustomed to distinguish between the uses of the two words, so that
proportion indicates a quantitative likeness of relation (e.g.\ $1/2 = 2/4 =
3/6$, and so on), while analogy denotes a qualitative likeness of relation
(e.g.\ red stands to orange as orange to yellow). By the help of proportion
we can discover unknown magnitudes, if only we know that they are terms in a
proportion whose other terms are known to us. Analogy, on the other hand, is
a highly imperfect aid to our knowledge. It cannot give us any positive
knowledge of the unknown terms, but serves especially to express the relation
between terms we know beforehand. It moreover passes easily from being a
conceptual relation over to being a poetic juxtaposition, and it moves
generally on the boundary between thought and poetry.\textsuperscript{15}

In my work \emph{Den nyere Filosofis Historie} I have sought to show that
those philosophers who do not concede the correctness of the way in which
critical philosophy marks off the limits of knowledge constantly---with a
greater or lesser degree of explicit consciousness---operate with analogies
in their attempts to solve the riddles of existence. It is the only aid that
is left when experience refuses its service. One then conceives existence in
its wholeness, or existence's ultimate ground, in analogy with certain
elements, phenomena, or relations within the domain of experience. The
character of the different systems, of the different metaphysical world-views,
rests essentially on \emph{which} elements, phenomena, or relations one lays
at the foundation of one's analogizing.

It is often supposed that there is a forced choice between the two
possibilities, that existence's inmost essence is of a spiritual kind and that
it is of a material kind. But the division of the existent into spiritual and
material is purely empirical. No proof can be given that existence should not
be able to appear under other forms than these, so that the relation between
these two forms of existence should be an Either---Or (what in logic is called
a contradictory relation). Perhaps the great difficulty that shows itself to
be bound up with reaching a satisfactory hypothesis concerning the relation
between the spiritual and the material hangs together with the fact that
existence is not exhausted in those two forms of existence, but that there are
many---not to say with Spinoza infinitely many---other forms of existence. It
might be that the relation between the two forms of existence we know would
only be able to be understood if we knew several other forms of existence. We
have at any rate not all the data for the solution of the problem of
existence. Let us suppose that we knew only two colours. We should then be
inclined to believe that all colour that existed at all must be one of these
two. Now we know that light can be refracted in more than two ways---and yet
the physics, physiology, and psychology of the colours are so extraordinarily
rich in problems! Existence's light can surely be refracted in many more ways
than the metaphysical systematizers have imagined. This belief (and this
analogy) is at any rate as justified as that on which the idealistic world-view
supports itself.\textsuperscript{16}

When one does not take the probative force of analogies so strictly, it is
easy enough to speculate. But one then puts images in the place of concepts,
poetry instead of philosophy. Epistemologically it is of great importance that
a sharp distinction be drawn between concept and image. The image may have its
great value; but when it does not stand as an example or as a vivid sign for a
concept, the value the image has rests on its relation to other sides of our
nature than the cognitive.

\medskip

\noindent\textbf{16.} Only reluctantly does human consciousness give up
immediate intuition, as this appears in the form of sensation, memory, and
imagination, in order instead to travel the road of analysis and form
concepts, which can with difficulty fully make up for what was had in
intuition. But it is once for all the case that the objects of the religious
consciousness---when this consciousness has reached a higher stage of
development---can no longer be apprehended by immediate intuition. The
involuntary formation of images must constantly be altered and corrected
anew, and it shows itself at last insufficient to render the eternal,
infinite, and exalted essence of the religious objects. The religious
consciousness here oscillates between two tendencies. On the one hand it
scruples to use sensuous and human expressions of its infinite object; on the
other hand it scruples to deprive immediate feeling of the living and gripping
expressions in which it involuntarily finds vent. In a religious genius like
\textsc{Augustine}, who was at once thinker and prince of the Church, one can
study both tendencies in their mutual conflict. For him (as later for
Schleiermacher) the expression ``mercy,'' used of God, has only figurative
significance, because it presupposes a suffering at the suffering of others.
But in order that the ``souls of the unlearned'' (\textit{animæ indoctorum})
should not be hardened, Augustine nevertheless thinks that one ought to use
it---thus as a kind of pedagogical anthropomorphism. Only those, he adds, who
combine religion and inquiry (\textit{religio et studium}) can dispense with
the figurative expressions.\textsuperscript{17}---But how far now does the
figurative extend? Is there properly any expression whatever for religion's
objects that is not figurative?---If we look at religion's history, we find
especially two interesting processes of development which aim at eliminating
everything figurative, and which culminate in the rejection of all definite
expressions, because these are drawn from limited domains of experience and
therefore make finite that of which they are used.

The first process of development lies before us in the Indian
\emph{Upanishads.} The Vedanta philosophers came to the conviction that no
concept and no image could express the deity's essence. They set up their
metaphysical idealism only in order at last to take it back. The deity is
``without end, without age, without shore, not outside and not inside.'' It is
best denoted in a negative manner (\textit{neti, neti}): ``it is neither this
nor that!'' And silence is the right answer when it is asked what it is. If it
could be beheld by us, it would not be what it is. It is the unthinkable---and
yet that by which thought thinks! No twoness, no opposition has validity for
it; therefore it cannot be known. Neither sight, word, nor thought can
penetrate to it, though it works in sight, word, and thought alike. It is
neither being nor non-being; it is neither known nor not-known. Only
symbolically can it be expressed. But in order to keep off all sensuous ideas,
which are easily brought on by symbols that have a definite meaning, the
Vedanta philosophers used as an expression for the highest the syllable
\emph{Om}, which means ``yes,'' but does not lead the imagination in any
definite direction.\textsuperscript{18}

The second process of development takes place in \emph{Neoplatonism} and is
propagated from this (with the so-called Dionysius the Areopagite as
intermediary) to \emph{Christian mysticism} in the Middle Ages. Although more
recent investigations have shown that there is not the complete opposition
between Scholasticism and Mysticism that was earlier assumed, there is
nevertheless one point in which these two medieval tendencies diverge
decidedly from one another, and that is precisely the question whether there
are analogies that can rightly be used to think God. The Schoolmen held fast
to analogy's justification; the Mystics rejected it. \textsc{Thomas Aquinas}
teaches, indeed, that it is impossible to use any expression of God quite in
the same sense (\textit{univoce}) as that in which it is used of created
beings; but there are said to be expressions that can be used of God not
merely by a play on words (\textit{pure æqvivoce}). There is said to be a kind
of analogy that lies midway between univocity and mere likeness of word. It is
that which rests on a causal relation, as when, for example, both a medicine
is called healthy, and the one who is cured by it is said to become healthy. A
relation of analogy of this kind is said to obtain between the Creator and
created beings. Yet it is not a relation of likeness that could make it
justified to bring the Creator and the created under the same generic concept.
No, every determination---even the concept of being itself---is used in a
different sense of the Creator and of the created: God stands outside every
genus (\textit{extra omne genus}).---Very clear the knowledge that is thus
supposed to be possible evidently cannot be, and it is not easy to form any
idea of it. For analogy is (however much one fences it about) a relation of
likeness, and in all comparison there must be able to be a question of generic
and specific concepts: juxtaposition by the help of analogy is in reality the
formation of a generic concept. When Thomas Aquinas names, as an example of
theologically applicable analogy, the one that rests on causal connection, it
must be remarked that the more heterogeneous, indeed generically different,
the cause is from the effect, the less do we understand the causal relation.
Where the causal relation is not purely external and purely factually imposed,
there is a continuity and thereby a likeness between cause and effect which
makes it justified to subsume them both under the same generic concept. Such a
connection one will also on closer investigation find between the medicine and
the state it produces in the organism.---The Scholastic assertion of analogy's
justification therefore passes consistently over into the mystical denial of
this justification. \textsc{Hugh of St.\ Victor} already taught that God is
better thought by negative than by positive expressions. Later \textsc{Suso}
called God ``\textit{eine namenlose Nichtigkeit}''---``a nothing of all the
things one can think or name'' (though in himself he is ``a most essential
something''). And the author of \emph{Deutsche Theologia} says: ``If the
perfect, sole Good [God] were something, this or that, which the creature
understood, then it were not all, nor yet the only one, and it would then not
be perfect either. Therefore one calls it Nothing. One means thereby that it is
nothing of that which the creature can comprehend, know, think, or
name.''\textsuperscript{19}---Mysticism was at once a movement of feeling in
the direction of inwardness and an intellectual movement in the direction from
image to thought and from thought to thought's limit. In that one experienced,
in the strong movements of feeling, states in which no definite ideas could
form, one was prepared to exercise sharp criticism of every attempt to express
the highest in the form of a concept. Mysticism here meets with critical
philosophy, which maintains the insufficiency of ideas in the face of that
which does not appear in experience's limited form. Already in antiquity we
find this kinship between criticism and mysticism, for there is a certain
likeness between Carneades and Plotinus in epistemological philosophy of
religion.

\medskip

\noindent\textbf{17.} But neither Scholasticism nor Protestant theology will
allow Mysticism and Criticism to be right in holding that all religious ideas
are figurative. In a certain sense it is also a question of life or death. For
the concept of revelation (in the strictest sense) would fall away if at
certain points the difference between image and reality did not fall away.
Otherwise the old problem of the relation between knowledge and things
themselves would at last arise again, and the full conclusion and rest be
excluded.---That expressions such as Lord or Father, used of God as absolute
being, are figurative, and that the similitudes underlying these expressions
fail at definite points, could not in the long run escape reflection. The Lord
is Lord only in relation to slaves or servants. The word God, said Newton
therefore, is a relative word (\textit{Deus est vox relativa et ad servos
refertur}). And the concept of Father presupposes not only the concept of
child but also the concept of mother, although in the history of the popular
religions there shows itself a tendency to push out this latter
concept.\textsuperscript{20} And yet both the inwardness of feeling and the
need of the imagination lead again and again back to the old images, even
though one gradually acknowledges their symbolic character more readily, at
any rate in theory.

When reflection develops, one passes from the more vivid and concrete
designations to the more abstract. Such a more abstract designation is
``personality.'' One gives up the more definite characterization for a more
indefinite one, which yet is also drawn from human relations.

Psychologically regarded, the determination of what the concept of personality
contains will always come back to two chief traits.---Personality (the self,
the I), as we know it, is in the first place characterized by its unity. All
elements in its being are gathered up and connected with greater or lesser
energy, and that in a more intimate manner than the elements of matter are
connected according to the laws of outer nature. The very elements that are
connected in personality's unity are not absolutely produced by it; here it is
more or less dependent on the conditions of life. New elements may crop up in
it without its wish, though not without its own unconscious or involuntary
co-operation, and it then depends on how far and in what way they are brought
into connection with the elements that are already worked together.---In the
second place, in every personality's content there will be some elements that
occupy a more central place and have a more constant character than others.
Such central and constant elements are the ruling ends and interests, which
express the will's kind and direction. This is the real side of personality,
which especially conditions the individual's special distinctiveness.%
\textsuperscript{21}---Seen both from the formal and from the real side,
personal life as we know it presents an alternating relation of activity and
passivity. Activity shows itself in the working-up of a content received with
preponderant passivity. During this working-up the peripheral and shifting
elements are more or less energetically expressed and mastered by that which
in the individual personality (and in the individual period of life) is the
central and constant. We know personality only as striving and struggling, as
asserting itself against the resistance and the difficulties caused partly by
lack or superfluity of content, partly by conflict among the elements
themselves. A personality that produces all its content absolutely ``out of
itself''---that has no ends still unaccomplished---and that has no real
resistance to overcome---is not what we understand by a personality, when we
hold to psychological experience.

The dispute as to how far the deity (as this is conceived at higher religious
standpoints) can be called a personal being---a dispute often designated as
the dispute between theism and pantheism---rests on the fact that from the one
side one idealizes and extends the concept of personality by taking away the
moment of passivity, while from the other side one denies the justification of
calling a being a personality when it itself produces all its elements and
therefore knows no real resistance, so that it cannot struggle, cannot strive,
cannot hope.

Personal life is the highest form of existence that experience shows us. And
the unity we are always at last led to assume in existence on account of its
continuity and conformity to law recalls the unity in our consciousness, the
formal side of personality. But for all that it is only an analogy, and one
that fails at the most essential point, when one uses our idea of the limited
I, which is a single link within the great world-order, to express the
incompletable principle that manifests itself in the very fact that there is a
world-order. And it becomes so much the more difficult to hold fast to this
analogy, the greater the inclination is to think the deity as perfect, as
``absolute'' in the sense of finished and concluded, so that it (as Andreas
Sunesøn says in his Scholastic didactic poem) can receive no increase
(\textit{natura dei, cui nil accrescere posset}). By the dogmatic
absolutizing and the strong emphasis on unchangeableness, the difficulties in
holding fast to the analogy with the concept of personality are sharpened and
made into downright self-contradictions.

The difference between a theistic and a pantheistic conception rests, in the
individual case, not merely on epistemological arguments. Involuntarily other
motives play in as well, and to a great extent the difference rests on a
difference in the strength with which the need to form vivid images,
especially in the boundary regions, stirs in different individuals. From this
side we shall later, in the psychological part of the philosophy of religion,
come again to touch on this point.---As regards the expression pantheism, it
must be remarked that it is used in a somewhat different manner. According to
the one usage (which is maintained as the right one by an authority like
\textsc{Eduard Zeller}) pantheism denotes a world-view that assumes an inner
(immanent) relation between God and the world, whether one then conceives God
as personal or not. According to this usage a standpoint like Lotze's would
just as well be called pantheistic as Spinoza's standpoint. According to the
other usage (which I follow) the concept has a narrower extent, in that an
immanent world-view is called pantheistic only when it does not ascribe
personality to the principle of unity in existence (except in poetic
speech).---

To go to still more abstract expressions like force, life, substance (being)
does not help; they too denote relations and thus presuppose something to
which one stands in relation. Such expressions too must, when they are to be
applied in the way the religious consciousness in its mythical and dogmatic
forms wishes, undergo such an alteration that nothing but the bare word is
left.---

It shows itself then undeniably that all analogies fail
at the decisive point. We stand at thought's limit, and it is no wonder that
real concepts can no longer be formed, and that the images may indeed satisfy
the individual mood, but do not permit of any consistent application. This
unknowable is strongly emphasized by so-called \emph{agnosticism.} But this
tendency has an inclination to think that ``the unknowable'' is wholly
different from everything that experience shows, so that none of experience's
laws should have any significance whatever for that which is nevertheless
assumed to be existence's principle. This is a dualistic assumption that
cannot be grounded. Although our finite points of view and forms of thought are
unable to give a determination of it, let alone an exhaustive formulation of
it, they must nevertheless stand in inner connection with it, and we have no
right to declare them meaningless. Our thought, our knowledge, belongs once
for all to existence, has arisen within it, and is itself one of the facts that
must be taken into consideration before we halt our inquiry into existence's
nature. The ideas our experience leads us to develop are indeed probably very
subordinate and derived in relation to the great connection into which we
belong as single links; but meaningless they do not on that account become. Only
we are unable to state what degree and kind of metamorphosis our points of view
would have to undergo if they were to be taken up into the highest, most
comprehensive connection that we could think.

The expression ``agnosticism'' was in its time formed by \textsc{Huxley}, so
that he might have a designation ending in \emph{-ism} like all other decent
expressions for views. Were I to find such a designation for the conception I
strive to make good, I should prefer to call it \emph{critical monism.}

\medskip

\noindent\textbf{18.} In order that concepts, analogies, and images should be
able to be used to express that which at higher stages of development is
religion's object, they must be subjected to an idealization and sublimation
which for religious feeling easily brings coldness and remoteness with it. The
religious consciousness has, to be sure, a tendency to represent its object as
exalted as possible, to raise it far out beyond all finite relations. But by
merely following this tendency it works against itself, since the inner, living
relation to its object is made impossible. Of the Hindus an observer reports:
``I have often asked: Where is the temple of the highest God? And the answer,
plainly enough given with great astonishment at so unexpected a question, has
always been: A temple for the highest God! What is that to mean? There is no
such temple, because he can have none. He is without shape and without name,
inconceivable, and we cannot worship him. We need a visible object in which our
spirit can find rest.''\textsuperscript{22} Within Christianity we find again
the same two tendencies. In the Catholic Church there is assumed a whole
graduated series of beings to whom one can turn and get one's need of a finite,
limited object for prayer and devotion satisfied, if the highest link in the
series seems to us to be too remote and exalted. From God the series goes
through Jesus and the Madonna down to the special saint who is active in the
special kind of case or takes particular care of the individual person. In
Protestantism a limit has been set to this tendency, but it has not been
abolished. In a Copenhagen church it was once said in a funeral oration that
God cannot help us in our great sorrows, ``for he is so infinitely far away'';
therefore we must hold to Jesus.

Just as the religious consciousness oscillates between literal and figurative
meaning of the expressions it uses, so it oscillates between the tendency to
exalt and infinitize its object and the tendency to draw the object down into
the finite world as fellow-combatant and fellow-sufferer. Often the strongest
feeling, the most burning devotion, turns towards the struggling and suffering
God, not towards the eternal and perfect one. If now both tendencies are held
fast at once in their extreme forms, there appears \emph{the religious
paradox:} God is unchangeable and changeable, eternal and becoming,
victorious and beaten, blessed and suffering! Often pronouncedly religious
spirits heap up such oppositions by preference, sharpening the contrasts to
the utmost, partly in ecstatic admiration of the great, partly with scorn
towards the thought that vainly seeks to grasp it. The chief opposition to
which everything here leads back is the opposition between unchangeableness
and changeableness. The typical difference of the two greatest popular
religions rests on their relation to this opposition. And the opposition
recurs within one and the same religion. Even he who believes in a God ``with
whom there is no variation or shadow of turning'' may yet come to the result
that ``God is the most changeable of all beings.''\textsuperscript{23} The
paradox does not first come forward with special dogmas, e.g.\ the dogma of
the God-man, as \textsc{S. Kierkegaard} has sharply pointed out in his
\emph{Filosofiske Smuler.} The paradox lies properly already in every
anthropomorphism, in every image that is to be more than an image. For the
assertion that any image whatever suffices is an assertion that the unlimited
is limited, the infinite finite. One easily forgets the paradox through the
power of habituation, when one has let one's imagination settle itself
comfortably in the circle of images one has formed for oneself or lived into
by the help of tradition. But as soon as the surge of feeling or the speed of
thought becomes stronger, the limited form is burst; and when it is
nevertheless to be held fast, or rather, when one nevertheless \emph{will}
hold it fast, one approaches a \emph{credo qvia absurdum.}

\medskip

\noindent\textbf{19.} A Both---And has on this point too shown itself to be
distinctive of the religious consciousness. But where this Both---And sharpens
into the paradox, it becomes clear that the interest which has led to the
development of the religious ideas in their pronounced form cannot have been a
purely intellectual interest, not the mere need to understand. In every
religious standpoint, especially in the great popular religions, knowledge
does to be sure play a part. Religion is in its classical times everything for
the human being and satisfies his need for knowledge too. But the need for
knowledge is here nevertheless subordinate to the need for life, the need to
develop and preserve life. ``I seek thee, that my soul may live,'' says
Augustine to his God. Only thereby is it explained that ideas can be developed
and held fast which defy the intellectual interest in clear,
contradiction-free thoughts. We are led to see that the foundation of religion
must be sought in other sides of the life of the spirit than in the life of
pure thought.

Religion cannot be made or constructed. It grows out of life, springs from the
human being's fundamental mood during life's struggle, from his will to hold
fast in spite of everything to the validity of the highest valuable thing that
his experiences have led him to know.

The hypothesis that religion is essentially faith in the persistence of value
presents itself here naturally, after it has shown itself that religion's
essence cannot be to give an understanding or an explanation of existence such
as the intellectual interest demands. The riddles that science cannot solve,
religion does not solve either. It has been shown that it can give no
explanation of the special events (section A)---that its ideas are not fitted
for the conclusion of scientific thought (section B)---and that these ideas
have the character of the image, not of the concept (section C). If the
religious ideas are to have any significance, it can then only be that of
serving as symbolic expression for the human being's mood, longing, and wishes
during the struggle of life. They are thus of secondary, not of primary,
significance and origin. The philosophy of religion here combats dogmatics in
a manner similar to that in which art history combats the mythology of art.
What drives the artist to shape a figure is---the art historian
maintains\textsuperscript{24}---the living joy in figures that life has shown
him; he wishes now to give them new birth, perhaps in a more ideal form. But
what he \emph{calls} the figure he shapes is a subordinate question, which is
often only settled long afterwards.

The intellectual interest will naturally, once it has developed, not be able
to refrain from influencing the religious interest. If religion is essentially
faith in the persistence of value, then this faith constantly needs ideas in
order to become clear about what it means; these ideas will---perhaps slowly,
but yet steadily---be worked together towards greater consistency and towards
greater agreement with the ideas the human being comes to form by other roads
and under the influence of other motives. The question then becomes how the
two interests really stand related to one another. The intellectual interest
leads one to conceive existence as a great, unsurveyable system of causal
groups and causal series; the religious interest leads one to conceive
existence as a home for the development and preservation of values. Can these
two points of view be united? How does value stand related to causality, the
riddle of life to the riddle of the world?---Perhaps this problem is
insoluble, and its solution would be the finding of the philosopher's stone.
But the epistemological philosophy of religion has nevertheless, it seems,
helped us to see more clearly where the problem lies. It is indeed
predominantly by an indirect, negative road that epistemology can illuminate
the religious problem, but this road too must be travelled if clarity is to be
reached. We turn now to a direct study of religion's essence by the help of
psychology. What matters is how far the hypothesis of religion as faith in the
persistence of value, to which I have several times pointed, can receive a
satisfactory grounding and confirmation.

\begin{center}---\end{center}

%% ============================================================
%%  III. PSYCHOLOGICAL PHILOSOPHY OF RELIGION
%% ============================================================
\chapter{Psychological Philosophy of Religion}
\label{ch:psychology}

% Chapter III opens with an untitled introductory section (§20) standing
% between the chapter head and section A. It is a separate chunk in the
% e-book (s011-Chapter-005.xhtml); the skeleton had jumped straight to A.

\noindent\textbf{20.} Properly we have already in the preceding section made
use of psychology, since there is a constant reciprocal relation between
epistemology and psychology. Epistemology investigates the forms and elements
of our knowledge with respect to the question of their usability for
understanding existence, while psychology investigates them with respect to
their factual character and their factual coming-into-being, whether they have
any usability and validity or not. In the case of the religious ideas it has
now shown itself that their significance cannot be that of making existence
intelligible to us (in the sense in which our intellectual interest takes the
word understanding). Then the psychological investigation alone stands here
remaining. It may possibly lead us to see what positive significance the
religious ideas do have.

The psychology of religion is a part of general psychology, a speciality
within it. It stands in a reciprocal relation to general psychology.---It is a
means for the latter, in that within its own distinctive domain it collects
and works up material of great psychological interest. In the religious domain
human beings have made some of their deepest, most intensive experiences of
the soul. In religion (if it is genuine and original) all the elements of the
life of the soul work with an energy and an interplay such as scarcely in any
other domain. Therefore the study of the religious life of the soul is of such
great importance for psychology as a whole. It is not merely the history of
religion, but also---and quite as much---the individual religious life that
draws the psychologist's attention to itself. In the history of religion one
does not so easily get to see the forces of the soul in their original working.
It occupies itself predominantly with the great types, with the forms of
religion common to great groups of human beings. For the psychologist, good
biographies, and especially autobiographies, of religious personalities are
more instructive than the most learned works on the great popular religions.
Thus---to name a few single examples---Augustine's \emph{Confessiones}, Suso's
and Saint Teresa's autobiographies belong to the most important material of the
psychology of religion. The history of individual religiosity naturally also
belongs under the history of religion; but this has hitherto been
predominantly occupied with studying the great types.---But the psychology of
religion not merely furnishes material for general psychology. Conversely it
also uses the latter's points of view and laws in its own special domain. It
seeks to understand the phenomena of religious life by a purely psychological
road. Here general psychology is a means for the psychology of religion. And it
is especially from this last side that we must regard the psychology of
religion here, where it appears as a part of the philosophy of religion. It
uses the general psychological methods, so far as the subject permits, and
rests on the same presuppositions as general psychology.

\begin{center}---\end{center}

\section{Religious Experience and Religious Faith}
\label{sec:experience-faith}

\begin{flushright}
\textit{Qvid est qvod amo, quum te amo?}\\[2pt]
\textsc{Augustine}.\\[2pt]
\textit{(What is it that I love, when I love thee?)}
\end{flushright}

\bigskip

\subsection{Religious Experience}
\label{ssec:experience}

\noindent\textbf{21.} The word experience is here taken in the sense of
lived experience, the finding of states, as opposed to the working-up by
thought. It will first and foremost be a matter of becoming clear about what
and how much can be reckoned as religious experience, and what limits this
experience is by its nature subject to. I attempt in what follows to
characterize religious experience from a purely psychological point of view.

\medskip

\noindent\textbf{22.} Experience in the sense of lived experience cannot be
completed so long as life lasts. It can present a series of individual states,
but no completed totality, and no general proposition can appear as a fully
confirmed result of experience. Totalities and generalizations are due either
to a thinking working-up of what has been experienced, or else to an
involuntary emphasizing of a single experience or kind of experience as
decisive and all-determining. Prior to the arising of a conscious reflection it
will especially be this involuntary emphasizing that lies at the foundation of
what comes to stand as the results of religious experience. It lies in the
nature of feeling that, when it has been awakened by a single event, it seeks
to spread itself over the whole life of consciousness and to colour all other
elements in this, whether they stand in connection with that event or not. By
such a \emph{feeling-expansion} (thus I have in my \emph{Psychology} named
this phenomenon) inner states which in and for themselves have nothing to do
with one another can acquire a common stamp, and a whole can come out of the
series of inner lived experiences, although this series has not in itself
reached a real conclusion. For different individuals and under different
circumstances it may be highly different lived experiences that thus come to
determine the mood of a whole inner life. However unavoidable and however
significant this feeling-expansion may be, it nevertheless contains sources of
error that one must be attentive to when one proceeds to a reckoning-up of what
experience and lived experience have really yielded.---Naturally such a
feeling-expansion is itself a lived experience; but it is derived in relation
to the lived experiences whose result it is a matter of establishing.

In the debate between optimism and pessimism, for example, one will find that
single lived experiences, which are assumed to be typical, determine the
conception that is maintained. Although from both sides one appeals to
observation and thought, the result one reaches is determined not so much by
logical generalization or induction as by a feeling-expansion. The endeavour
aims at holding the underlying lived experiences fast in their freshness and
letting them keep their place at the centre of consciousness, fixing attention
on them; then the rest follows of itself.

If faith in the persistence of value is the core of all religion, religion will
not be able to be built on immediate lived experience. Immediately we come to
know only the individual, definite values, as they appear according to our
human and individual nature and under our special conditions of life; but about
their persistence no experience can immediately instruct us, any more than it
can teach us that what has the highest value for us is also the central thing
in the whole of existence. What is experienced and lived through may be a
\emph{motive} for believing in the persistence of value, but not the
\emph{content} of this faith itself. Personal life, especially as it expresses
itself in the great breakthroughs where new directions of life are introduced
and new forms of life are born, is for us the highest valuable thing we know.
The experiences we make of it we therefore use involuntarily to illuminate the
whole of existence for us; it is to us as though existence in such
breakthroughs revealed its hidden forces. It is a feeling-expansion that here
operates. If we try to transpose it into the form of thought, it is an
inference from analogy that we get---and an analogy that lies out beyond the
limit where thought passes over into poetry.

\medskip

\noindent\textbf{23.} In all experience a distinction must be drawn between
what is immediately given, what serves to explain it, and what serves to
express it. These three moments, the state itself, its cause, and its
expression, can lie so close to one another that they cannot be separated
except by closer and more exact reflection. The relation between them is a
similar one to that between a discovery, the proof of what has been discovered,
and its formulation. These three moments too can approach one another, though
it may be of great interest to keep them apart.

When one appeals to one's experience, one can provisionally support oneself
only on \emph{the lived state itself.} That this state may have \emph{a
definite cause} is not lived through in so immediate a manner as the state
itself is lived through. Precisely when something has intervened strongly in
our life of the soul and called forth great effects there, we shall be so taken
up with these effects, which are indeed the very thing we immediately live
through, that it will be impossible for us to form a clear and exact idea of
the state's cause. The real cause always consists of a plurality of conditions
and can only be established by critical investigation. It always rests on an
illusion when one supposes oneself immediately to experience ``the cause as
cause.'' That a causal relation is present one cannot immediately observe;
immediate observation always shows us only that the one follows \emph{after}
the other, not that it follows \emph{from} it. The presence of a causal
relation is established only when it can be shown that \emph{only} when this
definite event goes before does this definite state \emph{unavoidably} follow
after. But neither an \emph{only} nor an \emph{unavoidably} can be known quite
immediately; they must be ascertained by comparison or by experiment. In
doubtful cases it is also comparison and experiment that we resort to, and it
is indeed in doubtful cases that the appeal to ``experience'' especially
appears.

We do not begin with doubt, either in the domain of inner or of outer
experience. Every rising element of consciousness is from the first taken
involuntarily as valid, as the expression of reality. It presents itself with a
quality of existence that is shaken only when contradictory elements too
present themselves with their quality of existence. If such contradiction fails
to appear, no doubt arises, and no criterion of reality or validity is needed.
But when it does set in, the question becomes which of these elements possesses
the true reality, and this can only be decided by there being found an
unavoidable connection between one element and other elements whose validity
stands fast beforehand. Thus we proceed in doubt about something in the world
around us. The difference between a hallucination and a real event is
discovered by the fact that the latter, but not the former, can be inserted as
a link in a great, connected causal series that is tied to elements whose
validity is not drawn into doubt. The more firmly connected a dream is, the
more it stands for us as reality; and the more isolated a real event stands in
our experience, the more easily can we take it for a dream. What thus holds of
outer experience holds also of inner. Here too an immediate experience of a
causal relation is impossible. In every more deeply penetrating discussion of
the significance of religious experience it will at any rate show itself
expedient, if not necessary, to get the immediately lived experience itself set
out with the greatest possible abstraction from the supposed cause. The
strictly immediate lived experience itself cannot be any illusion. The
illusory arises precisely only through an incorrect causal explanation which is
confused with immediate observation. How causal explanation comes about at all
we shall later have occasion to discuss.

If there is present a strongly agitated state or lived experience, it will not
be exactly the need to find its cause that first of all announces itself, but
rather \emph{the need to find vent in outbursts, in exclamations, in words, or
in action.} The individual then seeks after thoughts or images that can satisfy
this need and give it a definite direction. It is a need for reaction that here
expresses itself---a need for reaction in fear or in thanks, in love or in
indignation and hatred. The manner in which this reaction takes place will be
determined by earlier lived experiences and by the tradition under whose
influence the human being stands. Only in especially original natures does a
new formation here take place. It is such natures that become prophets or
founders of religions.

By a natural displacement, that which satisfies the need for expression and
reaction easily coincides with that which satisfies the need for a causal
explanation, as soon as this announces itself. In the pronouncedly religious
states the need for knowledge does not stir independently, and therefore not in
a critical manner either; the need for expression and reaction, which is one
with a need for symbolization, has decidedly the preponderance and will
therefore become determinative for the content of the religious idea.

When immediate experience cannot give sure direction as to the cause of the
lived experience, it can still less contain reliable information as to whether
the lived experience is due to a natural or to a supernatural cause. In exact
descriptions of religious states it also often becomes clearly evident that one
infers a supernatural cause only because one cannot find any natural cause and
does not think that one has oneself been co-operative. Thus Saint Teresa says:
``I felt my soul kindled with burning love of God; this love was evidently
supernatural, for \emph{I did not know} what thus kindled it in me, and
\emph{I had not myself contributed} anything to it.'' It is a trait that
constantly recurs among mystics and pietists, that the more they hold back (or
suppose themselves to hold back) their own thought and will, the more will that
which is lived through in their inner life be assumed to be of divine
origin.\textsuperscript{25} There lies here the same dualistic conception at
the foundation as with the concept of miracle generally. Psychological miracles
stand quite on a line with physical miracles, though the belief in the former
may hold out longer than the belief in the latter.

\medskip

\noindent\textbf{24.} Religious ideas can stand in a somewhat different
relation to the immediate lived experience. It may happen that the lived
experience unfolds as a series of states, and that only afterwards is an idea
sought which interprets it (that is to say: explains it or expresses it or
both). Such an idea can be formed anew, under the impression of extraordinary
lived experiences and in original, creative personalities, although even
founders of religions are involuntarily guided by traditions with which their
lived experiences fuse, yet in such a way that a wholly new direction in the
life of the spirit can be introduced.---Most often the interpretation is drawn
without more ado from the circle of transmitted ideas---ideas that were indeed
possessed and known beforehand, but not particularly attended to. There are
examples of persons becoming clear only afterwards that they had been
``converted'' or ``born again,'' by comparing their lived experiences with
``the statements of trustworthy believers.''\textsuperscript{26} The change
that is noticed is interpreted by the help of tradition, for which one now
first gets a practical use and which one transforms from a dead treasure into a
living one. But very often the relation is such that the transmitted ideas
themselves call forth the lived experiences. For most human beings religious
experience is an attempt to live through, in their inner life, the states that
transmitted ideas describe and demand. The ideas are thus fully finished
beforehand and are merely to be converted into lived experiences. Most human
beings make their religious experiences within religious communities with a
more or less developed doctrine of faith. Thereby the extent of the experiences
is dammed in beforehand. Lived experiences that cannot find their expression or
their explanation within the transmitted circle of religious ideas stand as
uncertain, or as dangerous, or even as pernicious. In theological expositions
this confessional character, which is so frequent in religious experiences,
appears clearly and often with full consciousness.

The older Pietism laid---in contrast to the orthodoxy it stood over
against---great weight on the inner experiences. But since the Protestant
Church maintains the Bible as the highest norm, the ecclesiastical pietists
expressly declared that inner experience was to be subordinated to the Bible's
teaching. In his remarkable work \emph{Theologica experimentalis} (1715)
\textsc{Gottfried Arnold} taught that religious experience stands under
Scripture and does not denote a special principle alongside it. Experience's
significance he saw in this, that it was a practising of, or a living into,
what Scripture taught. Experience is thus not ground, but consequence and
fruit, and it is to be ``judged according to God's word.'' With the clarity
that distinguishes Arnold, he then draws from this the conclusion that
experience must occupy a different place in religion from that in science: in
religion it comes after faith, in science it goes before faith. More recent
theologians (for example \textsc{Henry Ussing} in his work \emph{Den kristelige
Vished)} have on the other hand attempted to parallel religious and scientific
experience without seeing the consequences, disquieting for themselves, to
which this must lead. At any rate they are compelled to concede that the
individual cannot by the road of experience come to know everything that he is
to believe according to his Church's teaching. And what significance has ``the
common experience of the church community'' for the individual, and can it
rightly be called a common experience when not all make it or can make it?
Besides, the Church's teaching contains a great deal, for example the creation,
the day of judgment, and so on, which no one can ever have experienced.

The confessional limitation has made religious experience, and indeed all inner
experience, one-sided and imperfect. It has kept the eye away from sides of the
inner life that had a full claim to come into consideration. And what in a
philosophy-of-religion respect is the most important thing: it has made it
psychologically inexplicable how the religious ideas came into being in the
first place. Religious experience must surely be able to lead to the formation
of new ideas, or at least to the valuing of other ideas than the traditionally
sanctioned ones. There must in the domain of inner life be discoveries to make
just as well as in the domain of outer life. The conditions of inner life
change just as well as those of outer life; that follows already from the fact
that there is a constant reciprocal action between the inner life and the
outer. One cannot then dogmatically assume that all essential spiritual
experiences should already have been made, so that all following generations
have now only to practise the forms and thoughts that at a certain definite
point of time were laid down. A proof that the world of the spirit is closed
will never be able to be given.

\medskip

\noindent\textbf{25.} The life-value of an idea is established both by the fact
that experience leads one to form or to choose it as the expression of what has
been lived through, and by the fact that holding fast to it and sinking oneself
into it helps to keep one's courage up in life's struggle, to preserve fidelity
to the best we know, and to set us tasks that carry us forward in personal
development. Whether therefore an idea stands as the effect of experience or as
the cause of experience, it is tested through experience. But it must be well
held fast that a complete experiential proof of an idea's life-value is present
only where there has been opportunity to test whether the same experiences
could not be expressed in another way, or whether the same or equally valuable
experiences could not be made in the attempt to live in ideas other than this
definite one. There must not merely be a test made, but also a counter-test---
and for that life very seldom gives the individual opportunity. The individual
will as a rule be restricted to experiencing that certain ideas under certain
circumstances have helped him or not helped him. He will not be able to prove
that he could not have been helped by other ideas than those that helped him.
Immediate lived experience can at any rate tell him nothing about that. And
just as little can it tell him anything about whether that which did not help
him might not be able to help others---or to help himself under other
circumstances. Experience will here always preserve a personal-individual
stamp, and ``common experience'' is more or less an illusion, because the
different individuals in reality interpret and apply ``the common'' each in his
own way. When there is thus already misgiving about the transition from the
individual experience to the confessional, it must awaken still greater
misgiving when confessional experience makes a claim to be accessible to all
human beings, if only they ``will'' receive it. Here again there lies at the
foundation the presupposition that one knows all possibilities, that is to say
here: all personal circumstances and conditions.

\medskip

\noindent\textbf{26.} Hitherto we have investigated religious experience
\emph{as} experience, that is, according to the conditions and peculiarities it
has in common with other experience. Now we shall enter upon what is
distinctive of \emph{religious experience} and makes it different from other
psychological experience.

Religious experience is essentially religious feeling. It is the mind's own
inner state during the course of inner and outer events that is its immediate
object. But there are other experiences of feeling than the religious. In
describing the nature of religious feeling it will be a matter of stating its
likeness to and its difference from the most nearly related feelings.

All feeling, that is to say, all pleasure and displeasure, of whatever kind it
may be, expresses the \emph{value} which something that happens in the inner
or the outer world has for us. And conversely: that has value which is the
object of immediate satisfaction, or which can also become a means to such.
(Cf.\ §3.) The concepts of end and means presuppose the concept of value, and
value again presupposes a subject that can feel pleasure and displeasure. When
we ascribe value to things and states, this property, like all properties,
denotes a definite relation to a subject, here to one that can feel pleasure
and displeasure. This property is derived in relation to other properties that
things and states have; for it must be these properties that make the things
and states valuable. Value is thus a property in the second power.

The different kinds of value correspond to the different kinds of feeling. One
group of values hangs together with self-assertion, from its most elementary to
its most ideal forms. Another group of values hangs together with devotion to
beings, relations, and tasks that point out beyond the conditions of isolated
self-assertion; to this belong the ethical, aesthetic, and intellectual values.
The possibility of a third group of values, the religious values, rests on the
fact that the two first groups of values, those of self-assertion and of
devotion, are to be acquired and held fast within existence, such as it once
for all is. Existence thereby comes to stand for the human being as a
battlefield on which the fate of values is decided. It is a great drama that is
performed, in which the human being is both fellow-player and spectator. Were
he merely fellow-player, his whole energy and all his interest would be taken
up by the place he himself is to fill in the drama, and he would not have time,
strength, and interest to be more deeply affected by the course of the drama as
a whole. Were he merely spectator, his mood at the drama's course would be
purely intellectual or aesthetic. But when he is both, he will himself have
values that are at stake in the struggle, and besides his own participation in
this, the picture he forms of the drama's course will also affect and attune
him. He will feel himself so drawn in, with his whole inmost being and with
respect to the highest values he knows, to the whole great order and course of
things, that there will arise in him a living feeling of displeasure or
pleasure, according to the fate the values suffer. In its most immediate form
this feeling finds vent in outbursts of hope and fear, admiration and disgust,
joy and sorrow. Such outbursts are in all domains the simplest value-judgments
there are. \emph{This feeling, determined by the fate of values during the
struggle of existence, is the religious feeling.} It is thus determined by the
relation between value and reality. This relation---as it stands for the human
being---determines the value the human being ascribes to existence. The
religious judgments are thus secondary value-judgments, are derived in relation
to the primary value-judgments in which the two first groups of values are
expressed.

The religious values can---in spite of their derived character---be felt just
as immediately and just as energetically as the primary ones. They will
themselves be able to be the human being's central values, since the relation
between value and reality is the strongest relation of tension in which the
human being is placed, and which can therefore lay claim to all his energy; for
what is at stake here is at last life or death for everything that has shown
itself to have value for the human being. In the religious problem the concept
of value appears in the second power, derived to be sure, but in the same sense
in which a concentration is derived in relation to the forces that are
concentrated towards one and the same point. In religion's classical times, the
time of beginnings and the time of organization, such a concentration of all
interests and forces towards one point is the most characteristic trait.

And just as religious feeling, in spite of its derived character in relation to
other feelings, can be just as immediate and energetic as they, so it can also
appear independently over against them, react back upon them. It by no means
stands merely passively related to them. There is here a constant reciprocal
relation, and the degree of independence may be highly different on both sides.
With greater independence the two first groups will be able to appear only at
more developed stages of culture, where the spiritual division of labour has
come into force. As we have earlier seen, the religious problem presents itself
precisely only where the special values acquire a greater degree of
independence over against that concentration of all values which is religion's
distinctive character, and which psychologically finds its explanation in the
fact that it is the fundamental relation between value and reality that in
religion is the really determining thing.

The psychological description and analysis of a feeling must naturally not be
confused with the feeling's own reality. In order to describe and analyse we
need a scaffolding of concepts and points of view which the real state of
feeling itself has no need whatever to know. It appears, or can appear, with a
closed character of wholeness; it grows forth in the manner in which an
instinct forms itself: under the influence of impressions and lived experiences
which the human being cannot always account for, and which only the comparative
investigation points out.

The feeling whose possibility and general character has now been described
constitutes the most essential distinctive mark of all religions and religious
standpoints. In relation to it all ideas are subordinate and conditioned.

One might object to this that in the great historical forms of religion it is
after all an essential trait that personal beings, on whom the human being
feels himself dependent, are represented as ruling in existence, and that the
concept of religion therefore presupposes myth, dogma, and cult.---How one will
define the concept of religion is at last a free matter. As with all
fundamental definitions, it is grounds of expediency that here decide. Now it
is not expedient to define a concept in such a way that one thereby cuts off
important questions and problems. The religious problem becomes undeniably very
much simplified if one will speak of religion only where myths, dogmas, and
cult can be pointed out. But if myth, dogma, and cult always have religious
significance only by being borne by a feeling such as that described above,
then there shows itself a possibility that this bearing element in religion can
also be and work without procuring itself expression in myth, dogma, or cult.
It is then natural to lay the chief weight on it, and the wider extent of the
concept of religion therefore becomes the most natural. Great importance this
question does not moreover have, since the dispute about it easily becomes a
dispute about words. If one would rather avoid the word religion than use it in
a wider sense, one might, as I have proposed in my \emph{Psychology}, call the
feeling that is determined by the relation between value and reality
\emph{cosmic life-feeling.} For just as the organic life-feeling immediately
corresponds to and is determined by the course of life in our organism, and
acquires a different character according as the organic life in us is checked or
furthered, so in the feeling described there will be given a symptom of how it
goes with life (so far as we know it and tie our weal and woe to it) in
existence as a whole (so far as we can form a picture of it). If one will not
take the word religion in the wider sense, one may call religion (in the
narrower sense) a \emph{distinctive} kind of cosmic life-feeling.---To use the
word religion in the wider sense would only be unfortunate if this usage were
used for unworthy accommodation to existing forms of religion. But the criticism
of the existing forms of religion becomes precisely sharper if it can be shown
that these in many ways come into conflict with religion's own essence and
therefore point out beyond themselves. Besides, those who wish to compromise
will always be able to find their ways out, and there is no reason to give up a
usage that is expedient in a scientific respect because it can be misused for
unscientific ends.

In his brilliant book \emph{De l'irréligion de l'avenir} (1887) \textsc{Guyau}
holds fast to the narrower usage as regards the word religion, in that he will
not speak of religion where myth, dogma, and cult are lacking. But he then
distinguishes between \emph{irréligion} (non-religion) and \emph{antireligion}
(religion's opposite), and his endeavour aims at establishing that ``the
non-religion of the future will be able to preserve the purest thing the
religious feeling contained.'' The question is, it seems to me, whether ``the
purest thing in the religious feeling'' is not precisely the really religious
thing itself. The American philosopher of religion \textsc{J. Royce}, who
stands at a different standpoint from Guyau, says in an interesting essay he
has written on the young French philosopher who departed so early: ``If Guyau's
opinions were my own, I should without hesitation call them religious, because
in them---like he himself---I should see the rationally formed consummation of
what humanity's religious instinct has sought after.''\textsuperscript{27}---We
find ourselves here in the domain of fine nuances. I for my part would not have
so great a misgiving as Guyau about the wider application of the concept of
religion, but would not be quite so sure of my case as Royce.

\medskip

\noindent\textbf{27.} The description of religious feeling attempted in what
precedes is confirmed partly by the fact that it lets the difference between
religious feeling and other feelings come out clearly, partly by the fact that
starting from it one can with ease explain the different forms that religious
feeling can assume.

The more the human being is taken up with the very work of asserting himself,
or the more he is absorbed in intellectual, aesthetic, and ethical interests,
the more does the really religious recede, if it does not wholly disappear.
Existence is then a resistance to be overcome, or an object that is understood
or contemplated, or a foundation for tasks of the will.---Especially
instructive here is the difference between the intellectual and the religious
relation to existence. Existence presents to knowledge an unsurveyable richness
and multiplicity; it contains far more than human beings will ever be able to
bring under the point of view of what they call values. It realizes forms and
stages of existence that go far beyond, and often stand in direct opposition to,
what that which from a human standpoint stands as the valuable would demand.
Intellectually regarded, this overflowing fullness is a good, because the
picture of the world at which our knowledge works thereby becomes the more
comprehensive, and because the endeavour is the more strongly awakened to
understand the individual form or the individual stage in its distinctiveness
and in its definite connection with the whole. The religious interest, on the
other hand, aims at seeing this whole richness of forms and stages as the
expression of the realization of great values in existence. For the religious
interest the great multiplicity contains a danger and a temptation, since the
reality of values cannot be pointed out so easily as if the content of existence
were less luxuriant. The unsurveyable world seems to be too great for the human
values. There may then here come a collision between the intellectual interest,
which is always on the track of new differences and new connections and
therefore rejoices in the fullness, and the religious interest, which in order
to be able to hold fast to the point of view of value has a natural tendency to
limit the picture of the world (which became a chief point in the epistemological
philosophy of religion). This opposition comes out clearly when one compares two
natures like \textsc{Pascal} and \textsc{Spinoza}. What terrifies the one,
nature's infinity, delights and inspires the other. And yet Spinoza was in his
way a religious nature. Intellectual interest can, namely, involuntarily bring
with it an evaluation of existence. When existence in all its fullness
nevertheless shows itself intelligible, it is indeed from an intellectual point
of view a home for values, comes in a manner to meet the human spirit. It
procures for us the deep joy of knowledge, which becomes of especially penetrating
effect when we understand ourselves as distinctive links in the great connection.

Then arises the intellectual love which was the highest thing for Spinoza, the
quiet, inspired contemplation in which everything was at last harmonized for him
into a total picture.---It will be easy in analogy with this to develop the
relation between religious feeling and the feelings that stand related to it in a
similar manner to the intellectual feeling. The religious relation is everywhere
characterized by a striving to hold fast the persistence of the valuable in
existence: whatever the valuable may be that the human being sets highest, and
whatever the conception may be that the human being has of existence. Life,
truth, beauty, and goodness must be presupposed to have been experienced in order
that a religious relation should be possible. The religious relation comes into
force when these values are compared with the rest of reality. From this one sees
at once the opposition and the connection between religious feeling and other
feelings. It must here be noted that what I here distinguish between as two
different acts, the experience of life, truth, beauty, and goodness, and the
experience of these values' relation to factual reality, need not appear at
different points of time, but may coincide in one act, so that the human being
himself does not become conscious of the difference. It is in the psychological
description that we make this difference; it need not be made in life itself. The
religious strife gets its sharp sting precisely from the fact that at the
classically religious standpoint there exists no valuable life, no truth, no
beauty, and no goodness outside religion; this has in its concentrated form taken
up all valuable elements, so that these do not come to appear in their
independent distinctiveness. This happens, as often remarked in what precedes,
only when the religious problem is posed as a consequence of a beginning division
of labour in the spiritual domain.

The differences between the religious standpoints may rest partly on the
different values that are presupposed to have been experienced, partly on the
different conception of existence that lies at the foundation, partly on the
different degree of energy with which value and reality are compared, so that
the relation between them can become the object of penetrating experience. These
three sources of differences all spring from the given description of religion as
a feeling determined by the relation between value and reality. Every exhaustive
characterization of a given religion or religious standpoint will lay these three
points of view at its foundation. Both in what precedes and in what follows there
are examples of these points of view's significance.

\medskip

\noindent\textbf{28.} If the conception of the world immediately showed us the
highest valuable thing realized, or at least clearly showed us existence as the
secure home of the valuable, no special religious feeling would arise. There
would be an immediate harmony between explanation and evaluation of existence:
what could be pointed out as an operative cause in the world would thereby at the
same time have been pointed out as a means to the realization of something
valuable. There would, to use an expression from political economy, be no
marginal utility. As it is, on the other hand---since the connection of values
with the rest of reality is so great a problem---a distinctive exertion is
required in order to hold fast the conviction of the persistence of the valuable.

Nor would any special religious feeling arise if the human being's capacity to
work for the valuable (in himself and in existence) were unlimited. Religious
feeling presupposes a striving to assert the values in us and outside us---a
striving that meets resistance. The human being feels himself dependent and
divided precisely in the inmost and most intimate domain. The feeling of
dependence first properly arises where there is a striving to press forward. The
chained animal will not feel itself dependent except when it strives to reach
further out than the chain reaches. When the feeling of dependence is something
other and more than a lassitude that may derive from dullness, it presupposes that
one has gone to the will's limit. The lion in the cage, pacing up and down along
the bars, has gone to the limit and feels its dependence; if it kept still at the
back of the cage, it would feel no dependence, but would find existence far
simpler and more harmonious than it now does. But in order that a religious
relation should arise, it is not enough to go to the will's limit. At the very
limit and in spite of the limit the wish for the persistence of the valuable must
be held fast. There must form ideas of a world of values whose reality is not
given up because human thought and human capacity have found their limits. There
may arise an intimation that the principle of this world of values is at last one
with the principle of existence's causal connection---that what makes it possible
for us to find values in existence is the same as what makes it possible for us to
find it intelligible.

In \textsc{Schleiermacher}'s famous description of religion as a feeling of
dependence it is not sufficiently emphasized that this dependence is conditioned
by an activity, and that it arises at this activity's limit. Nor is it
sufficiently emphasized that the dependence is felt with respect to the work for
those values that for the human being stand as the highest and as those which in
spite of everything persist and can therefore be asserted by human striving.

The role that the consciousness of sin plays in several religions is an example of
the character of the religious feeling of dependence. In the consciousness of sin
the human being feels the disproportion between the ideal of will that his
evaluation has led him to acknowledge, and the reality of his will. It is now
sluggishness and lassitude, now dividedness or dispersion, now one-sided or
impatient concentration, that makes his inner reality so different from the ideal.
In the consciousness of redemption or reconciliation there then dawns, often
suddenly, the conviction that the valuable in spite of everything persists,
persists in his own inner being and will come to be victorious there. This inner
psychological drama stands, in the highest popular religions, in Buddhism and
Christianity, as the real world-drama, to whose unfolding in souls the great
cosmic processes at last all minister.

\subsection{Religious Faith}
\label{ssec:faith}

\noindent\textbf{29.} Like all other experiences, religious experience too is
shaped in more or less definite ideas, whose content depends on the stage of
development and knowledge at which the human being stands. The harmonious or
disharmonious relation between value and reality gets its most immediate
expression in exclamations. When the human being commands a larger circle of
ideas, he will involuntarily compare the feeling that the relation between
value and other reality awakens in him with other feelings and experiences, and
thereby, through likeness (especially analogy) with, or opposition to, these
other experiences, closer expressions and designations for religious experience
will be able to be won. Thus there are used in religious language expressions
like life and death, health and sickness, light and darkness, truth and
falsehood, beauty and ugliness, justice and injustice, to denote the
oppositions that are lived through in the religious domain. The expressions are
drawn from the domain of the two first groups of values, from the domain of
life-preservation, truth, beauty, and goodness.

The religious life of feeling oscillates, like all life of feeling, between
pleasure and displeasure, between vehemence and lassitude, between inwardness
and outwardness. But since in every personal being there is a more or less
pronounced striving after unity and connection, there must in spite of these
oscillations be certain chief ideas to which the mind constantly returns, and in
which it finds expression for what is essential in its experiences. In the dark
and weak moments it seeks to hold fast what it has experienced and thought in
the bright and strong moments. It strives to preserve continuity in spite of the
oscillations and through the oscillations. In religious faith there is expressed
not merely the individual expression of religious feeling, but above all an
endeavour to hold fast that relation between value and reality which the most
strongly intervening experiences have taught the human being to acknowledge.

In the concept of faith there lies the conviction of a continuity, a persistence
out beyond the horizon experience has shown us, and through the interruptions
and leaps that experience presents. We believe in a human being when we hold
fast the conviction that through all changes he will be like himself, although
his life does not lie before us as completed or as fully known. We believe in
ourselves when we cherish confidence that we shall not fail the best in us, but
shall hold fast our decisive resolutions through all hindrances. And thus
religious faith too is a conviction of a firmness, a trustworthiness, an
unbroken connection in the fundamental relation between value and other reality,
however much the states of reality, and with them the empirical appearance of
values, may change. Faith is akin to faithfulness, and presupposes faithfulness
(in the sense of firmness) in faith's object as well. Faith is a subjective
continuity in disposition and will, which aims at holding fast an objective
continuity in existence. Faith's object is the persistence of value---but the
very existence of faith is a testimony to the persistence of value within the
individual personality.

Faith arises in its simplest and healthiest forms involuntarily. But that does
not exclude its being a matter of will, when one takes the word will (as is done
in my \emph{Psychology}) in the wider sense. An unfolding of activity is needed
when an expression or an explanation is to be held fast in the face of the
oscillations of states, in the face of what is incomplete and often
contradictory in experiences. But in its simplest forms this activity can appear
as instinctive confidence or secure expectation. Only when the opposition
between value and other reality assumes sharper and more glaring forms does
faith appear as desire, wish, purpose, or resolution, in that the concentration
of all the mind's forces upon the leading and saving thought is undertaken with
greater consciousness of what is at stake. In faith, moreover, what perhaps can
only be realized through long ages is lived through as present, just as in the
moment of resolution we already regard ourselves as acting, although there may
be a long road still to the carrying out of the action. Through this
anticipation the difference between means and end is abolished; the time of
preparation and expectation is no longer merely a means for a later time, but
the fruit of the work is lived through in the midst of the time of work, because
the work itself is the highest reward. In religion's own language this is
expressed by saying that eternal life is lived in the midst of time, does not
first come after long ages.

There is thus a close kinship between faith and will. In describing any
expression of will that lies beyond the purely instinctive, the chief point will
be the holding fast of the idea of a goal that is to be reached, whether this
goal is again a means for a more distant goal or not. It is by this
moment---the holding fast of an idea across greater or lesser distance and
resistance---that the concepts of faith and will pass over into one another, so
that they can even be said to be names for one and the same phenomenon of the
soul. If any difference is to be made, then faith denotes rather a rest, a
devotion---the rest and devotion that lie in concentration upon a single idea.
One may speak of a ``will to believe,'' in so far as a definite striving to
attain such a concentration may express itself. Conversely one may also speak of
a faith in the will, since a confident devotion to the endeavour of will and the
conviction of its energy may be a necessary condition for the will's work to be
done.

% NUMBERING. The e-book prints "30." at the head of the next paragraph AND at
% the head of "En karakteristisk og meget hyppig Type" further below. Only one
% can be right, and the cross-references settle which: §32's paragraph refers
% back to "den anden (31) og den tredie (32)", and §34's to "den første (30),
% dels om den tredie (32)", so the numbered types are the ones beginning "En
% karakteristisk og meget hyppig Type" (30), "I Modsætning til denne Type" (31)
% and "En tredie Type" (32). The paragraph below is therefore unnumbered and
% belongs to §29; its "30." is a stray numeral of the same defect class as the
% headless sentence at §14. See RESUME-NOTES.

Significant differences in the nature of religious faith are conditioned---like
the differences in the nature of religious feeling---by differences in values
and motives of evaluation, by differences in the knowledge of reality, and by
differences in the energy with which value and reality are compared and held
together. These differences are not all of equally great importance. It will be
easy to see that the last-named difference will be the most important.

The motive of evaluation, which determines what comes to stand as valuable for
the human being, can vary from individual to individual, from people to people,
from time to time, both in kind and in extent. It may be hope or fear; it may be
egoistic, or at any rate individualistic, and sympathetic; it may be simple or
composite, and if it is composite it may have a different timbre according to
the different relation between the elements that co-operate in it. In respect of
extent the difference between egoism, individualism, and sympathy will also be
decisive: the circle of values that is to be preserved may be attached either to
the individual himself, with more or less conscious isolation, or to groups of
human beings and endeavours within which the individual's own persistence and
striving stands only as a single link.

The conception of the world gets its significance for religious faith by
indicating the roads along which values can be unfolded and preserved. It
presents the whole stage on which the great drama takes place in which the fate
of the valuable is decided, and the stage determines in many ways the kind and
course of the play.

But the properly religious moment arises---if the preceding description is
correct---only through the comparing of evaluation and conception of the world.
Therefore the most important group of differences that acquire influence on
religious faith must be that which arises through the different inwardness and
energy with which value and reality are placed together in consciousness. One
might call them differences in religious synthesis. Naturally all three moments
always work together; but this last is the most central. In what follows I shall
bring forward some of the most typical differences that are predominantly
determined by the varying of the last-named relation.

\medskip

\noindent\textbf{30.} A characteristic and very frequent type of religious faith
is that which is determined by the need for rest. And what makes one weary and
tired during the course of life is above all unrest and dividedness. We are
affected from so many sides and have trouble collecting our thoughts; we are
drawn in so many directions and find it hard to hold the will fast upon a single
goal; so many different and shifting feelings are set in motion that the mind's
inner harmony threatens to dissolve. There is inner distress on account of the
feeling of the disproportion to our ideal, and there is outer distress in pain,
infirmity, and dependence on life's elementary necessities.

In the \emph{Upanishad} it is said: ``The Self (Atman), the sinless, that which
is redeemed from old age, death, suffering, hunger, and thirst---that whose
wishes are right and whose resolve is right---that one should explore, that one
should seek to know. Whoever has found and known this Self attains all worlds
and all wishes.'' And in another place: ``Save me, for I feel myself in this
course of the world like the frog in a dry well!''---Jesus of Nazareth says:
``Come unto me, all ye that labour and are heavy laden, and I will give you
rest\ldots. Learn of me, and ye shall find rest unto your souls.''---``Restless
is our heart,'' says Augustine to his God, ``until it finds rest in thee!''---

In passionate form this need for rest appears in natures like Saint Teresa,
Pascal, and Søren Kierkegaard.

There is indeed in Augustine too a mighty pathos; but one always notices in him
the Platonist and the prince of the Church alongside the intensely strained
seeker; it is precisely the union of all these elements that makes him so unique
a figure in the history of the religious life. Saint Teresa feels the need for
union with God so violently that only death can satisfy it: ``I did not know
where I should seek this life, except in death. Nothing could answer to my
wishes; my heart was every moment on the point of bursting, and it was to me as
though the soul were being torn out of me.'' ``I thirst to live, and I live
dying! I die because I cannot die! The fish that is drawn up out of the water
nevertheless sees its torment end; but what death can be compared with the life
in which I languish?'' For Kierkegaard too it was \emph{life's strife} that it
was a matter of being released from; the lines of the hymn he wished set on his
grave express it:

\begin{quote}
``It is a little while, then I have won, then is the whole strife at once
vanished.''
\end{quote}

Here in life the believer is in an alien element; there is a disproportion
between the inner and the outer, between life and its conditions. In him too the
image of \emph{the fish on dry land} comes forward in one place; it is
characteristic of this type that this same image has thrust itself forward both
in the ancient Indian of the \emph{Upanishads}, in the Spanish nun of the
sixteenth century, and in the Nordic thinker of the nineteenth. This trait is
instructive with respect to the psychology of religion. The human being's goal
stands as infinite, but the human being must lead his life in the world of
finitude, and from this it follows that his existence acquires a convulsive
character. In Kierkegaard, and already in Pascal, the opposition is emphasized in
a still sharper manner than in Teresa. In Teresa it awakens longing and inward
yearning; her ``will'' is laid claim to by the highest object, and only her
memory and her imagination work to interpret her lived experiences. But both in
Pascal and in Kierkegaard it is precisely the will that is again and again
summoned up, in order that they may hold themselves up under the cutting
disproportion between the true life and the actual conditions of life, and above
all in order that thought may hold fast to faith's object and not succumb to
doubt.---As an example of this type John Bunyan (the author of \emph{The
Pilgrim's Progress}), who has himself described the story of his conversion, must
further be named. Deep melancholy, ungovernable motor impulses, restless
vehemence here formed the psychological foundation for a series of spiritual
struggles, which were brought to an end only after a stage of despair had led
over into a state of complete passivity, in which strength could be gathered for
a new and higher life. Even after he had won a firm position and had begun his
vigorous activity as a curer of souls, he had again and again to summon up his
self-control in order to keep the old chaos down.\textsuperscript{28}

\medskip

\noindent\textbf{31.} In opposition to this type, which is characterized by the
fact that it is disharmony that lies at the foundation and awakens the need,
stands a type whose distinctive character is an inner need for self-unfolding
and devotion, which leads out beyond all oppositions and limitations, almost
without touching upon them. Religious feeling here becomes properly no
distinctive kind of feeling, but a heightening and widening of the feeling with
which the valuable is already beforehand embraced. Religious feeling here arises
through a further unfolding of the drive for self-preservation. There is here no
opposition between self-assertion and devotion, for devotion arises through the
human being's commanding far more strength and feeling than his own purely
individual interests can lay claim to. In overflowing joy and enthusiasm he then
embraces the life of others, at last the whole of existence. Religion springs
here from the strength to live and from the joy in life. It seeks rest here too,
but not rest from pain and dividedness---rather rest for a yearning that no
limited object can satisfy. The happy man will, said Aristotle, see life about
him, a life he can share and bear.

In an interesting manner this type comes forward in the modern age in
\textsc{Campanella}, \textsc{Joseph Butler}, and \textsc{Rousseau}.
Characteristic is especially Butler's statement (in a sermon ``on the love of
God''): ``As we cannot remove from this earth or change our whole business upon
it, so neither can we alter our real nature. Therefore no other spiritual
exercise can be required of us than the exercise of the faculties of which we are
conscious. \emph{Religion does not demand new feelings, but only the direction of
the feelings we already have,} though they are limited to objects that do not
satisfy their nature.''---This type lies at the foundation in Goethe (in Faust's
confession of faith and in the Marienbad Elegy) and in Schleiermacher (in his
\emph{Speeches on Religion}). It comes forward in a manner recalling the Gospel
of John in the English theologian \textsc{Fr.\ D.\ Maurice}, for whom hell
consisted in separation from God, and who---in contrast to those who ``construct
all theology out of sin''---says in one of his letters: ``How I long to say to
myself and to all, that the hell we are to flee is ignorance of the perfect Good
and separation from it, and that the heaven we are to seek is the knowledge of it
and the participation in it.''\textsuperscript{29}

\medskip

\noindent\textbf{32.} A third type is characterized by the role the
intellectual and the aesthetic element plays. For contemplative natures what
matters is to win \emph{a total conception} within which the relation between
value and reality finds its illumination. Now it is thought's need for
understanding, now the imagination's need for vivid images, that predominates.
Plato's doctrine of ideas satisfied both requirements for him and gave his mind
rest in the contemplation of the eternal ideas as the true reality, in relation
to which the sensible world in its constant change stood at last as an illusion.
A tendency in this direction occurs already in the \emph{Upanishads.} From
Platonism this type has passed over into Christian theology, where it is traced
in Augustine, in the mystics of the Middle Ages, and (in slighter forms, to be
sure) in more recent speculative theologians, as with us Martensen. A pronounced
representative of this type is \textsc{Spinoza}. Although his thought, so long as
it moves on the ground of experience, has a decidedly realistic character, the
highest consummation lay for him nevertheless in the contemplation of things
``under the form of eternity,'' a contemplation that could indeed only be reached
through serious work of thought, but which is nevertheless rather of an artistic
than of a scientific kind. The attainment of such a contemplation stood for
Spinoza as the highest good, as the only thing that could bestow lasting and deep
satisfaction. In this contemplation the most inward life was combined with the
most serious reflection upon life, and the joy of knowledge that was awakened by
it cast its radiance over the whole picture of the world and obscured the
disharmonious traits it presented when one remained at the finite and limited
standpoint. Every feeling and every experience can give its contribution to this
highest spiritual state, and access to this good stands open to everyone who
strives. The strife that finite goods can awaken between human beings falls away
here, and the joy in it grows the more there are who gain access to it.

Many of the representatives of this type have supposed that in the highest idea
or the intellectual contemplation, which for them was the highest thing, they had
a result of the highest science. This rested on an illusion, which was founded on
an imperfect investigation of knowledge's conditions and limits. If it is
philosophy that these spirits end in, then it is not philosophy as science, but
philosophy as art.\textsuperscript{30} But this type's great merit is then
precisely to have asserted thought's significance for life, even there where
thought moves at its own limit. The art of thought is of great importance for the
art of life, and it would lead to spiritual barbarism if this type came to lack
representatives.

With the need for contemplation there naturally combines the need for intuition.
In Plato and Spinoza the need for visual images seems to have ruled in the
imagination; they have been visual. The same holds of the author of the Gospel of
John. Although not only the opposition between light and darkness, but also the
opposition between life and death, frequently comes forward in his circles of
images, it is nevertheless the first opposition that rules, and life and death
are themselves often represented as day and night.---In contrast to this, Paul in
his highest states has been partly vital (in that the life of ideas was often
determined by the life-feeling corresponding to inner organic states), partly
motor (in that ideas of movement are frequent in him). The highest could not for
him be shaped according to the experiences of the eye or the ear, and the most
inward need brings forth unutterable sighs. He has a great capacity to ``speak in
tongues,'' a decidedly motor trait. Life and death appear for him as opposed inner
states that are sensed in the limbs. Chasing and struggle, longing and hope, flesh
and spirit lie at the foundation of characteristic turns of phrase in the Pauline
epistles. Visions and auditory hallucinations are subordinate in relation to
sensations of life and movement.---The hallucinations of presence which by some
mystics (Saint Teresa and Cécile Vé, for
example)\textsuperscript{31} are definitely stated to be different from
``imaginative'' hallucinations (hallucinations of sight and hearing) probably
belong under sensations of life and movement.

In contrast to contemplative or intuitive natures, in whom a concluding total
conception easily forms, stand the reflective natures, in whom the activity of
thought itself, not its conclusion in a total picture, is the chief matter, and
who perhaps even lack the capacity and the need to form and hold fast definite
images. Reflection in such natures turns quickly towards the points where the
total picture's insufficiency shows itself (where ``the similitude'' limps), and
it perhaps even feels a shudder at the contradiction that arises from the attempt
to present the infinite in finite form. Schleiermacher\textsuperscript{32} is a
characteristic example of a reflective type, but at the same time an example of
the fact that the intellectual dispositions need not determine the whole type.
There is a constant reciprocal relation between the intellectual, the emotional,
and the volitional dispositions. In mystics it will especially be the inward
feeling that brings with it the conviction of the insufficiency of concepts and
images to form an expression. Hereby there arises a certain spiritual kinship
between mysticism and criticism, even where they are not united in the same
person. Many of the arguments Carneades used against the dogmatism of the Stoics
are found again in the Neoplatonist Plotinus, and in more recent times the
combination of mysticism and criticism is not rare either. Living feelings and
clear thinking can far more often join hands than one as a rule supposes.

% NUMBERING. The e-book prints "32." again at the head of the next paragraph.
% It is a stray numeral, of the same class as the one at §29: this paragraph
% continues §32. The cross-references confirm it — §33 below refers back to
% "den frejdige Voven (32)", i.e. to the very type described in this paragraph,
% which therefore falls inside §32; and §34 refers to "den første (30) ... den
% tredie (32)", fixing 30/31/32 on the three type-paragraphs above. Numbering
% the paragraph separately would leave two sections competing for §33 and push
% §34 off the reading that two later chapters (III.D, III.E) cite by number.
% See RESUME-NOTES.

While the last-described form of religious faith is characterized by intellectual
dispositions, we come more over into the domain of the will when faith appears as
\emph{confident venturing.} Luther and Zwingli have, often in sharp opposition to
the Schoolmen's doctrine of faith and knowledge as differing only in degree,
strongly emphasized bold confidence (\textit{fiducia}, \textit{Vertrauen}) as
faith's essence. Especially in ``the Great Catechism'' Luther has with energy
uttered this thought, so that it even serves him as the foundation for a
definition of God: ``A God is that from which one is to expect all good, and to
which one is to take one's refuge in all distress, so that to have a God is
nothing else than to believe and trust him from the heart, as I have often said,
that it is the heart's faith and trust alone that produces both a God and an idol.
If the faith and the trust are right, then your God too is the right one, and
conversely, where the trust is false and wrong, there neither is the right
God.''\textsuperscript{33} In this bold confidence, which indeed presupposes a
feeling of disharmony and distress (whereby we are reminded of the first type, see
§30), but which aims at this distress's being able to be overcome, indeed at its
being overcome, lies Luther's chief thought, the fundamental thought of the new
religion he founded without knowing it. If the movement---as in the first
type---is to be made again and again from disharmony to harmony, there is no time
and strength left over to take part in the interests and works of human life. The
bold faith that one is---as regards the essential---well provided for makes it
possible, on the other hand, to place oneself in a more positive relation to
actual human life and to take part with a full heart in the development of
culture. In his little work on ``The Freedom of a Christian Man'' Luther remarks
that the Christian who has made himself the servant of all thereby attains
precisely the right to rule over all things. Freedom and the great strength spring
from bold confidence in life's ground. This confidence is won (or is presupposed
to have been won) through struggle; thereby this type recalls the first and stands
in opposition to the second (§31) and the third (§32), which are predominantly
characterized by an immediate unfolding of the soul's forces.

\medskip

\noindent\textbf{33.} Yet another type (of the many that are possible besides
those now mentioned) must be brought forward. It is characterized by the fact that
it gives up its own standard, but otherwise stands nearest to bold venturing
(§32). It is determined by \emph{adherence to a model, an authority;} faith is
here an echo, a reflex, but an echo that becomes possible through inward devotion
to the model. Faith is here therefore grounded not on direct and independent
experience, but on trust in the experience of others. One's own experience is not
wholly excluded: it must indeed be experienced that one can lead one's life by
building on trust in one's model, find light by being led by the model's
light.---Within the relation to the model, which is also a relation of faith,
highly different types---thus all those described in what precedes---can come
forward again.

There is hardly any human being whose view of life---whether it has the character
of religious faith or not---is built right through on his own first-hand
experiences. Models and traditions determine us all, and they determine
especially also the experiences we come to make and the manner in which we take
them up into our life. But where a decisive weight is laid on a certain sum and a
certain form of religious ideas as the only valid content of faith, there the
dependence on model and authority becomes unconditional and comes at last to
stand as the only virtue.

When the Church had by the road of theological speculation worked out a series of
doctrines which had to lie far above most human beings' powers of comprehension,
it was nevertheless so humane as to declare that in the case of the most
speculative doctrines (for example the doctrine of the Trinity) it was not
necessary to believe on one's own account, but that it was here enough ``to
believe something because one believes that the Church believes it.'' And in this
willingness to hold to the Church's faith one found something meritorious, because
it sprang from ``the love that believes all things.'' Innocent III and later popes
have acknowledged such an ``undeveloped faith'' \emph{(fides implicita)} in order
to secure the unlearned against becoming heretics through innocent errors in
theological specialities. It is this kind of faith that has received the name of
charcoal-burner's faith, after a tale (reported---somewhat differently---in Luther
and in Erasmus of Rotterdam) of a charcoal-burner who, on being asked what he
believed, had only this answer: ``In what the holy Church believes!''---without
his being able to say what it then was that the holy Church
believed.\textsuperscript{34}

In his Protestant zeal Luther declared that he who had no other faith than this
charcoal-burner's went to hell. But the Protestant Church cannot---as little as
any church---dispense with charcoal-burner's faith. When it distinguishes between
the individual believer's experiences and ``the church community's common
experience,'' and when something is to be believed that by the nature of the case
cannot be an object of personal experience, then it is impossible to avoid the
acknowledgement of indirect faith (\textit{fides implicita}). It is Protestantism's
honour that it has so energetically asserted the individual's duty to make his own
experiences and draw his own conclusions; but it has dragged so much along with it
from the old Church---for example the doctrine of the Trinity, precisely the
doctrine that gave the medieval Church occasion to set up the concept of
\textit{fides implicita}---that it is an illusion when one believes one can
dispense with this concept.---Besides, this concept has its great humane and
pedagogical significance. All human beings begin their development with childlike
adherence to authorities and models. Their first expressions of freedom consist in
their experiences leading them to choose other models than the original ones, when
they are not confirmed in adherence to these by the course of life. Development
can lead to an increasing predominance of the results of independent experience,
and in quite a few---the great ones in the world of the spirit---creative
self-activity goes so far in depth and extent that they themselves come to stand
as models for the race.

The spirit of Protestantism demands that access to free investigation of
experience, its foundation and its results, stand constantly open. Traditions and
models must be investigated with historical criticism. Psychology must test
whether the ``experiences'' are natural and immediate; logic must investigate the
inner coherence of the assumptions; ethics must discuss the genuineness of the
values. If this testing is ever cut off, one ends in barbarism or in convulsive
attempts to hold fast the absurd. The religious consciousness has a tendency to
drag traditions along with it which have neither religious, intellectual, nor
ethical significance, dead values that can no longer really be lived through by
anyone, but which one dares not throw away for fear that they should draw more
with them in their fall. As \textsc{Rudolf Eucken} has said: ``one holds fast the
impossible in order not to lose the necessary.''

Every personal development begins with childlike adherence to transmitted ideas;
but there may be many degrees and nuances in the manner in which one later draws
nourishment from the tradition or step by step removes oneself from it. Some
natures never notice any difference between their own lived experiences and the
transmitted symbols. Others notice a difference but shudder back from it, fearing
the apparently self-abandoned ``subjectivity,'' and throw themselves, with greater
or lesser personal truthfulness and intellectual honesty, back into the common
forms. There are peculiarly ecclesiastical natures who live and breathe immediately
in what is transmitted, and who are always able to find both the end and the
beginning of their experiences in the transmitted forms. An immediate unity of word
and experience, of name and thing, prevails in them. They have no need for
individual symbol-formation, which often hinders spiritual fellowship, but will on
the contrary be inclined to make the possibility of such a fellowship, of a church,
the criterion of the genuineness of experience and faith. Even for theologians of
liberal tendency the possibility of forming a congregation and founding a cult is
the mark that a real religion is present: ``Only in the capacity to found a cult
does the genuine religion, borne by unbroken conviction, make itself known''
(Troeltsch). What matters is then not what the individual has lived through and
uttered, but whether a congregation can gather and a tradition form about definite
facts and ideas (Kaftan).\textsuperscript{35}---In modern theological discussions
the invocation of the congregation therefore plays a great role.---A dominant
tendency in the history of religion is so far in agreement with this, in that it
designates its task as that of presenting the faith of peoples, of the great
masses, and not the faith of the individual great personalities, however
significant this may be.

\medskip

\noindent\textbf{34.} It is more individual types that we have hitherto brought
forward. I now pass over to the great differences that have found their expression
in the highest popular religions. These types have developed under reciprocal
action between older religious traditions, racial characters, and the experiences
of centuries. There appear here especially two chief types of deep and abiding
significance for spiritual development.

The one chief type is characterized by a need to raise oneself above the struggle
for existence, \emph{to be freed from change and opposition, from all ``twoness''
and difference.} With difference follows suffering, with change follows unrest,
and the suffering and the unrest bring it about that one sets oneself as an end the
attainment of better states---but these are attained only to prepare a so much the
more bitter disappointment! This whole surging is to be halted and abolished. Only
in the eternal and unchangeable is there rest! But since all ideas are drawn from
experience's world of difference and instability, no positive expression can
characterize the eternal and unchangeable state that is yearned for. And since all
change and movement is an illusion, this very yearning comes at last---when that
state is attained---to stand as an illusion. One yearns to be free of all yearning.

This type recalls partly the first (§30), partly the third (§32) of the individual
types. \emph{Indian religiosity,} especially as it was concluded in Buddhism's
doctrine of Nirvana, is the most characteristic example of it. Nirvana is not pure
nothing. It is a form of existence to which none of the properties that
ever-shifting experience presents applies, and which therefore, in relation to the
states we know from experience, comes to stand as a nothing. It is a liberation
from all need and suffering, from hatred and passion, from death and birth. It is
attained by the highest degree of concentration of thought and will of which the
human being is capable.\textsuperscript{36}---In \emph{Neoplatonism} and in
\emph{the mysticism of the Middle Ages} kindred traits are found.---In all
religiosity the opposition between the unchangeable and the changeable plays an
essential role; but in this chief type it is the wholly decisive thing.

With the second chief type it is likewise a matter of redemption from strife, but
the opposites do not here stand equal, for the human being attaches himself to one
of the contending powers, so that redemption is reached through this power's
victory, and not merely through the human being's being raised above the sphere
where there can be any question of strife at all. Life is a strife between good and
evil powers. This strife shall and must be carried to its end; it has a positive
reality, and everything depends on which side the human being has attached himself
to while the strife goes on.---\emph{The doctrine of Zarathustra} and
\emph{Christianity} are the great representatives of this chief type, whose
significance rests on the decisive weight laid on \emph{historical development} and
ethical striving. Time and \emph{life in time} acquire here undoubted reality. (Cf.\
§11.)

There are naturally many intermediate forms between the two chief types and many
combinations of them, already for the reason that the one chief type historically
works upon the other. The Indian-Greek type has through Neoplatonism worked upon
the Christian type, and the collision and reciprocal action of the two types is
noticeable in Augustine's thought, so influential for the whole European
development. The opposition between the unchangeable and the changeable has a
tendency to thrust itself forward in place of the opposition between good and evil.
The proposition that God is the most changeable of all beings, which has been
uttered by a religious personality within Christianity, could not have been uttered
by an Indian or a Greek, scarcely even by Augustine.

Common to both chief types is that there is experienced an opposition between
times in which the goal (whether it be redemption from the strife, or victory in
the strife) stands as reached or at least as certain and present, and other times
in which darkness, dullness, and hopelessness prevail. And these states will
mutually heighten one another by \emph{contrast-effect:} for him who has known the
high and bright states, where one feels oneself in unity with the highest, dark and
depressed states become doubly painful and comfortless, while conversely the former
acquire a double radiance by their opposition to the latter. Here there is use for
faith (as faithfulness, see §29), in order that there may be continuity in the
personal life. The bright moments easily come, in the dark moments, to stand as
illusions. What matters then is to hold fast the right connection, in firm
conviction that even the dark times can purify and strengthen, and that all value is
not extinguished even in times when one must ``do without all consolation.''---That
these oscillations must be stronger and more frequent within the second chief type
than within the first is a natural consequence of the two types' distinctive
characters severally. Within the second chief type, where the temporal relation
acquires decisive significance, there will be a far greater difference between the
different moments, and there will perhaps even be a need to live through and endure
extreme states, to oscillate between opposite moods, in order that the consciousness
of victory may be the more glorious. In cult this need will especially manifest
itself.

\medskip

\noindent\textbf{35.} To the concept of faith corresponds the concept of God. From a
purely theoretical (epistemological-metaphysical) standpoint the concept of God can
only mean the principle of existence's continuity and thereby of its
intelligibility. From a religious standpoint God means---as faith's object---the
principle of the persistence of value through all oscillations and all struggles---if
one will: the principle of faithfulness in existence. If the religious problem could
be completely solved, it would have to be demonstrable that these two principles---the
persistence of energy and of value---quite coincide. The analogy between them is
clear; but from that there is a great step to their identity, and it is a great
question whether this step can be taken in such a way that it leads to a
contradiction-free concept. We shall in another connection come back to this question.

Religious experience may be said to be an experience of ``God,'' in so far as it leads
to a faith in the persistence of value in spite of everything and through everything.
He who would experience ``God'' must train his eye to find a valuable kernel behind
reality's hard shell, and train his mind in patient hope that such a kernel will
always be able to be found. He who would work for ``God's kingdom'' must work to find,
bring forth, and assert the valuable. ``Thou feelest him each time thy heart beats with
holy desire; thou hearest him in thy best thoughts as an eternal voice.''

Of the justification for using the word God in this manner there holds what was
remarked above (§26) about the word religion. Psychologically the kinship between
different standpoints is often greater than it might purely dogmatically seem.

As regards the hypothesis of religion's essence as faith in the persistence of value,
the investigation in the section of the psychology of religion that we here conclude
has on the whole confirmed it, in that a discussion of the nature of religious
experience and religious faith has decidedly led us to see its probability. Religious
experience acquires its distinctive character by being experience concerning the
relation between value and reality, and religious faith acquires its distinctive
character by not merely itself being a firm, continuous direction in the mind, but
also by asserting a continuity of value through all existence's oscillations. But we
shall in a later section (D) return to it in order to test how far the most significant
forms of positive religion can be traced back to it.

\begin{center}---\end{center}

\section{The Development of Religious Ideas}
\label{sec:development}

\subsection{Religion as Drive}
\label{ssec:drive}

\begin{flushright}
\textit{Various enough have been the religious symbols, as men stood in}\\
\textit{this stage of culture or the other and could worse or better}\\
\textit{body-forth the Godlike.}\\[2pt]
\textsc{Carlyle}.
\end{flushright}
% Høffding gives the Carlyle epigraph in English; reproduced verbatim.

\bigskip

\noindent\textbf{36.} The philosophy of religion has neither any ground nor any
capacity for investigating religion's historical beginning. Its task is to treat
the religious problem as this presents itself for us now, under our present
conditions of culture. To that end it uses, among other things, the history of
religion; but it interests itself more in the question whether, in spite of the
changing forms of religion, a constantly operative foundation for religion can be
pointed out, than in the historical origin of religion. And even if one could
point out the historical beginning of all religion, this would hardly have any
philosophy-of-religion significance. From a being's first beginning one cannot
always learn anything with respect to its real nature, for the transformations and
displacements that take place during development can bring about properties that
could not be observed in the first beginning. A being's real nature consists in the
law of its whole development.

For the history of religion too, religion's first beginning in the human race will
probably continue to stand as an unsolved problem. Even the lowest-standing savages
we know have run through a long development. One merely cuts the knot if one assumes
that any of the now existing religions should be the primal religion. And at any
rate this primal religion, if it could be found, would not necessarily be the
``real'' religion, so that if it could for example be shown that religious
development in the human race began with fetishism, then all religion should
``really'' be fetishism. It is an inference of the same kind as that according to
which, by Darwinism, the human being is ``really'' an ape. Inferences of this kind
are drawn both from theological and from anti-theological standpoints. It rests on a
complete misunderstanding of the nature of all development. From the theological
side one has for the most part thought it necessary to assume a primal revelation to
the first human beings, a perfect religion which they lost by a falling away; and
this primal revelation one then constructed out of one's own religion, which was to
be its restitution. The fetishism theory has after all a more natural stamp. For, of
two things one, it is nevertheless easier to understand that the more perfect can
develop out of the more imperfect than that the more imperfect should spring from the
perfect. For when the perfect contains the possibility of, or the germ of, the
imperfect, it is not perfect; whereas the imperfect can, by supplementation or by
transformation, develop in the direction of perfection. History shows us no first
beginning of religion, but a series of different, lower and higher forms of religion,
a series that does not advance evenly, but swings back and forth, even though through
all the oscillations a general direction of development can be traced, especially as
regards the religious ideas, the most easily accessible side of religion.

When we start from the description of religious phenomena given in what precedes, we
must expect that religion everywhere unfolds on the basis of a certain conception of
reality. The human being cannot have religious feeling when he cannot to a certain
degree arrange his observations of the world. Religious feeling is indeed determined
by experience of the relation between value and reality, and thus presupposes that a
reality is known. In investigating the development of religious ideas I shall proceed
by beginning with the simplest forms of religious idea and advancing to the more
composite, everywhere trying to explain to myself the transition from one stage to
another on the basis of purely psychological laws.

\medskip

\noindent\textbf{37.} In characterizing the most elementary forms of religion one has
often laid too great weight on the ideas, in that one has made primitive religion one
with primitive conception of the world. In spite of their great merit in the science
of religion, the English investigators, especially \textsc{Tylor}, \textsc{Spencer},
and \textsc{Frazer}, have nevertheless given way too much to the tendency to conceive
religion as a matter of knowledge, and in the elementary religions have predominantly
laid weight on certain ideas, their associations and their displacements. Elementary
or primitive religion was then described as belief in spirits. In contrast to this
the French sociological school, especially \textsc{Émile Durkheim}, has maintained
that belief in spirits is indeed found everywhere where there is primitive religion,
but that it is not the foundation but is derived. The foundation is everywhere formed
by a deep need to feel oneself at one with what can bear and secure life. It is by
virtue of this need that the human being distinguishes between the holy and the
profane, the fundamental religious distinction, which at the most elementary stage
known, for example among the Central Australians, is a distinction between what
conditions the tribe's persistence and welfare and what is indifferent or hostile to
it. The individual feels himself at this stage in solidarity with the tribe to which
he belongs, and therefore with the totem, the natural object that stands as the
expression of the tribe and its conditions of life. Through the tribe's cult its totem
is secured, nourished, and honoured, and through the individual's participation in
this cult he becomes a partaker in the holy and in its sustaining power. Cult is a
social power through which influence is exercised not only on the minds of human
beings, but also on the course of nature, on fertility, rain, and successful hunting.

Through the characterization Durkheim (in his work \emph{Les formes élémentaires de la
vie religieuse})\textsuperscript{37} has given of the most elementary religion known, I
find confirmation on two chief points of the philosophy-of-religion theory I maintain.
It shows in the first place that religion's foundation does not lie in the intellectual
domain, but must be sought in a need that is not bound up with clear consciousness, a
need for the securing and unfolding of life. And it shows in the second place that
when we try to interpret this religious need, it is to be conceived as a need to hold
fast a persistence of the valuable that is experienced, out beyond the limits within
which human beings can work for it.

Since the totem is common to the tribe and stands fast through the changing
generations, the idea of the forefathers naturally acquires a prominent place; their
souls are assumed to be constantly at work, and a belief in spirits forms on the basis
of the customs of the cult and as an attempt to interpret them. In its further
development this belief in spirits can become a conception of the world, whose traces
can be pointed out right up into the highest religions. This conception of the world,
which is found again everywhere on earth among human beings who stand at the lowest
stage of development we know, one is accustomed to call (after
\textsc{Tylor})\textsuperscript{38} \emph{animism.} Animism is a conception of the
world, not in itself a religion. But it acquires interest for the philosophy of
religion because it is the most elementary conception of the world that religion has
taken into its service. Its distinctive character consists in this, that
events---especially striking events---are explained by the intervention of souls, of
personal beings. Its arising is due in great part to the influence of dream-ideas: in
dreams the savage gets experiences that must stand for him as just as valid as the
experiences of waking life, and in these dream-experiences he himself and other
beings---human beings and animals, the living and the dead---appear in a freer manner,
more independent of space, time, and material conditions, than the experiences of
waking life led him to regard as possible. The dream-world, which here stands as
indubitable reality, is then consistently used to explain and supplement the world of
waking life. Consciousness is especially taken up with the spirits of the dead; in
their lifetime these beings had intervened more or less deeply in what went on; it is
then natural that one supposes oneself to trace their doings and workings now too,
since from the dreams about them one knows that they not merely continue to exist, but
even exist under more perfect conditions. The involuntary tendency to personify, to
conceive things after the likeness of ourselves and therefore to conceive natural
processes as personal actions, had to be supported by these ideas, even if these alone
should not be enough to explain animism.

Animism occurs among all tribes round about on the earth at a certain lower stage of
their development. It is the human being's most elementary philosophy. And even in the
highest, most fully formed religions, with their ideal conceptions of gods, exact and
unprejudiced investigation meets many traces of this circle of ideas.

\medskip

\noindent\textbf{38.} In totemism the individual's ideas are predominantly determined
by the tribe's. But within the social tradition individual processes of ideation can
nevertheless be traced, and such are indeed also a presupposition of the social
tradition's coming into being, even if it hides in the darkness of the ages. The
individual does not always hold to the tribe's fetishes (the objects to which a magical
effect is attached), but may choose himself a fetish on his own account. It is then a
single definite object, a special occasion, that leads the human being to assume a
power that decides whether something valuable for him happens or not. In such a process
of choice is contained the simplest formation of religious ideas we can imagine. The
choice is quite elementary and involuntary, just as elementary and involuntary as the
exclamation that is the simplest form of value-judgment. The object that is chosen must
in one way or another stand in connection with what strongly occupies the mind. It
perhaps awakens the memory of earlier events at which it was present or played a role.
Or it presents a certain likeness, often very distant likeness, to objects that
formerly helped in distress. Or it is quite simply the first object the human being
catches sight of in his state of strained expectation. It fixes attention upon itself
and is involuntarily placed in real connection with what is to come about---with the
possibility of attaining the goal one has set oneself. It may be both hope and fear that
leads to this choice; probably it is from the first predominantly hope, for the human
being is sanguine by nature. There is no reason why it should be fear alone that creates
gods. It seems, to be sure, to be a rule that evil beings are worshipped rather than good
beings; but they may be worshipped in hope that one can make them favourably disposed.

In such phenomena religion appears as \emph{drive.} The first manner in which ideas (as
different from sensations) appear is as elements in drive (according to the usage that
distinguishes between drive and mere instinct). Drive contains an idea of something that
can satisfy a need or in general cause pleasure, and the idea has here therefore
immediate practical significance. The thirsty man's idea of water is quite practical,
expresses the goal of the yearning and is held fast by the need or by the expectation.
Thereby the thirsty man's idea of water is different from the chemist's or from the
painter's. Religious ideas are religious only by this connection with need and
expectation, thus as elements in a drive. One and the same idea can therefore look
highly different from the religious standpoint and from the theoretical-psychological or
the historical-critical standpoint, and one understands the significance of religious
ideas only when one holds fast the connection with drive that is their real soil.
Fetishism in particular is understood only in this way. When the young Indian is to
choose his ``medicine'' (thus he calls his fetish or talisman), he withdraws into
solitude and gives himself up to fasting; the first animal he meets when he again leaves
his solitary dwelling becomes his ``medicine.'' A negro goes out of his tent in the
morning to set out on an expedition and sees a stone lying glittering in the sunshine;
involuntarily he takes it to himself and relies on its help. In such examples the
religious idea is formed or chosen by a kind of \emph{improvisation.}

There need be no connection between the different religious improvisations, although
habit and tradition will soon exercise their influence. The fetish is provisionally
merely the momentary dwelling of a spirit. It is what \textsc{Hermann Usener} has aptly
called a \emph{momentary god:} ``In full immediacy the individual phenomenon is deified,
without any generic concept, however limited, playing a part. The individual thing that
you see before you, and nothing else, is the god!'' As examples Usener refers to the
custom among the old Prussians and Lithuanians whereby the first or the last sheaf in the
field is treated as the dwelling of a god who had to be celebrated if the corn harvest
was to be good, and to the bouquet of Saint John's wort that young Lithuanian girls
plucked and set on a pole at the entrance to the house, and which was worshipped with
reverence. Both the sheaf and the Saint John's wort were originally fetishes. Among the
Greeks similar examples are found, as when a hero in Aeschylus swears by his
sword.\textsuperscript{39}

The tendency to form such momentary gods can constantly be traced in the involuntary
personifying and symbolizing that can make itself felt when one stands in a living and
interested relation to things in the surrounding world. One may for example find an omen
for the outcome of an undertaking begun in the circumstance that one succeeds in getting
the fire to burn in the stove: the momentary expectation that arises during the work with
the refractory fuel passes over by feeling-expansion to apply to the greater undertaking
with which consciousness is constantly occupied. With this expansion there then follows
the idea of the successfully produced fire as a kind of helping and encouraging genius.
When one notices the kind of momentary mythologizing that does not occur seldom, one has
in it a hint towards understanding the most elementary formation of religious ideas.
Only that what is for us now merely a fleeting meteor on consciousness's horizon becomes a
guiding star at animism's standpoint.

\medskip

\noindent\textbf{39.} Momentary gods can historically be pointed out only approximately.
For on the same occasion the human being naturally returns to the same idea of a god, and
this is then not produced anew in each single case, for example at each new harvest. The
gods then gradually appear each with certain constant properties and with dominion over
certain definite domains. By the constant property and the definite domain the individual
god acquires a firmer form in consciousness. Just as development is needed to raise
oneself above the individual moment, so development is also needed to raise oneself above
the individual special manner of expression and the individual special domain. But it is
nevertheless an advance when what is constant in property and domain leads out beyond what
is momentary in the simplest ideas of gods. Therefore \emph{the special gods} (A. Lang's
departmental gods, Usener's \textit{Sondergötter}) denote an advance in relation to the
momentary gods. Thought does not yet, however, raise itself above multiplicity and
differences, but assumes a principle, a guiding power, for each distinctive domain. A New
Zealand chief said to a European: ``Is there among the Europeans only a single one who
produces all things? But there is one who is a carpenter, another who is a smith, a third
who is a shipbuilder. And so it was in the beginning too. One produced this, another that.
Tane produced trees, Ru mountains, Tangaroa fish, and so forth.''

Usener has especially pointed out this stage in the development of religious ideas in the
older Roman religion, as it was before the Greek influence. In the liturgical books of the
Roman priests (the so-called \textit{indigitamenta}) were found lists of the special gods
who from ancient times were invoked each on his special occasion and in his special
capacity (\textit{dii proprii}, \textit{dii certi}). ``For all actions and states that
could be of importance to the human being of the time, distinctive gods have been created,
and they are named by specially formed words; not merely actions and states as wholes are
in this manner deified, but also all more prominent sections, acts, moments within the
actions and states.'' Thus in the cultivation of the field one invoked not merely the
earth (Tellus) and the god of fertility (Ceres), but twelve special gods: one for the first
working of the fallow, another for the second ploughing, a third for the last drawing of
the furrows, a fourth for the sowing, a fifth for the ploughing-in of the seed, and so on
and so on.---Among the Lithuanians and the Greeks too such special gods, restricted to a
single property and a very limited domain, can be pointed out, and---we may add---they are
found among nature-peoples round about on the earth.\textsuperscript{40}

Such special gods correspond---like the momentary gods---to a need in the religious
consciousness to know a power in its neighbourhood that is exclusively taken up with
bringing it the help that is needed. The old special gods often pass under altered names
into the new religions. Christianity's saints often appear with properties and functions
that show that they are the heirs of the special gods; often a saint is worshipped at the
same place where in antiquity a corresponding special god was worshipped. The saints have
their departments like their predecessors, and they satisfy, like these, the need for
professional divinities. It is taken from life when a young Roman girl in one of Ludvig
Bødtcher's Italian poems answers the question whether she is not afraid of earthquakes:
``Ah, Madonna graciously protects our little house well enough, and when the danger really
approaches, St.~Emidius helps.'' St.~Emidius is precisely the special protector against
earthquakes.---An interesting example of continuity in the development of religious ideas
Usener gives in the idea of the mother with the child. She is first worshipped as a special
goddess under the name of ``the child-nourisher'' (Kurotrophos)---without a proper name.
Later this name (Kurotrophos) became an epithet of different goddesses (Leto, Demeter), and
her worship glided over into the medieval worship of the Madonna with the child. This
continuity testifies that precisely towards this figure of a god, which represented motherly
love, there has been present a distinctive need for adoration.

In the opinion of some, the Jewish people's oldest god is to be regarded as a special god.
``The Jahvism of the earlier time,'' remarks a historian of religion, ``was a power hostile
to culture. Jahveh meets us above all as destroyer: storm-god in nature, war-god in the life
of the people---the god to whom, after a victory won by fighting, everything living is
consecrated to death!'' Only through reciprocal action with the ideas of gods of other
tribes is the Jewish idea of God said to have acquired its universal significance and taken
up elements that could make Jahveh a divinity for a people of culture.\textsuperscript{41}

\subsection{Polytheism and Monotheism}
\label{ssec:polytheism}

\noindent\textbf{40.} The drive-like personifying that is characteristic of the
religious consciousness in its most elementary form does not yet produce ideas of
personal beings in the proper sense. A personal being possesses at once several
different properties; it is the bond of unity between them, just as it is the bond
of unity between the different moments in which its life unfolds. The formation of
ideas of such beings presupposes a spiritual development that cannot be present at
the most childlike stages. It is a distinctive mark of the ideas of children and
savages that they attach themselves to a single side and property of the thing and
recognize it by that; hence the baroque juxtapositions and the many confusions to
which both the child's and the savage's speech testifies. There must be acquired a
capacity to raise oneself above the momentary and the special and to connect the
individual experiences into a whole, for a personal being is not exhausted by a
single situation or a single property. It is the capacity to form typical individual
ideas that must be present. A typical individual idea is an idea that fits an
individual being in all the different states in which it appears. Great art is
needed to form such ideas when they are to hold for beings that are of a more
composite nature and in constant development, and a complete conclusion of such
ideas is often not possible. Our ideas of personal beings as a rule acquire an
artistic rounding-off that cannot be fully grounded by observation, and our knowledge
of a personal being's distinctive character is rather a faith than a knowledge.

Just as the transition from individual observations and individual ideas to typical
individual ideas is psychologically regarded one of the most significant transitions
within the life of ideas, so the transition from momentary and special gods to
properly personal gods is one of the most significant transitions within the history
of religion. It is the transition from animism to polytheism. A more definite
distinction is now drawn between the divine beings and the natural phenomena to which
they are attached. The figures of the gods themselves acquire a richer and deeper
content, and only now can there be any question of a proper faith, since there is now
a relation of distance between the god and his manner of revelation which could not
be present at the more elementary stages of development. There lies here one of the
rare advances in the history of religion that benefits at once science and faith. By
the more definite difference between the gods and the individual natural phenomena,
these latter can the more readily become the object of objective observation and
investigation. If science stands in any debt to religion, this was incurred rather at
the transition from fetishism to polytheism, not, as has been supposed, at the
transition from polytheism to monotheism. But first and foremost that transition
benefited religious faith, since the figures of the gods could now become deeper and
richer and be attached in a more inward manner to the life of feeling, when individual
natural phenomena no longer stood as the immediately corresponding expression for them.
The human being has now begun to live in an invisible world.

The significance of this epoch in the history of religion has already been strongly and
aptly emphasized by \textsc{Auguste Comte}. In later times \textsc{Hermann Usener} has,
on the basis of interesting historical and linguistic investigations, strongly stressed
it. When this excellent historian of religion reproaches the philosophers of religion
for letting religious development begin with polytheism, Comte's example already shows
that this reproach is unfounded. And no less unfounded is the remark that the
philosophers ``treat concept-formation and the gathering of the individual into species
and genus as a self-evident and necessary process in the human spirit.'' In Aristotle
already the problem of individual concepts appears, and in more recent philosophy it
has since John Locke's time been discussed again and again. When in psychology a
distinction is drawn between individual ideas (corresponding to a single trait, a single
property), concrete individual ideas (corresponding to a phenomenon in a single, definite
situation), typical individual ideas (corresponding to a series of different observations
of the same phenomenon), and general ideas (concepts of species and genus, corresponding
to a series of observations of related phenomena)---there is thereby given a satisfactory
framework within which the phenomena of religious history stressed by Usener can also find
their place. The transition from the one kind of ideas to the other naturally does not
take place without hindrances and always presupposes favourable inner and outer conditions.

Usener has especially shown that only at a certain stage of development, namely at the
coming-into-being of polytheism, do the gods acquire proper names. The special gods are
named only adjectivally, after the individual properties with which they appear. As an
example may be named Apollo, who is characterized as the light-bringer, the god of
prophecy, the god of song, the reconciler, or the god of healing. Usener ascribes to the
formation of proper names for the gods a decisive significance as a condition of the
transition from momentary and special gods to personal gods. But there is here as
everywhere a constant reciprocal relation between idea and word. The proper name becomes
intelligible only when several properties and states can be gathered up in one idea. The
word is a help and a support for holding fast and developing a result won in the life of
ideas, but cannot be the exclusive cause. Usener therefore also expresses himself somewhat
waveringly. Sometimes he lets ``concept-formation stiffen into a proper name,'' or he lets
``language help to pave the way''---here, then, the reciprocal relation between the
formation of ideas and the formation of words is acknowledged. With the bare word a
significant period in the history of religion cannot begin. The word cannot be the
beginning or be in the beginning.\textsuperscript{42}

Historically one naturally cannot sharply point out the beginning of polytheism. When
momentary and special gods already presuppose a ``personifying'' tendency and capacity,
and when---especially on account of the influence of habit and tradition---pure momentary
and special gods can hardly be pointed out, it already follows from this that there are
many transitional links between the stages that we here in theory place in sharp opposition
to one another. The same will show itself to hold of the relation between polytheism and
monotheism.

\medskip

\noindent\textbf{41.} Neither the development of ideas nor the formation of words gives in
and for itself an exhaustive explanation of the transition to polytheism. There constantly
co-operates an influence from the life of feeling. The historians of religion strongly
emphasize \emph{the conservative, restraining influence} of cult, and thereby of the life
of feeling, for in cult feeling has its immediate foothold. When a cult in a region or
among a people has been organized about a special god, it becomes a hindrance to the
formation and acknowledgement of more composite ideas of gods. The influence that cult has
on the religious life of ideas appears strikingly in the very word God, when we hold to
those etymologies of this word that are said to have the greatest probability on their side
linguistically, according to which the word properly means ``he to whom sacrifice is made,''
or ``he who is invoked.''\textsuperscript{43} The connection between the idea of God and
cult brings it about that a thoroughgoing alteration in the former must effect an alteration
in the latter, and here now, by the experience of all ages, the organized cult offers a
strong, often violent resistance. Old ideas maintain themselves long by being attached to
old customs, and old customs can maintain themselves longer than the corresponding ideas. In
a village church in Zealand it is said to have been the custom even in the nineteenth century
that one made a bow when one passed a certain place on the church wall, but no one knew why
one did it, until on the scraping away of the whitewash a picture of a saint was found on the
wall: the custom had thus survived Catholicism by three centuries; it was a part of the old
cult that had maintained itself. Here it was now a matter of mere habit. But the life of
feeling will on the whole only slowly leave the bed once formed for it. It has accommodated
itself to the transmitted ideas, and a transitional time full of unrest and disharmony must
be lived through before the accommodation to the new ideas can be accomplished. In such a
transitional time the two tendencies in the life of feeling, the tendency to remain in the
old bed and the tendency to expansion, wage a hard struggle with one another. Or rather: the
struggle is waged between the old and the new feelings' tendency to spread over and colour the
whole of consciousness; the feelings attached to what is transmitted have namely also an
expansive tendency, and they will constantly try to colour and transform the new ideas which
they cannot wholly displace;---in case of need the old ideas are transformed after the new,
when only in that way can they be preserved.

The restraining and inhibiting influence of feeling on the life of ideas is, however, only the
one side of the matter. The new experiences can work so penetratingly that an alteration of
the ideas takes place through \emph{selection, idealization,} and \emph{combination.}---%
\textsc{Usener}, who otherwise lays the chief weight on the development of words and ideas,
therefore also remarks that the different special gods cannot have the same value for
consciousness: ``The god of the heavenly light that bestows life and blessing, the tutelary
god of the house and of the peace of the house, the saviour, the averter of all evil, must each
have far greater significance than a god who makes the harrowing or the weeding succeed well.''
The special gods that awaken the strongest and most lasting feelings will come to rule over the
others.---As objects of particular attention they will acquire a distinctive place; comparison
with other divinities will be kept away, and consciousness will on the whole not dwell on what
could leave them standing as limited. There thereby takes place an involuntary idealization.---%
With the selection and the idealization there naturally combines a combination, in that all
valuable properties and effects that present any likeness to or contact with a god's original
attribute will be referred to him. He stands once for all as the author of something especially
valuable, and the new valuable thing that is experienced is then involuntarily referred to him
too.---

The history of the light-gods is an interesting example of this development. Light's great
significance for everything living, by effecting the more vigorous and healthy course of the
life-process and by making activity and security possible, and its ideal significance as a
symbol of truth and justice, condition transitions of ideas of the most manifold kind, which
are nevertheless always borne by the immediate experiences of the value of light and of the
forces symbolized by light.

Quite especially will those gods that stand as guardians of common goods---as the protectors of
the family, the clan, or the people---become the object of these different processes. When the
human being in his life together with others and in his own inner life experiences the
opposition between good and evil, and when this opposition acquires a predominant place in his
consciousness, he will also transfer these ideas in idealized form to his gods. The idealizing
personification will become of great ethical significance by creating radiant divine models for
human life. The transition hereby takes place from nature-religion to ethical religion, a
transition that has rightly been called the most important in the whole history of religion, but
which takes place according to the same laws as all other transitions. Just as little as human
beings would have believed in the thunder-god if they did not experience the power of thunder,
just as little would they have believed in good gods if they did not experience the good as a
reality in life. The idea of the gods as the bearers of the ethical values in existence
presupposes that the human being has himself come to know these values. (Cf.\ §26.)---

Besides the inhibiting, selecting, idealizing, and combining influence that feeling can exercise
on the development of ideas, the significance of \emph{contrast-effects} in the domain of
feeling must further be emphasized. The contrast between two moods will react back upon the
ideas in which the two moods severally find expression, or by which they are motivated. If an
ideal idea has thus been formed, the ideal will be further raised in power by the opposition
between the perfection represented and the human being's own limitation, and conversely this
limitation will be represented the more glaringly, the brighter the mood that the idea of the
ideal had called forth. This moment is therefore of importance not only in the development of
the figures of the gods, but also in the development of the consciousness of sin, which plays so
great a role in many religions.

\medskip

\noindent\textbf{42.} Outside the circle of the properly divine ideas, the belief in
the transmigration of souls presents a good example of the selecting and transforming
influence of the feelings. The idea that souls after death pass over into other bodies
has hardly arisen in the first instance from religious motives. It hangs together with
general animism and has arisen among different peoples round about on the earth at a
very early stage of their development. It is frequent among East Asian peoples; it is
found among the Central Australians, the Guineans, and the Zulus; among the Greenlanders
and among North American tribes we meet with it too.

Now we find that it plays a very great role in Indian religion after the Vedic age. In
the Upanishads it is of great significance, while it is not found in the Indians' earlier
religion as this appears in the Vedic poems. It has then been supposed that the doctrine
of the transmigration of souls arose because religious thought sought to explain to itself
the great inequality which, already as regards original endowments and conditions, shows
itself to be present between human beings in an ethical respect. What could not be
explained as an effect of the human being's conduct in this life, the old Vedanta thinkers
are then supposed to have explained as an effect of his conduct in earlier existences. By
a similar road the Pythagoreans and Plato are supposed to have come to the mythical idea of
the transmigration of souls. According to this conception, which is espoused by
\textsc{Deussen},\textsuperscript{44} the need for redemption among the Indians has partly
created the idea of a transmigration of souls, in order to explain the inequality in inner
and outer conditions, and partly developed the idea of absorption into Brahman (and later of
Nirvana), so that the two ideas arose independently of one another.

This conception in the excellent connoisseur of Indian philosophy hangs together, however,
certainly with his overlooking animism's constant after-effects, which can be traced in the
exalted Vedanta doctrine itself. In the Upanishads' emphasis on the significance of the
dream-state one notices one of the fundamental thoughts that are distinctive of animism.
``In sleep,'' it is said, ``the spirit casts off everything that belongs to the body and
floats up and down, creating for itself, like a god, all manner of shapes.'' Indeed, even the
naive consequence of the animistic dream-theory is not lacking, that one must not too hastily
wake the sleeping, since ``it is impossible to heal a human being whose spirit has not found
its way back to him.''\textsuperscript{45}---When even the Vedanta doctrine itself thus
contains traces of very elementary ideas that are akin to the ideas of the transmigration of
souls, then the most correct explanation of the great role that the doctrine of the
transmigration of souls came from a certain point of time to play is perhaps that this doctrine
was not produced all at once by an artificial speculation, but that it arose because older
ideas (perhaps among subjugated peoples), to which one had hitherto ascribed no religious
significance, now---as a consequence of an alteration in the religious life of feeling---were
taken up into the circle of the properly religious ideas and interpreted as the expression of
religious experience. Under the overwhelming impression of life's evils and restless doings one
could explain to oneself the unhappy conditions of life only as a consequence of the earlier
existences to which popular belief testified, and one besides took so dark a view of things
that not even death brought peace, when there was no security against a new life in finite form
following after. \textsc{Richard Garbe} thus connects (in his work on the Sankhya philosophy)
the significance that the old doctrine of the transmigration of souls now acquired with a
remarkable change in the whole feeling for life of the ancient Indians. ``In the old Vedic time
there ruled in India a bright view of life, in which we see no germ of the later idea [of life
as a burden] that dominated and cowed the whole people's manner of thought; one did not yet
feel life as a burden, but as the greatest of all goods, and one hoped for a life after death
as the reward for a pious life. All at once there enters---without our gaze being able to find
transitional stages---in place of this childlike joy in life the conviction that the
individual's life is only an agonizing wandering from death to death.''\textsuperscript{46} If
this explanation is the right one, there lies here a classical example of how an alteration in
feeling works selectively, idealizingly, and intensifyingly upon the ideas. At the same time the
contrast-effect here comes out clearly: for in opposition to the constant life-process, not
interrupted even by death, with its unrest and despondency, the complete abolition of temporal
and finite existence had now to become the goal. The doctrine of redemption of the Vedanta and
of Buddhism unfolds in a definite relation of contrast to the doctrine of the transmigration of
souls.

A striking analogy to this course of development is presented by the development of Greek
religion in the time between Homer and Plato. Here too it was under the influence of a new and
darker view of life that the belief in the transmigration of souls was taken up and applied
within the circle of religious ideas, and the concern for the soul's fate in another world, a
concern that was unknown in Homer, became prevalent in wide circles.\textsuperscript{47}---The
idea of personal immortality has on the whole hardly had religious significance in the first
instance; it is due to the circle of ideas from which animism sprang, and has only later, under
the influence partly of an altered feeling for life, partly of ethical ideas, furnished motives
for the development of religious ideas.

An instructive counterpart to the significance that the doctrine of the transmigration of souls
acquired for the Indians and the Greeks at a certain point of their development is afforded by
this doctrine's appearance among the Egyptians. The Egyptians assumed, indeed, that the soul
after death could pass over into different shapes; but this transmigration of souls was not at
all placed in connection with ethical ideas.\textsuperscript{48} One has here had no use for a
motive that became of such great effect in Indian and in Greek religion.

\medskip

\noindent\textbf{43.} I now return to polytheism. Its development presents a swarm of
psychological and historical processes which the history of religion does not give sufficient
material to illuminate, and which psychological analysis, even were it more perfect than it is,
would hardly ever be able to unravel. What can be traced nevertheless testifies that in the
religious domain no other psychological and logical laws prevail than in the other domains of
the life of the spirit.

It might seem to be a psychological absurdity in polytheism that the human being here feels
himself dependent on several different divine beings. The one of these seems as though it must
come to stand in the way of the other; the thought of the one seems as though it must displace
the thought of the other. The religion of drive already, with its momentary and special gods,
could call forth this astonishment.---The solution lies in the fact that the human being in the
moments of lived experience and emotion is so completely absorbed in what now stands before the
imagination that no juxtaposition and comparison, perhaps not even any recollection of other
gods, takes place. There is indeed on the whole in the spiritual life a rhythmic alternation of
immediate absorption and of reflection. In the moments of reflection we can juxtapose and compare
the ideas that have made themselves felt in us in the moments of strong and deep moods. There may
then perhaps show themselves contradictions that must be abolished, and that we strive to abolish,
in so far as new lived experiences do not seize us so strongly that we forget reflection's
difficulties and problems. This holds at all stages of religious development and for all religious
standpoints. During this rhythmic alternation the one wave will nevertheless be able to acquire
influence on the other. When it is then real difficulties that reflection brings forward, there
will slowly take place an alteration in the character of the lived experiences, just as conversely
reflection is altered by the experiences of the enraptured moments. There will always be an
involuntary tendency to bring about harmony between the ideas that in the different lived
experiences have come to stand as the expression of feeling's object. On polytheism's ground the
solution will here lie in the idea of a race of gods or a world of gods, within which each
individual god has his definite place. It will be so much the easier here to find a solution
because there stand ready to hand obvious analogies from human relations. Partly the family, partly
the state can furnish models for the harmonious co-operation of personal beings, and it is not
least through these analogies that human beings' ethical-social development acquires influence on
the ideas of gods. The relation between gods and human beings too is presented by the help of these
analogies, as when the god is called the father and lord of human beings, and when the world is
represented as a great state, common to gods and human beings.

Nowhere does the process of religious ideation run its course under the exclusive influence of a
people's own experiences. The ideas of other peoples work along with them by reason of peaceful
intercourse or conquest. The different worlds of gods collide when the peoples collide, and a
historian of religion has even set up the proposition that ``religion has never at any time
developed except when several different religions have come into contact with one
another.''\textsuperscript{49} Thereby the course of development of the religious ideas becomes
still more complicated and obscure. The possibility of combinations and assimilations,
contrast-effects and expansions is thereby increased beyond survey. Here lie a series of problems
for the history and psychology of religion of the future.

\medskip

\noindent\textbf{44.} The idea of a world of gods, a divine kingdom, is a station on
the road from polytheism to monotheism. When one has found a quite distinctive
difficulty in the transition from polytheism to monotheism, it is because one has looked
at the matter too externally. Within polytheism there already expresses itself in many
ways that need for practical and theoretical concentration whose full consequence is
monotheism. And the same psychological laws that lead from the religion of drive to
polytheism lead from this to monotheism---in so far, that is, as outside dogmatic
speculation (and even there!) there is any real monotheism in the word's strictest sense.

The transition to monotheism can take place under two forms, which nevertheless pass over
into one another through many nuances.---In the first place, within the world of gods a
single god can become predominant, so that he is raised high above all other gods and
above other peoples' gods and at last comes to stand as the only god. With this there
will as a rule also follow a purification and deepening of the idea of God; or conversely,
a purification and deepening of the idea of God leads in a monotheistic direction.---In
the second place, an idea of the divinity can be developed which does not particularly
support itself on a single figure of a god that is deepened and enriched, but is grounded
on the divine in all gods---on that in the gods which first makes them gods (to use
Plato's expression, \emph{Phaedrus} p.~249 C), and which in the different gods is
represented in different ways.
% The e-book prints the dialogue title as "Fajtros"; Faidros/Phaedrus is meant.

In the first direction the development goes among Assyrians, Babylonians, Egyptians, and
Israelites. Jewish prophetism presents the best known, the most fully carried through, and
the historically most significant monotheistic process of development.

The development of Jewish monotheism is one with Jahveh's development from special and
national god to universal god. It is due to the deep effect that the historical events
exercised on the minds of the prophets. The prophet stands in the midst of his people and
of its traditions; but in contrast to the people and its spiritual and worldly leaders, it
is the people according to its ideal essence and according to its constant persistence
beyond the momentary conditions of the time that is the object of his concern. With
ecstatic absorption in the great world-events he combines the capacity to assert in mighty
images the ideal conception of his people and its significance which he has drawn out of
its history. A distinction can be drawn between a predominantly world-historical and a
predominantly ethical consideration as the foundation of the development of Jewish
monotheism; but common to both is that it is the very endeavour to hold fast under new
conditions the transmitted religious ideas that leads to a higher, more ideal development
of them: only by becoming a world-god instead of a national god could Jahveh avoid sharing
the fate of his people in its downfall as an independent people; and since he came to stand
as the upholder of higher ethical ideas than had hitherto been attached to the idea of him,
that superiority became possible which stood in so cutting an opposition to the fate of the
people whose national god he had been.

The historical consideration first comes forward in Hosea and Amos (c.~760 B.C.). It was
not---they teach---by the help of their own gods that the Assyrians conquered, but by the
fact that Jahveh wished through their armies to punish his disobedient people. The great
unknown prophet who is usually called Deutero-Isaiah (c.~540 B.C.) goes a step further:
the people of Israel suffers not merely for its own sins, but in order that it may be
purified to bring peace and salvation to all peoples. There forms for the prophet's
intimating vision the picture of a figure who wins and brings the highest through the
deepest lowliness and abasement, because in inwardness and in fellow-feeling there lies a
greater value than in the great unfolding of power among the conquering world-peoples. The
historical and the ethical consideration here pass over into one another. In Jeremiah
(c.~620 B.C.) the ethical consideration in particular is carried out, though Amos already
knows it too. As the God of righteousness, Jahveh must rule everywhere where righteousness
holds or is to hold---thus he must be the God of the whole world, who is not limited by
national barriers. And since the true relation to Jahveh is an inner relation in the heart,
outer forms and their difference become at last without significance. The time shall come
when everyone bears the law in his heart, and already now the prophet praises other peoples'
faithfulness to their gods in contrast to Israel's unfaithfulness to Jahveh. Both the
universality of the ethical ideas and the inwardness they demand lead out beyond the national
limitation.

Jesus of Nazareth set up no new concept of God in relation to that which had worked its way
forward in the great prophets. But he himself stood with his personality as surety that the
time had come to which the prophets in intimating visions and thoughts had looked
forward---the time when an inward and infinite relation to God could come into force. By the
deep movement in the minds of human beings that the picture of Jesus's personality awakened,
a movement that let all outer conditions and differences vanish in value, it became possible
that monotheism could become a popular religion on a larger scale. The one true God was a God
such as the one to whom a human being like Jesus could testify, and whose emissary he could be.

In the second direction the development takes place among Indians and Greeks.---Among the
Indians it is not any single god that becomes sole ruler, but there is a going back to the very
force that makes all gods what they are: the inward yearning and need that finds vent in prayer
and sacrifice. By an extraordinarily remarkable development of thought, that which drives the
human being to worship gods comes itself to stand as the true divine power.

Bráhman means according to some swelling, thereafter yearning, then prayer, and at last the
principle of existence itself, in that the idea passes over from the movement to its goal, from
the need to its object;---according to others it means the magically creative word. If it is
asked more closely (as the gods themselves ask in the \emph{Upanishads}) what Bráhman then
properly is, no other answer can be given than that it is one with Atman, the soul, which the
human being knows from his own inner life. There lies here, as we have earlier (§15) seen, the
first attempt at metaphysical idealism (c.~800 B.C.). The gods of the popular religion stand
under Bráhman as its forms of revelation. The relation of the Indian gods to Bráhman recalls
the relation that according to Plato prevails between the Greek gods and ``the ideas.''---This
relation of subordination was acknowledged only after serious struggle. It was provisionally as
a ``secret doctrine'' (this the word ``Upanishad'' is said to mean) that this monotheistic
tendency came forward, and it was, as Deussen conjectures, first developed by thinkers who had
come from the warrior caste, while the Brahmins (the priests) long held fast the literal
understanding of the transmitted ritual and only reluctantly went over to an allegorizing that
made it possible to unite the new doctrine with the old cult.---Among the Greeks the
monotheistic tendency comes forward especially in the philosophers. \textsc{Xenophanes}
(c.~500 B.C.) combats polytheism partly through a comparison of the Greek gods with other
peoples' gods, partly through a criticism of anthropomorphism. Each people represents the gods
in its own shape: the negroes' gods are black, the Thracians' blond, and if oxen had gods they
would represent them in the shape of oxen. In his zeal against all finitizing and degrading
ideas Xenophanes raises himself to the thought of a principle of unity that cannot be the object
of vivid representation. But he also applies an ethical criticism to the polytheistic popular
faith, and it is especially from this side that Plato later continues his work.

In India the thinkers' endeavours after unity were outstripped by the popular need for a manifold
representation of the divine, and in Greece the result of thought was likewise unable to become
the content of a popular religion. This became possible only where---as in Jewish prophetism and
its conclusion in Jesus of Nazareth---the thoughts are developed through the people's own
experience, its own fate, and where they are woven into the picture of a personality who stands
as a revelation of the highest inwardness and love and thereby himself becomes a symbol of the
value of the thought for which he fights.---

The psychology of religion cannot here, any more than on many other points, get further than to
point out points of view that can make the transition from polytheism to monotheism intelligible.
But there lies here no special psychological mystery, as has sometimes been maintained in a
somewhat advertising manner. There may therefore both psychologically and historically be riddles
enough left.

But when is the transition from polytheism to monotheism accomplished? When is there a pure
monotheism, and is such a thing possible at all? When one represents the divinity as appearing in
several ``persons,'' and when it stands in opposition to a ``world'' that is its limit and its
hindrance, can the concept of God then properly be said to be purely monotheistic? When ``the
world,'' the other thing that exists besides God, is a real reality, not a nothing, then it has a
force that is \emph{not} vanishing in comparison with the divine force, and is so far itself a god
within its limits, just as well as ``God'' within his. The same holds when one assumes a ``free''
will (in the word's metaphysical, not in its psychological or ethical sense) or believes in a
devil. And since assumptions such as these are found in all popular religions, there is no reason
to lay so great weight on the opposition between polytheism and monotheism as one often does. Just
as there are monotheistic tendencies in polytheism, so there are polytheistic tendencies in
monotheism. \textsc{William James} rightly maintains that popular theism has a polytheistic
character.\textsuperscript{50}

The most energetic attempts to carry through a monotheism have been made by philosophers: by the
Vedanta philosophers, by the Eleatics, and by Spinoza. But it has precisely in these attempts of
thought shown itself clearly that the problem of the relation between unity and multiplicity in
existence cannot be solved by making one of the two terms of the relation the only one. One will
then never without witchcraft be able to get the other term forth. Whether the ``reborn'' knowledge
that some theologians have spoken of would here be better placed than natural knowledge will show
itself only when the logic of this ``higher'' knowledge is one day written, which it is not yet.---%
But a closer discussion of this would lead us back to the epistemological philosophy of religion.

\medskip

\noindent\textbf{45.} The two tendencies that appear, more or less pronounced and in
different mutual relation, in every religious circle of ideas---the polytheistic and the
monotheistic---point back to the two tendencies in the nature of the religious need that
we have already had occasion to point out (§18).

There expresses itself on the one hand a need for gathering and concentration, for absolute
devotion, for getting beyond all strife, resistance, and change. This need comes to
unfolding in mysticism and monotheism; it meets with the intellectual need for a conclusion
in an absolute principle of unity.

But in the religious consciousness there expresses itself at the same time, more or less
energetically, and in rhythmic alternation with that first need, the need to have a
fellow-combatant at one's side in the midst of the strife, a fellow-combatant who from his
own experience knows what it is to suffer and to undergo resistance. This need will
counteract the carrying through of monotheism. And it will be supported by various other
motives. There will stir a need for vivid representation, a need for images, which leads to
the formation of limited ideas of the divine---or rather to ideas of the divinity as
limited, since only the limited can be shaped into an image. It is what has been called
``the antinomy of religious feeling,'' that the infinite and exalted is nevertheless
represented in limited form. The need for images can be so strong that it becomes a shyness
of infinity and thereby comes into the strongest opposition to that first need.

The reason why this opposition is not always noticed we already know. The religious
consciousness is absorbed in the individual mood and lets it unfold and express itself as
fully and wholly as possible. Only where memory, juxtaposition, and comparison become
possible is the opposition or the contradiction between the different ideas discovered. It
goes with religious experience on monotheism's ground as it goes with it on the ground of
the religion of drive and of polytheism. Religious faith's best---and often only---weapon
against criticism of every sort is to demand that all juxtaposition and comparison be kept
away. Again and again back to the whole and full moods, to the moments of religious affect!
He who for example---as has been demanded---reads the Bible only ``on his knees'' will
escape all criticism's difficulties. He is absorbed in the individual account, the individual
content---and even of that he involuntarily takes only what he can use. He will not come to
investigate whether the different accounts of Jesus's resurrection agree, and so forth. The
problem of how God's omnipotence is compatible with the human being's ``free'' will naturally
falls away when, at the moment one thinks of the one of these things, one does not think of
the other at all. When one---as Rasmus Nielsen and Henry Mansel attempted---distinguished
sharply between religion and theology, the possibility of this distinction lay in this
opposition between states where the great differences are forgotten and states where they are
juxtaposed and compared. But without an intellectual asceticism this opposition will not be
able to be carried through. There are at any rate those who will not lie on their knees but
stand upright and look freely about them in the world, and who do not find that the highest
they know suffers thereby. It is orientalism that would force us to our knees; but it is as
the heirs of the Greeks that we hold fast the conviction that an upright posture in life's
highest questions is possible.

\medskip

\noindent\textbf{46.} But in the opposition mentioned between the different states there may
lie the germ of tragic conflicts. Juxtaposition of the ideas that express the opposite
tendencies in their extreme form leads, in intellectually developed and yet passionate
natures, to the convulsive holding fast of the paradox. In quieter form and more unconsciously
it works in what I will call religious bashfulness. For religious bashfulness stems not merely
from a natural shyness of expressing oneself about the inmost lived experiences, but also from
a certain indefinite intimation that every image and every word in which one would express the
highest can always give only something finite and limited, and thereby not merely provokes
criticism from others, which is the least of the matter, but awakens in the individual himself
doubt about the validity and value of what has hitherto stood for him as the highest. It brings
with it for many natures a great spiritual pain to make up the final reckoning on this point,
especially where piety towards the transmitted forms works along to keep these outside a
juxtaposition and comparison that would otherwise lie near at hand.

By taking these relations into consideration we shall at once understand criticism's
unavoidability and its slow influence. Perhaps it is the eternal strife that we here stand
before. But it is a strife that must be carried on, and which---even if it should never
end---is not at every time fought on the same ground and with the same weapons. The strife's
eternity therefore does not exclude progress. The battle of the Einherjar was constantly the
same, but the battle of which there is here question is never the same; thereby it stands above
the battle in Valhalla. That the strife constantly repeats itself comes from the fact that a
change in the life of ideas takes place more quickly than a change in the life of feeling; a
theoretical criticism of religious ideas therefore does not reach down to the bottom of the
problem; the decisive point is the valuation of the feelings that lie at the foundation of the
religious ideas, and that after every critical cut are able to bring forth new shoots on the
stem of the religious ideas. A constant relation of oscillation and reciprocal action takes
place between the different poles of the life of the spirit.---The valuation of religious
feeling itself belongs in the ethical philosophy of religion.

\medskip

\noindent\textbf{47.} An advance will already have been made with the serious and consistent
acknowledgement of this constant reciprocal action, which makes a definitive solution
impossible. Let us try to take this experience up with us into the religious reckoning. The
idea of a finished divinity then becomes impossible. If, on the other hand, the religious
consciousness can reconcile itself with the idea of a divinity that is not finished and
concluded, but in constant becoming and development---as the religious consciousness itself
is, but on a larger scale---then there will be the possibility of a religion of hope, which
assumes that new tasks are constantly set, but also new possibilities for the treatment of the
tasks. Every state then comes to stand only as provisional, at most as a station on the
road---on existence's road in the large, on the individual's road in the small; and one will
then not be surprised or offended that ``the goblin moves with us,'' since the old problems are
posed in new shapes. The new dwelling to which ``the goblin'' follows us may very well be a
better dwelling, and why should ``the goblin'' absolutely be an uncanny being? He expresses a
law for our thought and perhaps a law for existence. When the divine unchangeableness consists
in, or makes itself known in, the fact that all change takes place according to definite laws,
and that the law of development itself is a fundamental law in existence, then the contradiction
between unchangeableness and changeableness falls away. The unchangeable is then the law of
change itself, and even if the individual law is altered, it happens always by virtue of a
higher law. Only a few steps can our thought climb up the graduated series whose possibility here
shows itself; but this very capacity to climb is sufficient to call forth a concluding hope, a
faith in the persistence of the valuable, a holding fast of the great \textit{Excelsior!} which
is the noblest element in all higher religions, only that the dogmatic formulas so often hinder
its full unfolding. If the conclusion we can win were more than a concluding hope, more than a
subjective conclusion, then there would be no constant \textit{Excelsior!} But when neither our
experience nor our thinking leads to an objective conclusion, the reason for this might indeed be
that existence itself is never concluded. That would be the most thorough explanation of the fact
that no dogma and no image suffices.

If one will say that this thought leads out beyond what can rightly be called religion, I shall
not quarrel about it, but let the word fall, especially if one will give me another in its place.
We are free and will not make ourselves slaves of the word.

\subsection{Religious Experience and Tradition}
\label{ssec:tradition}

\noindent\textbf{48.} The form and content of religious faith will never be able to be
understood from the individual's religious experience alone. The preceding investigation
already teaches that: the development from the religion of drive through polytheism to
monotheism takes place through long ages and many generations, and each individual stands
at a certain point within this development, determined by what goes before and determining
what follows after. Even if he has made the deepest and most independent experiences of the
relation between value and reality, the manner in which he expresses and explains these
experiences will be conditioned by the circle of ideas that stands at his disposal, and will
thus be more or less determined by tradition, whether he himself notices it or not. Ideas can
be involuntarily interwoven with what he lives through and feels, so that their content
itself seems to be immediately lived through. Only on closer investigation is the difference
discovered between the immediately given elements and the elements interwoven with them.

To illuminate this relation more closely I shall bring forward a series of historical examples.

\medskip

\noindent\textbf{49.} Among \emph{nature-peoples} already the relation between the
individual's lived experiences (especially in ecstatic states) and tradition appears clearly.
Among a Siberian tribe the young man who would become a shaman goes out into solitude, roams
by night in woods and mountains, and here sees visions in which the tribe's gods and the
spirits of the forefathers reveal themselves to him. He now sees what he has previously only
heard about. It goes likewise with the shaman when he is about to lose consciousness by
whirling himself round: in the ecstatic state thus produced he sees before him what he
otherwise knows only from tradition. A strict and long time of preparation must be gone
through by the young Central Australian before he can be initiated into the tribe's secret
faith.

What thus holds for the crudest stages of ecstasy and vision holds in principle at higher
stages too. The visions that the enraptured and inspired human being receives are determined
in detail by the ideas he has had beforehand. One can even produce hallucinations of a definite
kind by commanding an individual who has been brought into a hypnotic state to see a certain
object before him at a later moment, or by affecting him with definite sense-impressions
(colours, tones, tastes, and so on). Charcot and Pierre Janet have in these ways been able to
determine their patients' visions. To this comes the fact that the hallucinations that lie at
the foundation of the visions are often very indefinite and elementary, so that only by being
combined with and interpreted by ideas that the individual involuntarily draws from his own
memory do they acquire definiteness and significance. The individual gets, for example,
hallucinations of points of light or of a gleam of light and involuntarily forms from these a
white figure, which at last perhaps even speaks to him. An example of such a transition of the
vision from the indefinite to the definite and articulate is presented by the peasant girl
Bernadette's vision, which gave occasion to the worship of the Madonna at Lourdes. Still more
clearly this characteristic mark appears in one of Swedenborg's visions. One morning he sank
into religious brooding, and when he raised his gaze he saw above him in the sky a radiant
light; when he fixed his gaze on it, the light drew aside; heaven opened for him, and among much
else wonderful he saw the angels standing in conversation; kindled by a lively wish to know what
they said, this too was granted him, so that he first merely heard a sound (\textit{sonus}) that
expressed heavenly love, and thereafter formed speech (\textit{loquela}) that was full of
heavenly wisdom. Thus both as regards the hallucinations of sight and those of hearing a
progressive articulation takes place in the course of the vision's unfolding. And although it
was unutterable things the angels spoke of, things of which most cannot be expressed by human
speech, Swedenborg was nevertheless able by the help of his earlier visions to understand a good
deal of it: it is a kind of Unitarian theology that the angels develop. An orthodox theologian
would quite certainly have heard the angels utter other dogmatic views.---

Purely elementary hallucinations are far more frequent than one as a rule supposes; but only
under special conditions, with strong movement and tension in consciousness and under the
influence of dominant, feeling-toned ideas, do they become proper visions. The value of such
visions may be highly different, although the psychological explanation of them is essentially
the same. A hallucination is called a vision only when it has a certain value; but in a
psychological respect there is no difference of principle to be made. All visions bear the stamp
of the visionary personality's character, memories, and stage of culture. As a single example of
how ideas that were previously alive in consciousness can determine an elementary hallucination
and give it its significance, I shall (besides the vision of Swedenborg mentioned) adduce a
vision of Vincent de Paul which he had when his pious friend Mme de Chantal died: he saw a
luminous globe rise into the air and up in the heights unite with another globe of light,
whereupon both together were absorbed into a still more radiant globe of light, and an inner
voice now told him that it was the pious woman's soul which, in union with the soul of the
earlier departed François de Sales, became one with God!---It will therefore especially depend on
which ideas prevail beforehand of the visions; and if these ideas are formed and organized, they
will be able both to interpret and to soften strong visionary and ecstatic movements in the
religious domain. It is a phenomenon that repeats itself again and again in the history of the
religious life. Paul saw the heavenly Christ in the shape he knew beforehand, although he had
regarded it with other eyes.\textsuperscript{51}

When the ecstatic \emph{cult of Dionysus} made its way to Greece from Thrace, the oracle at
Delphi succeeded in organizing the wild movement by bringing it into connection with the existing
cult of Apollo. Whereas those seized by the Bacchic frenzy had previously known virtually no rule
other than the promptings of the swarming feeling, they were now brought under the influence of a
clearly formed and harmonious circle of images. The Delphic oracle generally, or in other words:
the Delphic priesthood, stood as a kind of authority that expounded and ordered the movements of
the religious life. Plato calls ``the god at Delphi'' the interpreter transmitted from the
fathers. Since the oracle was consulted by all the Greek states concerning the ordering of cult,
the Delphic tradition acquired a significant influence on the development of Greek
religion.\textsuperscript{52}

\medskip

\noindent\textbf{50.} In the history of \emph{the Jewish religion} the finding of the book of the
law (Deuteronomy) in the Temple under King Josiah (c.~620 B.C.) is an important event. A
thoroughgoing organization now stepped forward over against the free individual movements in the
domain of the life of cult and of prophetism. A closed priesthood now stood as leader of the
religious life, whereas previously the head of the household had attended to the cult for himself
and his own. The opposition between clergy and people was sharpened; prophetism was soon
succeeded by the priestly system. Judaism got its canonical book, a codified expression of the
tradition; it became a religion of the book, something in which Christianity and Mohammedanism
later imitated it.\textsuperscript{53}

At the coming-into-being of Christianity the religious traditions co-operated in the people and
the time in which Christianity's founder came forward. We are cut off from knowing the inner life
of Jesus of Nazareth prior to his public appearance. But however his life of feeling may have
unfolded, he used his people's traditions for its interpretation and formation. In the Sermon on
the Mount he presented his ethical teaching as an inwardization and spiritualization of the Mosaic
law. In the chiliastic expectations, the hopes of a coming kingdom of consummation perhaps
transmitted from Parsism to Judaism, he found (or his disciples found, under the impression of his
personality) forms and images for the great hope he implanted in the race. Thus even the most
significant religious personality that history knows developed by means of tradition and exercised
influence upon it. This does not exclude independence, for the old becomes new by being taken up
and applied by a deep and original genius.

In the same way as Jesus was placed when he was to form his religious ideas, his disciples were
placed when they were to interpret and express the effect that the Master's personality, his life
and death, had made upon them. They conceived him either as the expected Messiah \emph{(the three
first Gospels)}, or as the Logos, the eternal Word \emph{(the Gospel of John)}, or as the
sin-offering (Paul), or as a spiritual high priest \emph{(the Epistle to the Hebrews)}. These are
different roads along which they sought to make clear what the Master was for them.

\medskip

\noindent\textbf{51.} The picture we get through the apostolic epistles of the first Christian
congregation shows us no fixed organization. Enthusiasm and ecstasy stirred strongly and freely.
But if a congregational life was to develop, it was not enough that individuals gave themselves up
to inner lived experiences which in violent rapture made it impossible to form definite thoughts
and words, so that the outbursts became unintelligible to others. ``Speaking in tongues'' has
consisted in inarticulate and therefore unintelligible outbursts. ``The spirit'' worked so
violently that ``the understanding'' did not get to speak too. Paul therefore admonishes the
Corinthians to limit these ecstatic states, when they hold their gatherings, to such cases where an
interpreter can test and interpret them, so that what is lived through in these states may benefit
others. \emph{The fourteenth chapter of the first epistle to the Corinthians} gives an insight into
these matters (see also 2~Cor.\ V, 13). The interpreter must probably have been able to enter so
far into the state of the one ``speaking in tongues'' that he could find the thoughts and words
through which what moved in the mind in this state could be expressed. An immediate sympathy, an
involuntary inference from analogy guided by sure tact, formed the presupposition of the
interpretation, just as everywhere where one is to be guided by others' expressions and movements
to understand what goes on in their inner life. But at the same time, certainly, partly the Old
Testament writings, partly the apostolic teaching furnished norms for the exposition. Perhaps this
lies in the demand set up in the Epistle to the Romans (XII, 6), that the gift of prophecy is to be
exercised ``in agreement with the faith,'' if ``faith'' here denotes the content of faith. Thus the
individual's inner lived experiences were early brought under traditional norms.

Later on---as the Church became organized---the stream of individual enthusiasm was more and more
dammed in.

In the apostolic age it was especially regard for the edification of others that motivated the
limitation. But now everything depended on preserving agreement with tradition and excluding heresy.
The free individual movements were then pushed back as dangerous---in so far as they did not cease
from lassitude. Among lay people the capacity for visions and for speaking in tongues fell away
first---quite naturally---later among monks and priests. Comparison and misgiving began to prevail.
One did not rely on the immediate promptings, but tested them according to norms that became ever
more precise. In the so-called \emph{Pastoral Epistles} (the epistles to Timothy and Titus) ``sound
doctrine'' is already inculcated over against heretical errors. The history of the Church shows how
organization and the fixing of dogma advance in the course of the centuries. In place of expansion
came organization and articulation. The setting up of a canon \emph{(the collection of the New
Testament writings)} was the most important link in this course of development. It is the Church
that here, at a certain point of its development, once and for all established which books are to be
regarded as genuine testimony to the true doctrine.

Only through the ecclesiastical authorities could it from now on be decided whether a formation of
religious ideas was valid, whereas earlier the revelation was constantly continued in the inner life
of individuals. The Church now became the guardian of the gifts of grace. It went here as with the
miracles: only the Church is to decide whether a religious truth is present. Even as regards
religious conduct of life the Church demands that its tradition be followed: when Francis of Assisi
wished to renew the apostolic life (after the tenth chapter of the Gospel of Matthew), the Church
carried through the formation of an order after the model of the older orders (\textit{exempla
antiqvorum}). The Church here turned, as so often, against the renewal of the original when this
conflicted with the traditions that had formed in the meantime.

\emph{The Catholic Church} nevertheless maintains that new dogmas can be formed. When the religious
need has led to the formation of certain ideas, the Church's head can declare such a content of
ideas to be ecclesiastical doctrine. The new dogma then stands as the valid expression of something
the Church has fundamentally always believed, but of which it has only now become fully conscious.
Before the proclamation of the dogma of Mary's immaculate conception, Pius IX sent out a circular in
order to learn whether the congregations longed for this dogma, and the result of this approach was
such that an inscription could afterwards be set up in St~Peter's declaring that the Pope by the
proclamation of the dogma had fulfilled the longings of the whole Catholic world (\textit{totius
orbis catholici desideria explevit}).\textsuperscript{54} Catholic writers sometimes regard the whole
ecclesiastical fixing of dogma---right down to the dogmas proclaimed in the nineteenth century---as
further unfoldings of what Jesus in the forty days between the resurrection and the ascension spoke
of with his disciples.\textsuperscript{55}---The history of the dogma mentioned affords an
interesting contribution to the psychology of religion by the clear manner in which the relation
between need of feeling (\textit{desideria}) and idea there appears. One felt the need for an
intermediate link between the Jesus exalted to God and humanity. And besides, there was felt a need
to have the element of womanliness, and especially of motherly love, represented in the religious
circle of ideas. In an extraordinarily beautiful and instructive manner \emph{Yrjö Hirn (Det heliga
Skrinet,} 1909, chs.~11--12) has depicted the development of the worship of the Madonna and thereby
given an example of how a religious need leads to cult and through this in the course of time to a
dogma.

But as in every psychological process, manifold factors here too work together. No dogma can be
explained out of a direct religious need alone. The Church's inner and outer circumstances in
particular play an essential role. The Church feels itself bound by its traditions (as it itself
understands these), whether it follows them in blind habit or seeks to form new ideas in analogy
with, or as a supplement to, the ideas already formed. It combats a ``false doctrine'' and therefore
strives after a formulation that rejects it as sharply as possible. Or it seeks to unite and
reconcile contending tendencies by a shrewd compromise, often a shrewd indefiniteness or ambiguity.
Or it strives to form ideas that can correspond to its cult, as this has once for all developed.
Harnack finds in his history of dogma that it was eleven (11) different factors---namely ten besides
the properly religious need---that co-operated in the development of the ecclesiastical dogmas (the
content of the present orthodox dogmatics) in the first centuries.\textsuperscript{56} To these many
elements there ought perhaps to be added an aesthetic element. It is according to \emph{Yrjö Hirn}
not merely religion that works upon poetry, but pictures and poetic works can often only be
understood through ``the great, anonymous, and collective poetry that one can read out of the whole
ecclesiastical system of readings.''\textsuperscript{57}

\medskip

\noindent\textbf{52.} What thus appears in the history of a religious community shows itself
still more clearly in the religious life of a single personality, when this personality has
the capacity and the need to describe its inner development. A great example of this we have
in Augustine's \emph{Confessiones.} When Augustine composed this work he stood on the firm
ground of ecclesiastical tradition. He believed, by his own statement, in the Gospel only
because the Church transmitted it. But the central religious process---a feeling that strives
after expression in a content of ideas---can nevertheless be traced in him. The tenth book of
the \emph{Confessiones} in particular is characteristic in this respect. Augustine here seeks
to become conscious of \emph{what} it is that is the object of his inner need. ``What is it
that I love, when I love thee?'' he asks his God. He goes through a long series of beings and
forms that experience presents; but none of them satisfies what his feeling yearns towards.
Only in his own inner life, in his own spirit, does he find an analogy he can use. In a
profound development, extraordinarily interesting for the history of the fundamental
psychological concepts, he points out memory as the distinctive characteristic of the life of
the spirit. God must be a being who works after the likeness of the force that works in our
memory. And he works indeed also in our inmost soul---far further within me than the inmost
that I myself know! In the inner man truth dwells. But (now there is a passing over from the
psychological to the ecclesiastical consideration) nevertheless: God is not merely something
thus inwardly operative; he is first and foremost himself the object of a memory---the memory
of the ecclesiastical preaching (\textit{ministerium prædicationis}), from which fundamentally
we know him alone!\textsuperscript{58}

Never have deep personal experience, energetic thought, and ruthless faith in authority entered
so close and characteristic a combination as in this greatest teacher of the Church. In him too,
indeed, the opposition between self-absorption and dependence on outer tradition appears clearly.
Go not outside yourself, for truth dwells in your inner life! he exclaims. But it is the Church's
tradition that he thinks he owes everything to. As if that did not also come from outside! Yet
the strong dependence on tradition is unable to veil the fundamental psychological relation
between feeling and idea that is distinctive of the religious domain. Especially in the utterance
that was above set as the motto for the first section of the psychology of religion does this
fundamental relation appear definitely: a predicate is sought that can serve to determine what
makes itself felt in the feeling. If doubt during the spiritual struggles of his youth had not
led him to cling to the Church's authority, and if he had not later, when he stood as a churchman,
felt so strong a need to assert this authority against ``dissolving'' elements, the movements of
his inner life could have led to freer and more universal results. To this came the fact that he
lived in a time when the world---by his own statement---had grown old (\textit{senuisse jam
mundum});---the boldness for a free and independent shaping of the inner life does not thrive in
such a time.---

A certain analogy to the development in the tenth book of Augustine's \emph{Confessiones} is
presented by the line of thought in the \emph{Upanishads.} Just as Augustine asks: what is it that
I love, when I love thee?---so it is asked in an Upanishad: what is Bráhman? And the answer is:
Bráhman is Atman (breath, spirit, soul)! Thus the principle that bears the whole of existence, and
which itself (if Bráhman means yearning or willing) is a projection of feeling's need---this finds
(at any rate provisionally) its best expression in ideas drawn from the life of the soul
itself.\textsuperscript{59} The old Indian thinkers stand in a freer relation to tradition than
Augustine was able to; therefore the psychological process appears more clearly in them.

\medskip

\noindent\textbf{53.} In the mystics of the Middle Ages we find interesting examples of the
reciprocal action between feeling, the formation of ideas, and tradition.

According to a terminology frequent among the mystics, ``the will'' or love is placed in opposition
to memory and understanding. In the highest states all these faculties slumber, since the will is
overwhelmed by overflowing blessedness; all suffering is forgotten, and no search after causes of
the state, no comparison with other states, becomes possible. Not even any vision is possible in
the highest, ecstatic moments. And in such a state the human being cannot communicate himself to
others either. ``The will'' remains in this state longer than memory and understanding, and as soon
as these can work, the formation of ideas and communication become possible. The mystic can then
have written down what he thinks he has experienced. But the difference between what is really
lived through and the thoughts and words in which it is expressed is very strongly emphasized by
the mystics.---``It is felt, but not expressed,'' says \textsc{Hugh of St.\ Victor}, and using the
images of the ``Song of Songs'' he says: ``Thy beloved comes invisibly; he comes hiddenly; he comes
incomprehensibly; he comes to embrace thee, not to be seen by thee!''---\textsc{Angela di Foligno}
speaks of a vision in the following manner: ``What did I see? God himself! I can say nothing else
about it. It was a fullness, an inner, filling light, against which neither words nor comparisons
avail. It was the highest beauty, that which makes all lips close, and which contains the highest
good!'' When she read what had been written down about what she had lived through at her dictation,
she often could not recognize it again.\textsuperscript{60}

But although the mystics sharply emphasize the difference between lived experience and description,
they nevertheless maintain no less strongly that what they live through agrees with the teaching of
the Bible and the Church. It is a psychological riddle how they can convince themselves of the
agreement between the content of their lived experiences and the Church's teaching, when there is so
great a difference between lived experience and description; for it is only by the help of a
description that the comparison with the Church's teaching can take place. It also appears clearly
enough that the agreement with the Church's teaching is expressly willed, not merely found. In the
interpretation and description of the lived experiences one watched attentively that nothing came in
that conflicted with the teaching of the Bible and the Church. There are express declarations about
this from some of the greatest medieval mystics. \textsc{Suso} thus says of himself: ``In the working
out of his chief work [the Book of Eternal Wisdom] he was not in a state like that of one who is
active and dictates, but like that of one who is seized by the divine \emph{(divina patiens)}. When
what was inspired in him had been formed in writing, he came to himself \emph{(ad se revertens)} and
investigated with care \emph{(diligenter rimatus est)} whether there was anything found in it that
deviated from the utterances of the holy fathers.'' The difference between rapture, self-collection,
and anxious submission to authority appears clearly here. And self-collection here works not, as in
Augustine, to find by the road of one's own reflection an expression for the lived experience, but is
used predominantly to see to the agreement with the Church's teaching.---In the female mystics
obedience to the Church's authority appears still more definitely than in the male mystics. The
apostle Paul had indeed forbidden women to speak in the congregation and referred them to enquiring
of the men in religious matters. And since they lacked the male mystics' scholastic education, they
had to feel themselves easily uncertain. However sure they were that their highest lived experiences
had a divine origin, they nevertheless acknowledge throughout the Church's authority as decisive.
Angela di Foligno's visions and lived experiences were written down by a monk who declared that he had
added nothing, but had left out ``many things that were too exalted to be contained in his wretched
understanding''; thereafter the book was repeatedly examined by learned monks. And her confessor even
found authority for her strong erotic expressions in the description of the ecstasies: ``the Song of
Songs belongs to the Bible too!'' (This last trait I owe to a young Swedish friend who has
particularly occupied himself with the subject).---\textsc{Saint Teresa} describes how, in order to
avoid going wrong, she questions everyone who could instruct her, and she clings---she says---so
firmly to her Credo that even if she received all possible revelations, indeed even if she saw heaven
open, her faith would not waver concerning the least point that the Church teaches! And if the thought
occurs to her for a moment that what she hears might yet be just as true as what the fathers had
heard, then she is sure that it is the devil tempting her! The decision whether her book agreed with
the Church's teaching she laid in the hands of a learned and pious Dominican monk.\textsuperscript{61}

There is here a reciprocal relation present. For just as the lived experiences are interpreted and
tested by the help of tradition, so they are also often called forth by living into the transmitted
ideas. It is then reading that sets the mystical absorption and lived experience going. Both Hugh and
Teresa therefore also mention the significance of reading. One is reminded here how the persons in the
van Eycks' and Memling's religious pictures are frequently so absorbed in their books that they forget
everything and everyone around them.\textsuperscript{62}

And yet the mystics had also a purely subjective criterion of the validity of their lived experiences:
the inner peace and rest in the mind that was produced. So long as the need for submission to Bible and
Church was strong and living, this inner peace and rest could indeed only be reached when one had
brought one's ideas into agreement with the transmitted teaching. But there was nevertheless a
possibility that the immediate trust in the experiences of the inner life, and the intellectual activity
in the work of shaping them into thoughts, could become so vigorous that submission to tradition was no
longer possible, but the old leather bottles were burst by the new wine. In this way the way can be
paved for great religious conflicts, and this way all founders of religions and reformers have gone,
often after an honest attempt to find peace within the old forms. It is then discovered that a new
living garment must be woven for the divinity. In a Protestant mystic of our own time (discussed in my
book ``Oplevelse og Tydning,'' especially pp.~96--100) the opposition between lived experiences and
tradition appeared---in contrast to Angela and Teresa---very sharply and decisively.

\medskip

\noindent\textbf{54.} In \emph{the Reformers} too one notices the psychological process in which the
immediate experiences of the personal life meet, are shaped, and are limited by ecclesiastical
traditions. In Luther's case we have already seen how he---in a manner recalling the line of thought in
Augustine---seeks an expression for what moves in him. God is, he says indeed, that in which you have
the highest trust. But the closer development of his concept of God Luther wins only by simple adoption
of the old ecclesiastical doctrinal determinations. He let the scholastic doctrine of the Trinity stand
as it stood, while he set his personal stamp on the religious ideas that stood in immediate connection
with his personal lived experiences, thus justification by faith and the sacrament of the Lord's Supper
in a form closely akin to the Catholic conception. That Luther thus took up other and more than
corresponded to personal lived experience was a half-measure that finds its explanation in his faith in
the Church, but which became fateful for Protestantism's later development. A clear principle that can
be consistently carried through Protestantism gets only when it supports itself on the free conscience
and the personal experiences that individuals make concerning the relation between value and reality.

Still more clearly than in Luther the psychological foundation of the religious assumptions appears in
\textsc{Zwingli}. While Luther halts his independent train of thought when he has defined God as the
object of the highest trust, Zwingli goes a step further and thereby comes to his absolute doctrine of
predestination. Zwingli stresses more strongly than Luther the unconditional certainty of faith that the
individual can have of attaining salvation, and to this subjective but unshakable certainty corresponds
the faith in divine election. And if this election is to reach its goal, it must stem from a power that
makes all resistance give way, an absolute, infinitely exalted power. Even if Zwingli in carrying through
this line of thought is also under the influence of his philosophical studies (of Plato and the Stoics),
it is nevertheless the personal experience that gives the impulse by which the line of thought gathers
speed. Only an experience like Zwingli's could introduce a line of thought of this
kind.\textsuperscript{63}

Within Protestantism attempts have been made, still more definitely than the Reformers made them, to
trace religion back to the mind's own experience, to let all religion be a cult of the heart. Rousseau and
Schleiermacher are the most important representatives of this endeavour. But they too, when they attempt
to shape their lived experiences into thoughts and words, are referred to ideas they take up from the
spiritual atmosphere in which they live, and these ideas they are inclined to regard as just as
immediately given as the lived experiences themselves. Rousseau ``felt'' immediately ``the truths of
natural religion,'' and although Schleiermacher proceeds with greater psychological sense and critical
circumspection, the obvious confusion of feeling's expression and its cause is nevertheless found in him
too.

Personal experience for the most part stands immediately and naively over against the transmitted ideas
and forms in which the religious need of earlier times found expression, and it thinks it takes them as
they are. Only a closer consideration discovers the constant displacements that here take place. The mere
fact that the emphasis is laid on other elements in the content of ideas than before brings with it a
nuance. Something similar happens here as when in the history of art a type of picture is taken up by a
new race of artists. Of the older Italian painters' relation to the transmitted Byzantine Madonna type,
\textsc{Julius Lange} remarks that they did not place themselves revolutionarily towards it, but little by
little carried it over into a wholly different and opposite direction: ``Into the fundamental features of
the transmitted lifeless mask they begin to insert their personal opinions about the gracious beauty that
was to be characteristic of God's Mother. Thus Guido of Siena. It is still the almost circular contour of
the crown of the head, the face's elongated oval\ldots. But how fresh, how lively and amiable does this
Madonna not look out of her eyes!''\textsuperscript{64} In the history of religion the displacements are
not to be pointed out in so vivid a manner. They take place partly by one's letting those elements in the
tradition come into the foreground to which feeling particularly finds attachment, partly by certain
elements' being involuntarily reinterpreted, understood poetically, symbolically, or rationalistically, or
else wholly omitted, without one's being provisionally conscious of thereby deviating from the tradition's
original meaning. It was a great relief to Augustine when in his youth he heard that the bodily expressions
the Bible uses of God could be understood figuratively, not purely literally, and he later applied the
allegorical interpretation confidently on a large scale when there was something in the ecclesiastical
teaching from which he could not otherwise wring religious value. That God's spirit at the creation floated
over the waters means according to Augustine that the human being by the gift of the Spirit, which is love
of God, raises himself above the bodily; that God ``created [both] heaven and earth'' has for him the
symbolic meaning that within the Church there is both a more spiritual and a more literal understanding.
\emph{The mystics} went still further along this road, in that they, as has been aptly said of them,
``wished to abolish every barrier between themselves and Scripture's content, to know everything [that
Scripture relates] as an eternal present.'' Meister Eckhart for example interprets in a sermon the idea of
the raising of the widow's son at Nain thus: the widow is the soul, and the dead son is the reason.---In a
similar track the edifying religious discourse moves more or less at all times, led by the endeavour to find
in accounts of what happened millennia ago immediate nourishment for the spiritual life of the present. A
double meaning is then ascribed to the old utterances, without one's as a rule forming a clear idea of the
relation between the two meanings. Even in philosophers of religion in the nineteenth century---in men like
Schleiermacher, Hegel, and Coleridge---we find the assurance that the conversion of the content of
ecclesiastical tradition into the forms that agreed with their own life of feeling and thought does not
alter this content's original meaning. Even \textsc{Søren Kierkegaard}, who wished precisely to deliver the
dogmatic content from the speculative evaporations, nevertheless comes to push the objective content of
dogma back; he appropriates at any rate only what his passionately strained life of feeling consumes, and in
such a meaning (often a reinterpretation) as he can use it in.\textsuperscript{65}

Where such an accommodation is no longer possible on account of the inner life's energy, the old forms are
burst, and if the religious life is to be continued, new forms must be produced. We then stand---as already
remarked---at a breakthrough where something new comes into being. The question then becomes how the
psychology of religion is placed towards such breakthroughs and towards the prophetic personalities in whom
the new comes into being.

But before I enter upon this question, there is still a general consideration to be made concerning the
relation between feeling and idea in the religious domain.

\medskip

\noindent\textbf{55.} Religious judgments differ---in a manner similar to aesthetic and
ethical judgments---from other judgments in that the element of feeling, which is never
quite absent from any formation of a judgment, plays a quite distinctive role in their
coming-into-being. Religious judgments are value-judgments. They express, on the
hypothesis I lay at the foundation, the human being's experiences of the relation between
value and reality, a relation that can itself acquire its immediate value. In their
simplest form religious judgments appear as exclamations in which an inner state of wonder,
love, hope, or fear finds vent. In more articulate form the religious judgments appear when
consciousness seeks clarity about \emph{what} it really experiences---what it admires,
loves, hopes, or fears. What is the significance, it is asked for example, of this wonderful
event? Who is this human being who intervenes so deeply in our life? Now the more conscious
formation of the judgment takes place. Every judgment is a connection of concepts that has
come about by attention's starting from an idea that is given in more indefinite form and
seeking a closer determination of it. The starting idea becomes the judgment's subject and
finds its closer determination through concluding ideas, which become the judgment's
predicate. The psychological process in the judgment's formation takes place through the
transition from the first of these ideas to the second. As regards the closer description of
this process I must here refer to \emph{``Den menneskelige Tanke''} pp.~71--99.

With value-judgments the starting idea stands in greater obscurity than with other judgments,
because the element of feeling has so strong a preponderance. Indeed, when the underlying
feeling is new and original, or when for other reasons it stirs with particular strength,
virtually all clear ideas may be lacking. Every strong emotion is indeed characterized by
ideas' disappearing more and more from consciousness. In what precedes, cases have been
brought forward in which all articulation and analysis, and thereby all formation of
judgment, was impossible; examples of this were the primitive Christian ``speaking in
tongues'' and the mystical ecstasy. Since the formation of a proper judgment is possible only
when there is a starting idea, however obscure it may be, it must be assumed that such
religious judgments as express what has been lived through in strong emotion are formed only
in later moments, when the movement has faded, so that the human being can ``return to
himself'' and in memory characterize what has been lived through by juxtaposing it with other
lived experiences and by thinking over its significance for the life of the personality.

When great religious personalities have called the object of their highest trust and love
\emph{God,} this is understood when by God one means the principle of the persistence of
value in reality. What bears and gathers up in itself all values, and stands as the source
and consummation of all values, must be the object of the deepest feeling. And it is this
that the religious judgments, breaking forth with involuntariness and originality, say: what
I here experience is only a single expression, a single testimony, of the power that bears
all value in the world of reality. It is then this power itself that I here experience!---The
concept of God is \emph{the fundamental predicate} in all religious judgments, the predicate
that religious thought seeks after, setting out from the experiences that have been made
concerning the relation between value and reality. The conclusion of religious thought would
be a complete determination of this fundamental predicate, such that the relation between
value and reality became quite clear and transparent. This it would especially become if the
principle of the persistence of value and the principle of existence's law-governed connection
could be brought to coincide in one principle. In a form satisfactory to scientific thought
such a conclusion is certainly impossible. But the religious problem persists as long as a
striving to hold fast the religious predicate-concept, and thus as long as a religious
questioning, is psychologically possible.

The history of religion and of philosophy shows that a closer development of the religious
predicate-concept (or, to use a Kantian expression, the religious category) most frequently
takes place by the help of tradition's images and forms. Only where the individual has the
capacity and the courage for independent formation of images and for speculation does the
development strike into wholly new paths, and that will especially happen where a breakthrough
takes place in religious development as a whole.

But it is of great psychological importance to hold fast that the predicate-concept or the
category is not the first thing, the originally given, but precisely that which is striven
towards. Such a determination of the religious category (``God'') is sought that it can be
understood what is meant when the individual lived experience is subsumed under it, seen as a
single example of its validity, or---by poetic personification---as an action of God.
\emph{The primary religious judgments have ``God'' as predicate, not as subject.} The concept
of God, the religious category, is subject to the same rule as all concepts and categories:
our concepts do service as predicates before they can be used as subjects. The simplest
judgments are predicate-judgments (judgments without a subject). But there is both during the
formation of the judgment and in the exact formulation of the finished judgment such a
reciprocal relation between subject and predicate that not only is the subject determined by
the predicate it receives, but the predicate is also determined by being referred to this
subject. To take an example: when I learn that Amphioxus is a vertebrate, I not only learn
something about Amphioxus, but I also learn something about the concept vertebrate, namely for
example that its content does not have among its constituents the possession of a brain. Thus
the religious category (the concept of God) also gets its closer determination through its
applications, that is to say, through the different occasions that lead to setting up this
concept. But will it ever become finished, so that one can operate with it as with a clear and
distinct concept? That is the great question.

Orthodox and speculative dogmatics have supposed that in the religious predicates they had a key
not merely to the understanding of religion's essence, but of existence as a whole, and have
attempted to construct a ``higher'' science by the help of these unconcluded and incompletable
ideas. But from the other side too the nature of the religious predicates has been misjudged. A
school in the history of religion, which for a time was dominant in the science of religion,
supposed that in the purely etymological investigation of the religious designations, especially
the names of gods, it had the sufficient method for coming to know religion's essence. Etymology
was then to be the true philosophy of religion. By this road one can also obtain interesting
information. Thus it is of interest that the very word ``God'' seems to indicate that the
religious fundamental predicate is formed under the influence of the act of cult and means the
one towards whom this act is directed. But this example also shows precisely the method's
insufficiency; for the religious fundamental predicate is in fact used far beyond the individual
occasion that the act of cult gives. The concept of cult must at any rate be considerably
extended if it is to become intelligible what lies in the expression mentioned. Religious
development is far more complicated, manifold, and deep-going than the purely linguistic, and the
relation between the inner states and their expression is in the religious domain still more
composite than the relation otherwise is between psychical phenomena and their linguistic
expression. It has therefore also shown itself that inner and outer kinship can be pointed out
between ideas of gods which among different peoples are expressed in words of wholly different
stems, while conversely names of gods that one was inclined purely linguistically to put together
denote gods of wholly different kinds.\textsuperscript{66} The linguistic designation is only a
support, which is acquired at a certain point of religious development, but which can thereafter
be attached to ideas that remove themselves far from the stage of development at which the
religious life of ideas stood when the expression was formed. So long as a religion persists,
there are constantly released anew processes that seek a provisional conclusion in the religious
predicate-concept, and this can therefore constantly be determined in new ways.

Is there not here something that the human spirit has at all times, more or less inwardly and
energetically, wished to seek, but for which it has never found and can never find expression and
form? Does there not hide here a riddle, a secret, on which the human spirit constantly broods,
and which presents itself in ever new ways, according to the alterations of experience and the
differences of personalities?

\begin{center}---\end{center}

\subsection{The Scientific Conclusion of the Psychology of Religion}
\label{ssec:conclusion-psych}

\noindent\textbf{56.} We have in what precedes come to a point where the psychological
method is put to a decisive test. So long as the arising of religious feeling, and the
ideas in which it finds expression, are clearly and distinctly determined by the historical
circumstances, the psychological method works as well as scientific method in the domain of
the humanities is on the whole able to. But where something qualitatively new comes into
being---where a deepening or widening of feeling in principle takes place, or where feeling
appears with a wholly new nuance---and where new ideas are formed to be feeling's
expression---is it there possible to give an exhaustive psychological explanation?---This
question presents itself particularly when we stand before the breakthroughs mentioned in
what precedes, where the constant reciprocal action between personal experience and tradition
does not seem to give the sufficient explanation. The prophetic personalities, who make new
experiences and bring forth new symbols, do not seem to be explicable by application of the
psychological method, which functions tolerably so long as that reciprocal relation holds
good. Such personalities seem to go off at a tangent instead of moving in a closed curve
about a given centre: they lead us into wholly new regions in the spiritual world.

This point has so much the greater interest because that tendency in the most recent theology
which stands nearest to the scientific way of regarding things lays strong weight on the view
that all revelation takes place in the inner life of the prophetic personalities and consists
in the life and the originality with which their religious experience appears, and which lets
them find thoughts and symbols that can bring clarity and peace in wide circles. While the
older theology in a rather external manner placed faith and faith's object in opposition to one
another, there is now a tendency to see in the very coming-into-being of faith in the human
soul the greatest miracle, the real revelation, or at least to emphasize that it too belongs
to revelation. Theology is here brought into as close a connection with psychology as it has
not been since Schleiermacher's time. But so much the more pressingly, then, does the question
present itself of the possibility of a scientific conclusion of the psychology of religion.

Before I myself attempt to state psychology's position on the problem indicated, I shall discuss
earlier psychologists of religion's position on this point.

\medskip

\noindent\textbf{57.} In his \emph{Esquisse d'une Philosophie de la religion d'après la
psychologie et l'histoire}, \textsc{Auguste Sabatier} seeks, as the work's title indicates, to
place himself at a purely scientific standpoint. He will use the general psychological and
historical methods, and his standpoint as well as his results are therefore in several respects
akin to the fundamental thoughts in the presentation I have given in what precedes. Dogma is
for him never the original and fundamental thing in religion. It arises comparatively late in
religion. The prophets always go before the rabbis. Dogmas are, like all ecclesiastical customs
and forms, derived in relation to religious feeling, and it is the task of the psychology of
religion to show how this derivation takes place. Thus first and foremost the concept of God is
a symbol in which religious feeling has found expression.

In carrying through the psychological-historical method, Sabatier holds fast that it is no
explanation to appeal to God's intervention when a single definite phenomenon is to be
understood. ``Since God,'' he says, ``is the ultimate cause of everything, he is not the
scientific explanation of anything whatever.'' It was the same thought that already at the
founding of modern science came to determine its relation to the concept of God (see §5). It is
just as justified and just as necessary in psychology as in physics.

But Sabatier nevertheless applies, when he is in difficulties, the theological explanation in the
psychology of religion. It is, according to his teaching, God who produces the religious feeling,
which feeling then in turn sets the life of ideas in motion and thereby produces symbols and
dogmas. ``God enters,'' he says, ``into reciprocal relation and contact (\textit{en commerce et
en contact}) with a human soul and lets it make a certain religious experience, out of which dogma
then proceeds by reflection. What constitutes revelation and is to be the norm for our life is the
creative and fruitful religious experience that was first made in the souls of the prophets, of
Jesus, and of the apostles.'' Here I now cannot see anything but that the French philosopher of
religion comes into conflict with his own scientific principle: he here explains a special
psychological phenomenon by reference to God, although such a reference according to his own
definite statement does not explain any special phenomenon. And when the concept of God itself is
a symbolic idea formed on the basis of the experiences of the life of feeling, then the
``explanation'' of religious feeling that is reached by its help can only be a symbolic
explanation, and the gap in the psychology of religion (if there is a gap) is in reality not
filled. Every dogma presupposes a religious feeling, and no dogma can therefore be used to explain
this.

If one assumes that there is a gap in the series of psychological (and the corresponding
physiological) phenomena which can only be filled by the use of theological concepts, then one
has the obligation to state definitely where this gap lies. Where is the point at which the
supernatural intervention is to take place? All development of the soul shows itself indeed, the
better we come to know it, to be so connected and to rest on such an interplay of elements that
it is just as difficult to find a gap where an influence could intervene as it is to follow all
the transitions and unravel the individual threads. The assumption of a break in continuity brings
with it no fewer difficulties than the assumption of a complete continuity. As a rule those who
maintain an interruption of continuity do not pose the task sharply enough. It is only seldom that
one institutes so energetic a hunt for the point at which the interruption could lie as the older
Catholic theologians did when they---with full consistency, but with little taste---urgently
discussed whether the supernatural act that effected the Virgin Mary's freedom from sin was
inserted at her very conception or only at some point during the development of the foetus, and if
so at which. A similar investigation is incumbent also on a standpoint like Sabatier's and other
modern theological philosophers of religion's. It has not been made, and it would undoubtedly be
hopeless; but at the standpoint named it is a necessity.

Theological philosophy of religion involves itself in an epistemological circle similar to
materialism's. The concept of matter is formed by thought on the basis of sense-sensations, but
materialism then uses this concept dogmatically to explain the coming-into-being both of
sense-sensations and of thought. Just as sense-sensation and the logical principles stand as the
presuppositions that bear our theoretical conception of existence and therefore cannot in turn be
grounded by it, so religious feeling stands psychologically as a presupposition that cannot be
``explained'' by \emph{the} ideas it itself leads one to form.---Perhaps it is due to this analogy
between theology and materialism that they so often mutually understand one another better than
either of them understands critical philosophy.

\medskip

\noindent\textbf{58.} Half a century already before Sabatier's book appeared, \textsc{Ludwig
Feuerbach} had maintained that the psychological method is completely sufficient in the religious
domain---that religion is nothing but a psychological product---that all theology is psychology.
The psychological philosophy of religion owes Feuerbach an extraordinary amount, and his brilliant
analyses and characterizations stand in many respects still unsurpassed. But with respect to the
method and to the critical limitation of the results there are various misgivings to be
urged.\textsuperscript{67}

Feuerbach is fond of short and striking formulas, which hangs together with the fact that in him
the psychologist must often give place to the agitator. When he has found an essential factor in
religion's psychology, he is inclined to make it the be-all and end-all. Thus when he lets the wish
appear as the theogonic principle, in that it ``out of itself and only out of itself'' produces the
idea of God. Feuerbach here takes no account of the complicated circumstances under which the
formation of religious ideas takes place. However great the influence that must be ascribed to
feeling, it is nevertheless in the religious domain not merely active in relation to knowledge, but
also dependent on it. A significant side of religious development consists precisely in knowledge's
quiet influence on feeling. There is here---as I have also in what precedes sought to show---a
constant reciprocal relation present, even though the preponderance lies on the side of the
elements of feeling.

Without such a reciprocal relation religion would be able to have no practical significance for the
human being. As it is, on the other hand, myth and legend can become forms in which the conception
of nature and the memory of great personalities exercise their influence on the life of the race,
and in the symbolic language of dogmas the content of the weightiest and most important experiences
of life can in condensed and concentrated form benefit later generations.

Even with these modifications the proposition that all theology is psychology will not be able to
be the object of strict proof, any more than, in the domain of outer nature, the proposition that
all material phenomena are due to material causes. But the psychological-historical method
nevertheless stands as the only one we can apply, the only one that---where it can in the
particular case be applied---can give us a real explanation. Where it cannot be applied, we can get
no explanation. What we cannot explain by this method stands as a mere fact that can be described
but not explained---whereby of course it must be well attended to that one does not smuggle an
explanation into the ``description'' itself. To assert that all religious phenomena are subject to
psychological laws would be just as much dogmatism as to assert that there are phenomena that
cannot be explained psychologically at all. Psychology can give us working hypotheses, and we do
not give up these hypotheses until better working methods have been pointed out to us---but neither
do we confuse them with proved truths. In this spirit I have sought to work in the preceding
investigation of the development of religious ideas. We are on this conception compelled at many
points to concede that we stand before ``mere facts.'' But we add: every ``mere fact'' is (when it
is rightly conceived and described) a problem, so long as it is not placed in close and law-governed
connection with our other experiences.

\medskip

\noindent\textbf{59.} As regards now the problem of new formations in the religious domain, as these
especially appear in the prophetic personalities, it is not an altogether unique problem, but only a
special form of a problem that appears in different degrees and nuances in all domains of
experience. The more individual, and the more qualitatively new, a phenomenon is, the more difficult
it is to bring it into that continuity with other phenomena which is the condition of scientific
understanding. It is the opposition between continuity and quality or individuality that in all
domains sets knowledge its tasks. The religious problem is therefore not any altogether distinctive
problem, but presents a definite analogy with thought's other problems.

The concept of personality sets inquiry the double task: to find the continuity within the
distinctive little world that each single personality is (for without inner continuity, no
personality)---and to find the continuity between the individual personality in its distinctiveness
and the rest of existence. It is a task similar to that which natural science has towards the
qualitative distinctive characters of forces and substances; but the concept of personality
undeniably brings with it far greater difficulties for the humanities than the concept of quality
brings for natural science. Personality is itself the most pronounced quality that is known. The
humanities are therefore more difficultly placed towards their problems than natural science towards
its.

The problem becomes particularly sharpened when it is a matter of such personalities in whom in a
distinctive sense something new arises---the men of genius in the word's widest sense---those in whom
a distinctive genius rules, that is to say, in whom the involuntary and half-unconscious life of the
spirit produces effects that surpass what clear consciousness and sober work would have been able to
accomplish. The prophetic personalities belong to this kind by reason of the life and the originality
with which their religious experience appears. It is life's problem in a special form that is here
posed. Whether it is a question of Buddha's, Socrates's, or Jesus's personality, it is one and the
same great problem that we stand before.

But because the problem is sharpened, it is not thereby said that it must be treated by wholly other
methods than we otherwise apply. It would then at any rate have to be demanded that a new method be
stated---indeed, a wholly new epistemology. Merely to lose oneself in wonder does not help, and
besides, the astonishment and the admiration may be quite as great in him who quietly works to find as
much psychological and historical connection as he is able in the world of personality, as in the
theologizing Romantic who thinks he must at this point assume a wholly different principle of
knowledge from the one he otherwise applies. That in the domain of the humanities we more often than
in that of natural science experience our knowledge's limitation is something that the study of
religion is not alone in teaching us. It is not merely our imperfection, but it is first and foremost
existence's fullness, the revelation of existence's inner richness, to which this limitation is due.
It is not merely barriers to our inquiry, but material and task for it, that we here find. Where
something new breaks forth, it is a testimony that there were more hidden forces present than
experience previously seemed to teach; and even if we do not succeed in fitting this new thing into
such a continuous series as it is our thought's ideal to form, this testimony nevertheless keeps its
own value---and keeps it best when one gives it no dogmatic label. It will be seen that the psychology
of religion here in a certain sense takes the concept of ``revelation'' more seriously than orthodox
theology, which always confuses its dogmatic concepts with what is immediately and really given. So
long as existence itself is placed in becoming (and that it is must be inferred from the fact that we
find neither in the world of the spirit nor in that of nature anything altogether unchangeable and at
rest)---so long must something new come forth which---at least provisionally---stands as a kind of
paradox. Thereby nothing is altered in principle in our knowledge's position towards existence, any
more than in our knowledge's method or in our personal life's relation to our knowledge. Who has said
that we should be able to explain everything, or that it would be impossible to live if we do not get
everything explained?---

To seek an explanation in the unconscious life of the soul does not help. The unconscious is, like God
or the soul, a very indefinite concept. Psychology assumes potential life of the soul with the same
right as physics potential energy. But in both places one must build on the work that lies before us,
and by which the words soul or energy first acquire a meaning. Besides, in the unconscious life of the
soul there lies the possibility both of the lowest and of the highest in the life of the soul as it
actually lies before us, and special causes must therefore be sought in each individual case. It was
precisely this last point that made \emph{William James}, who himself laid so great weight on the
unconscious life of the soul (the back door, as he called it), place himself definitely in opposition
to \emph{Myers's} attempt to found all psychology on that concept.\textsuperscript{68}

\begin{center}---\end{center}

\section{Dogmas and Symbols}
\label{sec:dogmas}

\begin{flushright}
\textit{Du kerkerst den Geist in ein tönend Wort,}\\
\textit{Doch der Freie wandelt im Sturme fort.}\\[2pt]
\textsc{Schiller}.\\[4pt]
\textit{(Thou imprisonest the spirit in a sounding word,}\\
\textit{yet the free one journeys on in the storm.)}
\end{flushright}

\bigskip

\noindent\textbf{60.} In discussing the development of religious ideas the weight falls
quite decisively on feeling's selecting, strengthening, or inhibiting influence. The
religious consciousness gets its content through a constant choice of quality. But in this
process the relations that obtain among the ideas themselves also co-operate, and besides,
the reality ascribed to the content of ideas won may be of somewhat different kinds. In
these respects \emph{the myth} is different from \emph{the legend,} and both from \emph{the
dogma,} which again is different from \emph{the symbol.} The difference between dogma and
symbol is the most decisive for the religious problem in our days; but in the psychology of
religion it will also be of importance to consider the relation of both to myth and legend.

\emph{Mythology} and religion do not coincide. A mythology can develop on religious ground;
but religion---when we conceive its psychological nature as has been done in what
precedes---need not lead to the formation of myths. On the other hand there are myths that
have no religious significance, or that acquire it only later. The myth is formed
involuntarily, as a form in which the human being makes vivid to himself relations and events
in the world, especially in the outer world. The myth has an animistic character, in that an
event is made into a story that takes place between personal beings. Mythology can be taken
into religion's service; as religious feeling is awakened, it emphasizes individual
constituents of the myth and thereby transforms it more or less. In the Revelation of John,
for example, old mythological ideas are taken into the service of a new religious feeling;
especially in ch.~12 this is clear.\textsuperscript{69} Mythology can also, independently of
every religious motive, be taken into art's service, in that the imagination shapes its
originally crude or naive figures and traits into beings and actions with a clear stamp of
individuality and brings a more definitely motivated coherence into the narratives.

\emph{The legend} stands nearer to religion than the myth as a rule does. The legend is the
religious saga. Its kernel consists in the idea of a superior personality who has intervened
deeply in human life---has awakened enthusiasm, furnished a model, and paved new ways. There
then arises under the influence of memory a strong feeling-expansion, and as a consequence of
it a need for making vivid and for giving an account, which can set an image-forming process
in motion. At the individual points in this image-forming process, historical tradition and
the formation of myth work together, often in a very complicated manner.

In the myth there prevails especially the need to shape images by connecting as many individual
traits as possible. The formation of myth has the stamp of luxuriance, spreads according to the
laws of association by contiguity. In the legend the central interest in the subject prevails
more, the centre-seeking force; it rests more on inwardizing memory than on naive personifying
and painting-out. It is due more to analogizing than to combining.---In the myth, things and
relations are personified; in the legend one starts from the idea of a personality who is
illuminated by traits and relations from other domains, this illumination serving essentially to
express his value.---While the immediate feeling of unity with the highest prevails in mysticism,
it is in myth and legend the imagination that prevails. And the legend stands again in closer
relation to that feeling of unity than the myth, which often strikes into ways and directions
that concern rather life's outer circles than its centre. But as the legend can pass over into
myth (without its being able to be asserted, however, that all myths have arisen in this way), so
the myth can also be taken into the legend's service, especially in such a way that it serves for
the working out of individual traits in the legendary personality's character and fate.---

In the myth and the legend immediate intuition prevails, and the association of ideas works
involuntarily, only with the limitation that interest in the subject---with the legend, then, in
the individual person who calls forth the legendary process---effects. \emph{The dogma} stands
related to myth and legend as thought proper stands related to intuition and the association of
ideas. The dogma presupposes analysis---comparison, especially distinction. The images are
compared, and their individual traits brought forward and weighed. The one image must be brought
into agreement with another, either so that they cover one another, or at least so that they
become compatible without contradiction. And here in particular an involuntarily developed cult
and liturgy can co-operate (as with the dogma of Mary's immaculate conception). In the formation
of properly personal gods out of momentary and special gods---in the transition from polytheism
to monotheism---and in the development of the ideas of the divinity's different properties and
different revelations, there is use for this process of thought. It is during this process of
thought that---with greater or lesser consciousness---use is made of philosophical ideas and
reflections. Here the points are to be sought where the interweaving of religious and
philosophical elements takes place, often so intimately that a later critical investigation can
only with great difficulty distinguish the individual elements again. In the meantime the
religious consciousness has often to such a degree lived itself into the results of the dogmatic
process of thought that it thinks it here stands before immediate results of experience; it
confuses the feelings that the formed dogma (especially in its connection with acts of cult) is
able to call forth with the feelings by which in its time the dogmatic process was set going, and
it then thinks that in its dogmatically determined experiences it lives through anew what happened
at religion's original breakthrough.

The dogma forms an analogy (often also an impulse) to the artistic application and working-up of
myth and legend. But it is, as we have already had occasion to remark, not a purely theoretical
interest that prevails at the dogma's coming-into-being. The dogma comes into being by reason of a
need for fixed ideas in contrast to the myth's and the legend's changing forms and manifold
displacements. And this need announces itself not merely in individuals, each for himself. It makes
itself felt especially where a community has formed which will guard certain religious ideas and
reject ``false doctrine.'' Then there are needed distinctions that were formerly neither necessary
nor intelligible. One sets up dilemmas that formerly could not come forward, and one seeks by the
road of logical thought to decide these dilemmas by interpretation of the traditions. Where such an
interpretation is approved by the religious community's highest authority, the dogma is completed.
The concept of dogma therefore presupposes the concept of Church and rests on a distinctive mixture
of reflection and authority.

So long as the dogma is still living, the interest of feeling that already stirred at the
coming-into-being of the myths or legends it presupposes will always be traceable. When a dogma is
formed and asserted, or when a dogma is derived from another dogma, it is not purely disinterested
inferences that are drawn. A more or less pronounced religious, or at any rate ecclesiastical,
interest will always be co-operative. The ideal of a dogmatics it would be if each single dogma sprang
immediately from religious feeling, and yet all dogmas stood in the most beautiful logical harmony
with one another. This ideal was set up in the nineteenth century by Schleiermacher and Newman. But it
is an ideal that in this definite form becomes possible only on the basis of the sharper distinction
between knowledge and feeling that more recent psychology has inculcated. The history of dogma shows
us that strict logical thinking and purely religious feeling have not been alone determinative. Partly
mere associations of ideas, partly the ecclesiastical interest in mediating and in procuring uniformity
in doctrine and cult, have decided the direction of dogmatic development, quite apart from still more
``human'' motives and means that worked along with them. The dogma is then a product of highly manifold
elements.

A distinction has aptly been drawn between dogmas of the first degree, which express a religious
experience as immediately as possible, and dogmas of the second degree, which serve partly to connect
the dogmas of the first degree with one another, partly to bring coherence between them and the outer
conditions of the world.\textsuperscript{70} In Christian dogmatics, for example, the dogma of the
Trinity is a dogma of the second degree: it brings connection between the ideas that had formed about
God, about Jesus, and about the Spirit prevailing in the congregation.---There might still be reason to
speak of dogmas of the third degree. These would be such as established guarantees for the truth of the
two first groups of dogmas. To these belong the dogma of the Church, of the infallibility of the Pope
or of the Bible. These dogmas of the third degree have a---psychologically easily explicable---tendency
to spread at the expense of the two first groups, at any rate to come more and more into the
foreground, so that it is through the third degree that one comes into contact with the two first
degrees.

\medskip

\noindent\textbf{61.} The difference between the dogma on the one side, the myth and the legend on the
other, rests according to the foregoing especially on the fact that the dogma presupposes reflection
and authority. But this difference is not complete. In myth and legend already, reflection can often be
traced, even though it there expresses itself in a more subordinate, naive, and sporadic manner. Nor is
the element of authority quite lacking in myth and legend. These thrive only within a community where
the need to communicate can find vent and give the imagination life and strength for vivid forms. And
they live in the tradition from generation to generation, whereby they---besides the weight they owe to
their subject---acquire a special stamp of venerability that awakens immediate adherence, and which
furnishes the greatest strength to the resistance that historical criticism meets when it seeks to reach
back to the first origin of all these spiritual processes.

On the other hand the dogma does not quite lack that mark of vividness which is so strongly prominent
in myth and legend. The need for images lies deep in the essence of all religion, whether it leads to
favouring outer, artistic image-formation or not. In the dogma this need co-operates with the need for
conceptual fixing and limitation. Every image points (logically regarded) out beyond itself and has
(psychologically regarded) an inclination to draw other images after it; besides, every image has a
defect when it is to express a content that is in itself unsurveyable; every similitude limps. Therefore
the content of the legends and the myths is fixed in the dogma's form. That the figurative does not
thereby fall away, a closer investigation of any dogmatic concept can show.

In the philosophy of religion we cannot therefore make any difference of principle between myth, legend,
and dogma, such as every confessional theology must make from its standpoint. But as the most important
of these three forms the legend will come more and more to stand. Its kernel is the effect that a
personality outstanding and exemplary in the originality and depth of its spiritual life has made on the
human mind. Long after the dogma is extinct, and after the myth has been relegated to the literature of
fairy tales, the legend will preserve its value for all personal and serious conduct of life. The
Gospels will maintain themselves even when the dogmas that are supposed to be built on them, or that
they themselves contain, have been pushed aside.

It is well intelligible that the religious consciousness, when its ideas have once crystallized into
dogmas, and when these are fixed by the venerability of tradition and through the experiences they have
effected, shrinks from giving up this crystallization. If the dogma is given up, or regarded as mere
symbol, then the complete conclusion, the clear definiteness, and the firm foothold are lost, and one
stands before the indefinite and fluid. The symbol differs from the dogma in that it more clearly makes
known the difference between the original and proper religious lived experience and the ideas in which
this experience is expressed. That is philosophically regarded a great advantage, though it is an
advantage that the religious consciousness as a rule does not acknowledge. One might perhaps assert that
the difference between dogma and symbol is not so great as it often seems, and that the religious
consciousness itself acknowledges it. For involuntarily the religious consciousness nevertheless ascribes
to the dogmatic truths a different kind of validity from the scientific truths and the experiences of
everyday life. Often enough one gets the impression that the Sunday thoughts belong in another region
than the weekday thoughts, and often enough one lives in the prose of weekday life as though the Sunday
thoughts were ``only'' poetry. Besides, an investigation of the manner in which the dogmatic ideas are
applied in practical religiosity will show that the application is constantly symbolic; it will show
itself especially when sermons are preached on the accounts of miracles. But however much the religious
consciousness---partly in more prosaic and philistine, partly in more ideal forms---may approach the
acknowledgement of the dogmas' symbolic character, it will nevertheless start and take offence when this
symbolic character is definitely asserted. The firm, limited horizon is abolished; the distinction that
is made between the lived experience itself and its expression awakens doubt, discord, and unrest;
security is gone; it is as though one should now sew without tying a knot in the thread. The very concept
of revelation now seems to fall away: for does not this concept presuppose that there is a greater or
lesser domain where the image and the thing in itself completely cover one another?---It is from this
consideration that the Pope has denounced ``modernism'' as heresy, just as it was from it that Luther in
his time broke with Zwingli.

It is undeniably a great question whether there can be a religious standpoint where no dogmas, but only
symbols, hold. By symbols one must then think not merely of the historically transmitted ones, but also
of freely chosen symbols in which the personal experience, delivered from the dogma's barriers, can
express its new lived experiences---new leather bottles for the new wine. It is a great question---but
great questions can nevertheless sometimes be answered in the affirmative. The answer will depend on what
result one comes to concerning the hypothesis that the essential thing in all religion is not a solution
of riddles, but a faith in the persistence of value.---

\emph{The symbol} arises by a kind of analogy of feeling. A feeling that is determined by the lived
relation between value and reality seeks expression, and finds this through ideas that already stand as
expressions of analogous experiences. It is especially an analogy between the secondary and the primary
value-judgments (see §26) that here operates. The religious experiences are expressed by the help of
ideas drawn from the domain of self-preservation and of devotion. In his practical relations, during the
struggle for life's primary values, the human being forms his circle of ideas, and the religious
experiences are shaped by the help of this circle of ideas, without wholly new ideas being formed---or
being able to be formed. Thus the idea of the Father is a symbol whose use is due to the influence of
experiences of the furthering and shielding that existence can afford to what is humanly valuable; the
idea of the devil is a symbol that is applied under the influence of the opposite experiences. In all
symbolizing, ideas from narrower but more vivid relations are used to express relations that on account
of their exaltedness and ideality cannot immediately be made clear. In religious symbolism the analogy
rests on the relation to feeling. But a direct and positive determination is not reached by this kind of
analogizing. It brings it ``only'' to a poetry, not to an objective doctrine.---The word ``only'' must
here be thought of as used from the dogmatic standpoint. I myself stand at a standpoint where it is
precisely a sign of the highest that the poetic form is the only possible one.

Where the symbol is genuine, it springs from the immediate experience and the need this has awakened. It
is drawn from all domains that are accessible to human experience, but it will especially be the great
fundamental relations of nature and human life---light and darkness, power and powerlessness, life and
death, spirit and matter, good and evil---that furnish means for symbol-formation. The individual element
in existence is raised up to stand as a characteristic mark of the whole of existence---as comprehending
this in itself according to the relation experienced in it to obtain between value and reality. The
symbolic conception regards existence \emph{as if} its essence were exhausted in the individual element,
the individual experience of a harmonious or disharmonious relation between value and reality. Here there
is a kinship between the religious-poetic symbolism and the scientific use of analogy, alongside the
difference I have earlier brought forward.

The need from which the free formation of symbols proceeds has stirred in human beings at all times. It
lies at the foundation in animism, and it works in myths, legends, and dogmas, in all fairy-tale poetry
and in all art. By the acknowledgement of the symbolic in all religious ideas the religious consciousness
approaches the artistic conception, for which too the individual phenomenon stands as typical---as a
nutshell in which a world-content is hidden. But it approaches quite as much the more primitive
expressions of religious life in myth and legend. In common with these it has the immediate relation to
experience and feeling, while it separates itself from the dogma on account of the latter's reflected and
authoritative character. It does not have dogmatics' suspicion towards the involuntary formation of
religious ideas. If only there is immediacy and personal genuineness in the image-formation, the symbolic
conception sees no danger in the formation of myth and legend; they are natural forms for human spiritual
life at certain definite stages of development; indeed---what is more---they are perhaps the forms in
which the most intensive spiritual production of which human beings outside the strictly scientific domain
are capable takes place. All conscious art and speculation feeds on the content that first came forth in
the form of myth and legend---analyses it and combines its elements in new forms. Therefore religious
symbolism, even if it succeeds in freeing itself from the confessional and dogmatic barriers, will
certainly always continue to draw from the treasure laid down in myths and legends. At any rate this
treasure will assert its place alongside the new production of symbols, whose possibility cannot be
dismissed. That the capacity to form new symbols is nowadays so slight finds its explanation partly in
dogmatism's having cowed it, partly in the fact that the criticism of dogmatism has developed analysis and
doubt at the expense of secure, positive producing. The struggle for and against the dogmas has for so
long a time laid claim to the most important spiritual forces that it will perhaps be long before
sufficient freedom and strength of spirit can develop, not merely to preserve an upright posture in life's
highest questions, but in this upright posture to perform the great spiritual work: to shape the most
inward and essential experiences in mighty and unforgettable images. What worked in a childlike way in
myth and legend will then come again in manly shape.---That must he hope who believes in the persistence
of value, and who finds contained in the religious phenomena values that must be preserved in new forms
when the forms in which they have hitherto predominantly been clothed fall away. Faith in the persistence
of value holds fast that in spite of the division of labour in the spiritual domain which causes the
religious problem's coming-into-being, the real values that the spiritual life prior to the division of
labour possessed will not be lost. It is a faith in analogy with the faith we have in a friend of our
youth: that in manhood's dividing and divided work he will not deny what was essential in the ideals of
his youth, however many metamorphoses these must undergo in experience's purgatory.

\medskip

\noindent\textbf{62.} For the closer illumination of the coming-into-being of dogmas and
of their transition into symbols, a series of examples of the formation of religious ideas
shall here be brought forward, which take place under feeling's constant influence, but are
also clearly determined by the ideas' own laws. These examples are especially to show how
ideas that first appear in narrower and lower forms can be altered, extended, transferred,
and acquire ethical or cosmological significance.

Such a change does not, as the history of religion shows, take place under all conditions. A
characteristic difference here appears between Romans and Greeks. The Romans held tenaciously
to their special gods, as these had once been represented and worshipped. They were satisfied
when they merely had a reliable catalogue of the gods who were to be invoked, and their
religious imagination was not so lively and productive that it felt a need for a further
flight out beyond the circle of what was transmitted and what was shaped for cult. A richer
mythology the Romans got only through Greek influence. Among the Greeks the religious
imagination had worked luxuriantly. It supported itself on cult, and it was regulated by cult.
But the very customs of the cult are used by the imagination as motives for new formations of
ideas---especially when old customs seemed to need explanation, or when it was precisely the
ceasing of old customs that seemed to need explanation. In the myth of how Prometheus
cunningly gets Zeus to choose the bones of the sacrificial animal while he himself keeps the
meat, there is traced a need to explain why human beings themselves ate the best parts of the
sacrificial animal. A whole series of myths among Greeks and Romans, Indians and Jews, furnishes
an explanation of the ceasing of the old human sacrifices by the gods themselves wishing animal
sacrifices instead. Animal worship, which hangs together with totemism, is altered by the animal
shape's being demoted from being a real image of the god to representing the god's companion. A
remnant of animal worship is had when the god himself appears with more or fewer of the animal's
organs or properties, which then naturally receive a symbolic interpretation. The worship of the
heavenly light and of the sun is at a later stage motivated by light's symbolic significance as
an expression of purity, truth, and salvation.\textsuperscript{71}---It has already been
mentioned above (§51) how the liturgy (in the worship of the Madonna) can become an intermediate
link between legend and dogma.---

An interesting example---probably the very most interesting example in the history of
religion---of a transition from liturgy to cosmology is presented by the conception already
mentioned (§44) of Bráhman as the principle of the whole of existence. For Bráhman means first
yearning, prayer (or: the magical word). The motive for the transition from liturgy to cosmology
must be sought in the magical power that was ascribed to prayer: he who can rightly pray gets
power over the gods;---prayer, or the force that expresses itself in prayer, must then be the
real world-power! Even the gods notice this: they fear the mighty prayer, and they pray
themselves! The power of prayer then becomes the god above all gods.---Later on (in the
Upanishads) a transition is then made (cf.\ §52) from cosmology to psychology, in that Bráhman
is identified with the soul, which everyone knows from himself.---In both transitions, both in
the transition from liturgy to cosmology and in the transition from cosmology to psychology,
there works the need to represent the world's inmost ground as one with the highest goal of all
striving, or as a striving.---When finally in Indian religion the last step was taken: to declare
every expression and every thought insufficient, then free religious symbolism lay so
extraordinarily near that it can only be explained by the distinctive Indian conditions of culture
that it was not carried through. There set in at this point in the Indian development partly a
mere repetition of what had once appeared, partly stagnation, partly a sinking back into fetishism
and polytheism. No necessity can be pointed out for its having to go thus under all historical
conditions.---

In the opinion of some investigators the idea of a judgment of the world and a coming kingdom of
God, as it appears in post-exilic Judaism, can only be explained by the influence of Persian ideas.
But even if one accepts this opinion, which has been maintained in detail by \textsc{Erik
Stave},\textsuperscript{72} it must nevertheless be held fast that the ideas taken up acquire an
altered form and significance through the deeper conception of evil in the world and of the human
being's position towards it which had formed beforehand among the Jewish people. Persian dualism had
to be sure an ethical character, but the ethical was here nevertheless throughout grown together
with more external elements. The struggle against evil was in Parsism to a great extent a liturgical
or magical struggle, in that it was carried on by prayer, sacrifice, and purification; to a great
extent it was also a work of culture, in that what was harmful in nature was driven back by
agriculture, cattle-breeding, and the hunting of beasts of prey. The ethical element did not yet
appear so purely as it appeared---already prior to the taking up of those Persian ideas---among the
best in the people of Israel. For the quiet ones in Israel the whole question lay essentially in the
domain of the inward life, and in their expectation they looked more and more towards a purely
spiritual redemption.\textsuperscript{73} The thought of God's kingdom became more and more an
ethical thought.

When this thought was gradually developed into being the fundamental thought of a whole world-view,
thus into having not merely national and ethical but also cosmological significance, that became
possible only under the influence of rabbinic and Greek speculation; and at the definitive fixing of
the Christian Church's dogmas it was Greek speculation (in part perhaps also, as Harnack has shown,
Roman jurisprudence) that in fact furnished the conceptual forms that were used. Plato and Aristotle
would hardly have acknowledged their concepts as the Church Fathers and the Schoolmen applied them.
But the Christian Church's intellectual requirements were satisfied under the given conditions by
this ``pagan'' thought, or by what was taken for it. There is lacking in this whole development of
dogma---however interesting it is from certain points of view---the mighty and inward energy of
thought that in India led from liturgy to cosmology, thence to psychology, and thence again to the
limit of all thought. It appears far too strongly in this ``Christian thinking'' that it is alien
forms of thought that are being used.---

As a last example of the formation of religious ideas and its relation to the dogmatic I shall refer
to the history of so-called ``natural'' religion.---So-called natural religion is the result of
attempts at reduction that have at different times been undertaken with respect to the positive
religions, thus by Stoics and popular philosophers in antiquity, and among the so-called Deists in
the sixteenth, seventeenth, and eighteenth centuries. It retains a few dogmatically shaped ideas
(especially of a personal God and a personal immortality) and regards them as ``rational'' or else as
immediate experiences. Even where---as in Kant and the Kantians---the significant step has been taken
of regarding all religious ideas as symbolic, the dogmatic tendency can nevertheless still show
itself in this, that one regards certain definite symbols---and these are then as a rule those that
the ideas of ``natural'' religion furnish---as alone justified, necessary, and
``correct.''\textsuperscript{74} But in principle ``natural'' religion stands before the same problem
as the positive religions: before the problem of a connection between the ethical-psychological ideas
and cosmic reality, or between value and reality. When certain symbols are said to be ``more correct''
than others, this can consistently only mean that they acquire greater worth for us by resting on
analogies that lie nearer to us.---

The transition from dogma to symbol hangs closely together with the more definite carrying through of
the difference between feeling and idea. This difference (which in a more indefinite manner was known
to the Neoplatonists, Augustine, and the mystics) hangs together in turn with the differentiation of
human spiritual life, which stands generally as the cause of there existing a religious problem. The
evaluation and the explanation of existence no longer immediately coincide: the explanation does not
follow without more ado from the evaluation, nor the evaluation without more ado from the explanation.

That differentiation was in turn conditioned chiefly by the coming-into-being of the modern scientific
conception of the world. The extension and infinitization of the world-system that became a consequence
of \emph{Copernicanism} had to make it impossible to regard human life and its ideals, in the same
manner as hitherto, as the central thing in existence. And the clearer conception of the difference
between the spiritual and the material that \emph{Cartesian philosophy} brought entailed a criticism of
the animism on which religion constantly supported itself. Finally \emph{critical philosophy's}
investigation of knowledge's nature and limits marked a decisive turning point as regards the
intellectual elements in religion.

So much the more definitely, then, does the question arise anew what is the essential thing in religion,
and how far it can preserve its vital force under the new spiritual conditions.

\begin{center}---\end{center}

\section{The Proposition of the Persistence of Value}
\label{sec:persistence}

\begin{flushright}
\textit{Das Sein ist ewig, denn Gesetze}\\
\textit{Bewahren die lebendigen Schätze,}\\
\textit{Aus welchen sich das All geschmückt.}\\[2pt]
\textsc{Goethe}.\\[4pt]
\textit{(Being is eternal, for laws preserve the living treasures}\\
\textit{with which the All has adorned itself.)}
\end{flushright}

\bigskip

\noindent\textbf{63.} In the preceding sections reference has been made in several places to a
definite hypothesis about the essential content of all religious myths, legends, dogmas, and
symbols. The religious fundamental proposition in which all religion's inmost tendency can be
expressed has been maintained to be the proposition of the persistence of value. If this
conception is correct, the religious problem forms an interesting analogy to all other
fundamental problems. What again and again sets human thinking in motion in the different domains
is the relation between unity and multiplicity, or between continuity and difference. The
religious problem stands---if its content is the possibility of asserting the persistence of
value---as a distinctive form of the one great riddle that splits into distinctive forms in the
special domains. The proposition of the persistence of value is a special form of the principle of
existence's continuity. It forms an analogy to the causal principle, which indeed also---in its
way---asserts an inner connection in existence in spite of all differences and changes. In
particular it forms an analogy to the proposition of the conservation of energy, which sets up a
connection in nature amid the interplay of the forces of nature. Naturally this analogy to other
problems and propositions is not enough to prove that the proposition named really is the religious
axiom. It speaks in its favour, since on account of the human consciousness's identity with itself
in all domains it may be conjectured that the human problems must present common traits. But the
grounding that may perhaps be said to have been given already in the preceding presentation for the
assertion that the proposition of the persistence of value is the religious axiom is nevertheless
essentially of a psychological nature. By description and analysis of religious experience and
religious faith it was attempted to show that that proposition corresponds to the religious need,
in that this aims at holding fast a persistence of the highest valuable things out beyond the
barriers experience sets, and in spite of all the vicissitudes experience brings. Or---as it was
expressed---faith is faithfulness, and faith's content is that there is faithfulness in existence.
Faithfulness is persistence, continuity through all changes. This grounding presupposes that the
given description and analysis of religious experience and religious faith is correct. But even if
one were to concede this, the grounding given hitherto would be incomplete. An objective historical
confirmation is required, a demonstration that the actually appearing forms of religion, especially
the great positive religions, as regards their essential content, can be said to rest on the
proposition of the persistence of value as on a final presupposition. Such a historical verification
of the proposition shall now be attempted; but beforehand I must (in connection with what was
developed in section A on religious experience and faith) give the proposition a closer
determination.---If it then shows itself that history confirms what the purely psychological
investigation led one to conjecture, I shall pass over to discussing how far it is possible to hold
fast the proposition of the persistence of value at a standpoint that lies outside all positive
religion.

\subsection{More Precise Determination of the Proposition of the
  Persistence of Value and of its Relation to Experience}
\label{ssec:value-determination}

\noindent\textbf{64.} Religion presupposes that the human being has experienced that there is
something valuable. Whatever else one may think about religion, one must concede that it does not
itself produce all values in the first instance. If the human being believes, for example, in
another life full of goods or evils, he must indeed know from his own experience that there are
goods and evils; otherwise he can connect no meaning with his faith. And when the human being
ascribes to his gods certain excellent properties, he must indeed have come to know these
properties and their value from experience. If that is not the case, there is no connection between
the human being's religious state and his other states, and therefore no connection in human life as
a whole either. Religion's content always depends on the human being's experience, and especially on
what he has found to be valuable. Which values a human being finds depends in turn on the motive of
evaluation that prevails in him. The motive of evaluation may appear as a momentary need that is soon
succeeded by another need. Or it may lie in a deep, constant yearning that is one with the instinct of
self-preservation. It may be conditioned by the human being as a single isolated being, and thus have
an individualistic, perhaps even egoistic character, or it may spring from the fact that the human
being cannot separate his cause, his weal and woe, from the great connection in which he feels himself
to stand as a member of comprehensive human groups and as a participant in great, common interests.
The values in whose persistence the human being believes will be those that for him are the highest;
but they may in the different cases, in different human beings and under different historical
conditions, be highly different. The Greenlander's faith in the persistence of value is highly
different from the Greek's, the Indian's from the Christian's. The egoist and the pleasure-seeker can
have their heaven, just as well as the one seized by beauty or the ethical idealist; but these heavens
will look very different. These differences strictly speaking do not concern us here, where what
matters is to find what is common to all religion---that which brings it about that we ascribe a
religion to the Greenlander just as well as to the Greek, to the Indian just as well as to the
Christian. The opposition between ``lower'' and ``higher'' values is to be sure of great importance for
religion; it is this that makes us distinguish between nature-religions and ethical religions, and
within the ethical religions again between lower and higher forms. But it does not become decisive for
the religious axiom itself.

Since all religion presupposes experience of values, the religious values themselves must in a certain
sense be derived, namely conditioned by the interest in the primary values that the experience of life
has led one to know and to assert. Religion presupposes a special experience that in the course of
things in the world it is the fate of values that is at stake. A life of the soul of the most elementary
kind cannot know religious values, but will be limited to the simplest forms of self-assertion and
devotion. The animal cannot (probably cannot) be religious, because it cannot (probably cannot) make
special experiences of the fate of its values in the world.

But---as already remarked in an earlier connection---from the fact that the religious values are derived
in relation to other values it does not follow that there must always be a temporal distance between the
experiences in which they make themselves felt and those in which the primary values make themselves
known. Often primary values first come into being in religious form, so that the two kinds of experience
are made at once. The difference is due to an abstraction, or at least a differentiation, that need not
set in. I can perfectly well immediately conceive a beauty in nature that is new to me as testimony to
existence's glory; the primary value and the religious value---beauty as such, and existence as that
which makes beauty possible---then appear for me at one stroke.---And even if the difference between the
primary and the religious values makes itself felt, the religious value can nevertheless preserve its
immediacy and its independence. It is not necessary that the object of religious faith stand merely as a
means for the persistence of the primary values in existence. It can itself appear as the highest good,
as the object of immediate admiration and enthusiasm, trust and love. The difference between lower and
higher religions will rest on whether the religious values stand only as mediate, or whether they appear
as immediate. The transition from the lower to the higher forms can here, as in other domains, take place
through displacements of motive and of value. What at first has value only as a means can later acquire
value as an end; what at first has value from a certain cause can later acquire value from a wholly
different cause; and from what originally stands as a final end there can develop effects that far
surpass the original end in value. Such displacements are of the greatest importance not merely in the
ethical domain, but also in the religious domain. The continuity in humanity's religious development is
due to a great process of displacement which was supported by significant historical and personal lived
experiences. Often the religious consciousness oscillates between the mediate and the immediate value:
when Augustine says to his God, ``I seek thee, \emph{that} my soul may live!'' then in this utterance he
ascribes to the object of his faith only a mediate value; but when in other statements he calls God the
highest or only good itself, goodness and truth itself, then the object of his faith appears with
immediate value. For Spinoza ``the intellectual love of God'' arises through the fact that the full
understanding of himself and of his unity with the whole of existence awakens a deep joy of knowledge,
and that with this joy there combines the idea that it is due to God (the principle of unity in
existence) as cause; but on closer absorption in this thought it presents itself thus, that in the
understanding, the joy of knowledge, and the intellectual love, God himself is present as operative, so
that the merely mediate relation falls away. All mysticism asserts---in contrast to the external,
mechanical, and dualistic character of ordinary religiosity---the immediate character of the religious
values. The principle of the persistence of value in existence comes at last itself to stand as the
highest valuable thing. On this rests the fact that religion can exercise a reaction upon the other
spiritual domains and enter into more or less pronounced opposition to them. There is a reciprocal action
between primary and secondary values that makes the religious problem quite particularly complicated. It
may be difficult to unravel what in the individual case is the original and what the derived. And above
all: when the form in which the persistence of value appears itself has immediate value, the sum of the
values whose persistence is in question is of course increased, and the very persistence of religion
becomes a part of the religious problem. It is not enough to show that the primary values do not suffer
by positive religion's falling away; with it there fall away immediate values, for which equivalents
must be pointed out if one will maintain that positive religion can disappear from human spiritual life
without any loss being suffered. It would be a confirmation of the religious axiom's validity if such
equivalents could be pointed out. The continuity of value would be preserved---and if this is the
properly religious thing, ``the purest thing in religion,'' then religion would persist even if positive
religion fell away. The religious problem here shows itself in its unavoidability and in its great
importance, however one may stand towards the existing forms of religion.

\medskip

\noindent\textbf{65.} We possess much that is valuable which we do not uninterruptedly enjoy. A state
that makes us happy need not be uninterruptedly present in its whole strength. There may be stops that
do not mean the melody's cessation, but are merely temporary conclusions, stages of rest or of
preparation. The pause is itself a member of the melodic sequence and makes its definite effect within
it. Such a pause can only be appreciated by one who preserves the continuity with what precedes, and
who also lives through what follows. He who arrives just as the pause begins experiences nothing; for
him the pause is a nullity. And for him who departs before the pause is succeeded by new tones, it is
one with an absolute cessation. The pauses of the world's course may be very long, and only he who can
place them in the most intimate connection with what goes before and what follows after can understand
the pauses' value and see that they are more than interruptions. Whether in the individual case we have
a pause or an absolute cessation may stand as an insoluble question. If continuity is to be asserted, it
can only be by a faith. It may look as though there were absolute losses of value.

We need here---if continuity is to be held fast---the distinction between potential and actual value.
Faith in the persistence of value does not presuppose that there is always just as much actual value in
existence, if only there is always just as much possibility of value's arising. But if all values passed
over into being potential, it would be the same as a dead state of equilibrium setting in. The
possibilities would become impossibilities. Religious faith will in its special forms acquire a different
character according to the relation it presupposes between actual and potential value in actual
existence. It will acquire a bright and optimistic character if it relies on the potential values' being
able to be released without all too bitter resistance; a dark and pessimistic character if it assumes
that existence is of so tragicomic a constitution that the richness of potential values it contains in
and for itself will never be able to be released, or that every single process of release will entail a
disproportionate loss of value at other points.

For a complete carrying through and verification of the proposition of the persistence of value it would
be required that nothing in the world's course be left standing as mere means, or as mere possibility,
let alone as mere hindrance, but that what had mediate value should always at the same time have
immediate value, and that all hindrances should at the same time be means. In the religions that assume
the divinity to be both almighty and all-good, such a presupposition of the persistence of value lies at
the foundation. And it may be said to be an ideal by which every religion can be measured.

It is of distinctive interest to see that the ideal form of the proposition of the persistence of value
forms an analogy to, and to a certain extent coincides with, the highest ethical principle. The ethical
ideal---as it shapes itself on the basis of universal sympathy---appears in the thought of a kingdom of
humanity in which each single human personality always stands as an end and never merely as a
means,\textsuperscript{75} thus as one that has immediate value and never merely mediate or potential
value. Just as this ethical ideal appears with greater or lesser clarity as a presupposition, or at least
as a tendency, in the special ethical rules and laws and at the different stages of development, so that
ideal form of the proposition of the persistence of value in existence appears with greater or lesser
clarity as a presupposition or tendency in the special forms of religion. And if it should now later show
itself that the final measure of value for all religion is of an ethical nature, then this analogy between
the religious axiom and the highest ethical principle is of great importance, since it might indicate that
religion tends to stand as a projection of the ethical. In the religious fundamental presupposition itself
it then lies that religion must at last be an object of ethical criticism. It must for example be
investigated whether the all-goodness with which almightiness is said to be combined really is---goodness!---

Every single religion will naturally lay quite particular weight on the form in which the religious axiom
appears in its doctrine and cult. It will have an inclination to confuse the clothes with the person, or at
least to make the clothes a part of the person. That the proposition of the persistence of value is the
religious fundamental proposition is not disproved by a single religion's refusing to acknowledge that
other religions in their way espouse and express the same proposition. This can be shaped and expressed in
many different ways, in myths, legends, dogmas, or symbols of highly different kinds. But there must---
philosophically regarded---be a distinction drawn between the proposition itself and the special ways in
which it is expressed. The burden of proof in any case rests on him who asserts that it can only be
expressed in one definite way. The proposition that is to express religion's essence in all religion's
different forms must naturally have an abstract character. What matters then is that one does not confuse a
philosophy-of-religion proposition with ideas that appear immediately and consciously in and for the
religious consciousness. The religious consciousness need not know the proposition in its abstract form, and
will perhaps even reject it in this form. Thus the ordinary human consciousness need not know the laws of
the association of ideas at all, without this telling against the fact that these laws lie at the foundation
of every connection of ideas, are fulfilled in this same human consciousness. Thus we draw our breath
without needing to have the least intimation of the physiological laws of respiration.---

\medskip

\noindent\textbf{66.} Where attempts are made to derive reality from value, the religious problem is
sharpened. For even if all hindrances, all resistance could be shown to be necessary means for the
development and persistence of the valuable, the question will nevertheless always arise why means are
needed at all---why the valuable is not immediately and does not immediately rule---why there is a
difference or a strife between value and reality. In the problem of evil the religious consciousness has
constantly had to struggle with this difficulty, and at this point there is also sharpened, as will show
itself, the difficulty for the philosophy of religion in pointing out the fundamental proposition of the
persistence of value as a presupposition or as a tendency in all religion.

It might seem as though the hypothesis of the persistence of value must be incompatible with a pessimistic
view of life. And if pessimism too is a kind of religion, then it would furnish an objection to our
philosophy-of-religion hypothesis.

Pessimism assumes a fundamental disproportion between value and reality. But a disproportion is also a
relation, and when religious experience---on the description given earlier---concerns the relation between
value and reality, then pessimism is grounded on religious experience, only that its experience leads to a
different faith from that which most frequently appears in religion's history. And yet there must at the
foundation of pessimism too lie a faith in the persistence of value: for if all value fell away, the relation
between value and reality would naturally fall away too. The pessimist must assume that there is something
valuable in the world, but that it persists only under hard struggle and constant suffering, and it is this
struggle and suffering to which he has first and foremost directed his attention. Struggle and suffering
presuppose life and a need for self-assertion. If all life, all need in existence could be excluded, there
would no longer be any question of pessimism. In Leopardi's poem ``The Broom'' (\emph{La Ginestra}) humanity's
relation to merciless nature is presented by the broom-bush's relation to the lava on which it grows. When
the lava at some point burns up or overwhelms the bush, the struggle is over. But the plant's life is the
valuable thing whose opposition to crude reality calls forth or expresses the poet's mood. The most radical
pessimism would, as \emph{Simmel} has shown, be expressed in the assertion that the very fact that there is
pain and unhappiness in the world at all makes the world valueless.\textsuperscript{76} But no pessimist can
hold fast such an assertion. And when Leopardi does not---any more than Schopenhauer and Buddha---in spite of
the strong emphasis on the disharmony between value and reality, point out the way out of stopping life, thus
abolishing the valuable, in order to escape the disharmony, there lies in that necessarily a faith in the
possibility that the valuable can persist under the constant struggle and suffering. An absolute pessimism has
never at any time been presented, either in any religion or in any philosophy. Buddha points to the possibility
of attaining Nirvana, which is by no means one with annihilation, and for Schopenhauer the disharmony is
abolished in artistic and scientific activity, in compassion, and in religious asceticism. There is in all
pessimism something valuable that persists---that can assert itself and perhaps even increase.---An absolute
pessimism would hold only for those beings that in some religions are thought to be in eternal torment. But
this eternal torment is in no religion the only content, and there are made---as we shall later see---
characteristic attempts to show that the highest valuable thing persists not merely in spite of, but precisely
by reason of, the eternal torment in a part of existence. Precisely these attempts, which are due to the
reflection of pronouncedly religious natures, will furnish a testimony to our religious axiom's historical
validity.

One might imagine a pessimistic conception that assumed a constant diminution of value in existence. Such a
conception too would be determined by religious experience, although religious faith is lacking to it. If one
wished to lay the weight on the constant tendency to diminution of value, then such a pessimism could be called
religion with a negative sign. It would constantly be determined by the relation between value and reality and
would not denote the complete opposite of religion. The complete opposite of all religion would be neither
pessimism nor optimism, but neutralism: all evaluation beyond the domain of human actions falls away, and the
human being stands as a mere spectator at the course of the great world-process, convinced of its infinite
indifference to what human beings call value. I consider it very difficult to decide whether such a neutralism
can really be given. Even for the mere spectator there will in every intellectually or aesthetically enraptured
(let alone in an ethically enraptured) state arise a mood with a religious gleam: it gives existence a
distinctive value that the understanding or the contemplation of it causes us joy, and our view of the world is
involuntarily stamped by the advance or retreat of such intellectual or aesthetic values.

\subsection{Psychological-Historical Discussion of the Proposition
  of the Persistence of Value}
\label{ssec:value-psychological}

\noindent\textbf{67.} In two different forms we find in the historical religions the faith
in the persistence of value. We meet here with the same chief forms that we already stood
before in the description of religious experience and religious faith (§34). They are not
complete opposites in their essence, and historically they have in several ways influenced
one another; but they stand nevertheless as two types, such that the special forms of
religion, the different religious standpoints, approach more or less pronouncedly the one or
the other of them. Since it will be impossible to go through all special forms of religion
and religious standpoints, I hold to these two types, in the confidence that religious life
has in them expressed its most characteristic distinctive features.---Before I pass over to
considering them severally, I must nevertheless recall that the presupposition of the
persistence of value reveals itself even in the most elementary religion one knows (see §37),
namely through the preservation of ``the totemic principle.''---

According to the one type the highest value is always actually present. What hides it from
human beings' gaze is the motley multiplicity and constant change of the world of experience,
and the illusions about the reality of the sensible in which human beings are thereby ensnared.
When these illusions fall---which presupposes a serious work of thought and will---human beings
become clear about the eternal reality and about their immediate unity with it. All multiplicity
(or at least all outer, sensible multiplicity) falls away; the temporal relation loses its
significance, and the valuable shows itself as the only thing that always is and always has
been, the only true reality.

According to the other type the valuable persists only through the constant struggle against
forces that would check and abolish it. It has besides a development to go through before it
exists in its whole fullness;---it has a history. Only when the historical course of the world
has reached its conclusion is the valuable all in all. The temporal relation acquires here real
and decisive significance; it is not, as for the first type, an illusion that falls away with
true understanding. New values really come into being, and the work that must be done in order
to have a share in the valuable that persists is just as much a reality as the goal that is
reached through the work.

\medskip

\noindent\textbf{68.} The first type appears historically in the Indian Vedanta doctrine, in
Buddhism, in Platonism, in the mysticism of the Middle Ages, and in Spinoza's system. The
persistence of value is here asserted by the view that what in \emph{time} appears in scattered
multiplicity and in successive unfolding is in \emph{eternity} concentrated in an absolute unity.
Eternity lies here not \emph{after} time, but \emph{in} time, and it discloses itself in the human
being's inner life as soon as the human being reaches the deepest self-collection. This type must
consistently assume that there is a relation of equivalence between time and eternity, but this
relation of equivalence could not itself have any history, since all differences of time lose their
significance, and the very work of becoming a partaker in the eternal could have no reality.

One might think that this type is not suited for a popular religion. And yet Greek religion, as it
appears in the Homeric age, is essentially to be referred to this type, even though it here appears
in a childlike form. For the Homeric Greek, in the midst of human life's toil and struggle, Olympus
stood in an eternal clarity, untouched and unchangeable. Olympus is the gods' eternal and fixed
dwelling, always encircled by shining radiance, raised high above the regions that storm and rain
visit. In this radiant picture the Greeks saw the valuable's eternal reality expressed, and over its
glory they forgot the shadows of their own life, or they took, with melancholy and resignation, the
opposition between the Olympian and the earthly as something that was once for all so.

While there is here a naive but insuperable opposition between eternity and time, a more inward and
more thought-through relation appears in the Vedanta doctrine, and especially in Buddhism's doctrine
of Nirvana. Nirvana is, as already several times remarked, not to be thought of as pure nothing. It
is a state of peace and exaltedness, the opposite of desire, hatred, and disappointment, which thrive
in the world of sensible illusions. It is a state where holiness is consummated, where the work is
performed, and where ``this world'' no longer exists. In order to reach this state the individual must
withdraw from the ordinary human relations; but this happens with the most inward enthusiasm: ``Like
swans from the swamp, the understanding ones fly away from house and home in order to be free of
sorrows, illusions, and sensible enticements. Bold, with quiet peace in their minds and with senses
that have come to rest like the charioteer's tamed horses, they are the greatest of all heroes, envied
by the gods themselves. Desire and hatred the Buddhist has cast far from him, as the jasmine bush
shakes off its withered flowers. An unearthly blessedness seizes him when he steps into his quiet
cell, where he sees through life's vicissitudes and clearly knows the eternal.''\textsuperscript{77}

Within the Greek world a kindred line of thought comes forward in Plato. Only the world of ideas is
true reality, and everything depends on whether the knowledge of the ideas (things' essence and
models) dawns on the human being; only then is he freed from all discord and imperfection. Especially
in the dialogue \emph{Phaedo} this type appears in a characteristic manner. Through Platonism it has
also intervened in the development of Christianity, which otherwise belongs predominantly under the
second type. In the Gospel of John there are traits that go in the direction of the first, the
Indian-Greek type; but it is not certain that they have arisen through direct influence from it. On
the other hand such an influence is both clear and certain in Augustine. In him is found the Platonic
concept of the eternal as the unchangeable (\textit{semper stans æternitas}), which \emph{is} to an
equally high degree at every time, and within which nothing arises or perishes. It is found again
further in the mysticism of the Middle Ages, as when Meister Eckhart says: ``What is eternity?
Eternity is a present now that knows nothing of time. The day that passed a thousand years ago is not
further from eternity than this hour in which I stand here, and the day that shall come in a thousand
years is not further from eternity than the hour in which I now speak.'' To this belongs also Jacob
Böhme's motto: ``\textit{Wem Zeit ist wie Ewigkeit, und Ewigkeit wie Zeit, der ist befreit von allem
Streit}''---and Spinoza's profound instruction ``to see things under the form of
eternity.''\textsuperscript{78}

It is indubitable that all these different standpoints each in its way assert that all true value is
imperishable. A difficulty nevertheless arises on a closer consideration of the relation between
``time'' and ``eternity.'' If the proposition of the persistence of value were to be held fast
strictly, ``time'' and ``eternity'' would have to stand in a relation of equivalence to one another,
be two forms of value between which a conversion could take place without loss.

But that is not how it is. ``Time'' is at all these standpoints supposed to rest merely on an illusion,
an illusion that must dissolve into nothing if the true value is to be---not: come into being, but---
unveiled. Therefore the work and the development in time---however necessary they may be---acquire no
real significance, no reality, but are to be regarded as the exertions we make in dreams. The very work
of abolishing the illusion, of letting the dream-image sink away into the nothing that it is, and of
living oneself into the true reality, is an illusion. At most the work can come to stand as a mere
means that is forgotten when the end is reached, a ladder that is pushed away when one has climbed to
the height. From these standpoints no immediate and independent value can consistently be ascribed to
the work and the development---and thereby to the goods and the tasks that the world of time brings.

There is at the same time lacking here a natural motive for the individual to set himself other ends
than his own personal liberation from the illusion. When the individual has himself reached the
consummation, has torn asunder the veil of illusion, why should he then linger in the illusory world,
even if he could thereby help others to reach a similar perfection to that which he has himself reached?
Why should the swan return to the swamp?---When Buddha had reached the consummation in spite of all the
hindrances that the evil spirit (Mara, Buddhism's devil) laid in his way, the latter urged him at least
not to show others the way of consummation he had found. Buddha did not obey, partly out of ``compassion
for beings,'' partly also simply because his perfection suffered no diminution by being communicated to
others. But in and for itself the consequence lay nearest which is also uttered both by Buddha himself
and in the Buddhist poem \emph{Dhammapada} (the Way of Truth), that he who would reach the highest, which
is liberation from sorrow and fear, must love no one, since separation from those one loves causes
sorrow.\textsuperscript{79}

Besides the dualism between time and eternity, another dualism comes forward at these standpoints. It is
namely not everyone who attains Nirvana, and Buddhism therefore quite naturally developed a doctrine of
hell, or took it over from popular belief. Nor is it everyone who attains to knowing Plato's eternal ideas
or to acquiring Spinoza's capacity to see things under the form of eternity. In one form or another the
old Greek opposition makes itself felt between the eternal, radiant Olympus (the analogy to Nirvana, to the
world of ideas, to the form of eternity) and a dark human world. Only a few, the favourites of the gods,
are raised above this opposition. For the rest there is loss of value all along the line.

The result then becomes, as regards the first religious type, that the proposition of the persistence of
value is a tacit presupposition, but that in the closer development of the standpoints it is not unfolded
with full consistency. But the objection that thus strikes this type---that the proposition of the
persistence of value is only defectively carried through---is no alien objection. It arises through the
application of a standard that this religious type itself acknowledges. Human life outside Olympus and
outside Nirvana is nevertheless not absolutely valueless: the Greek has his childlike joy in life, and the
Buddhist must at any rate ascribe a real value to the yearning out beyond time that can lead to Nirvana.
Olympus and Nirvana cannot then consistently occupy a uniquely privileged position, untouched by what goes
on in the finite world. Only through living and positive connection with this can they mean the highest
valuable thing. We measure here religion by religion itself, in that we compare a religious standpoint's
individual traits with the goal it has itself set. Religion can be attacked from outside, from standpoints
alien to its inmost goal, its real tendency; but such a criticism would not help us to find the religious
axiom, the fundamental proposition that religion cannot deny without denying itself. If the proposition of
the persistence of value is a proposition that no religion can deny without denying itself, then this
proposition is the religious fundamental proposition.

\medskip

\noindent\textbf{69.} What distinguishes the second chief type is its historical character. It
believes in a world-history. In Parsism this historical character first appears in pronounced
form. The struggle between good and evil is carried on in the whole of nature as well as in
human life, and one looks towards a great world-struggle after which the consummation is
expected. In the Jewish prophets we meet an independent historical faith, grown up under the
influence of the inward need to hold fast the hope of the people's future and the conviction of
the national God's faithfulness through all vicissitudes. From the tribulations of the time one
gazes towards a day on which it may be said: ``Behold, this is our God, whom we trusted in, and
who has saved us!---Jahveh is a rock of ages!'' (Isaiah XXV, 9---XXVI, 4). In post-exilic
Judaism's messianic expectations and in Christianity's doctrine of God's kingdom and its speedy
coming, the same historical faith is expressed. It has set its significant stamp on the thought
of European humanity. Augustine's work on ``the City of God'' (\textit{civitas dei}) and
Bossuet's ``\emph{Discours sur l'histoire universelle}'' form the religious preparation for what
a later age has called the philosophy of history, sociology, or the history of culture.

The temporal relation acquires here positive, decisive significance. The difference between past,
present, and future does not rest on an illusion. And definite tasks are set that can only be
solved through work and development. There is a real ``not yet,'' as well as a ``no longer.'' It
is through the changes of the times and the phases of development that the highest is realized. A
comparison between the ``Upanishads'' and the New Testament is in this respect highly instructive.
Whereas there no relation of equivalence could obtain between time and eternity, here there is
supposed to be a relation of equivalence between past and future: life is won by being lost; there
shall come thousandfold compensation for what is sacrificed; Jesus has not come to dissolve but to
fulfil; the kingdom has been prepared from the foundation of the world, but is an inheritance that
is only now to be entered upon. Not even suffering and destruction testify against the persistence
of value, for the very suffering in the service of the highest and with the gaze turned towards it
testifies to its power. ``Thou art victor because thou didst let thyself be sacrificed!''
(\textit{victor, qvia victima!}) Augustine calls out to his Saviour.

And in the idea of the kingdom that is to come there lies not merely the significance of history,
but also that of solidarity, of the life of community, and the possibility of a more positive form
of love of humanity than Buddhism could consistently accommodate. The individual's consummation is
tied to that of the others.

There is thus a more realistic stamp on this type than on the first. And with it follows a greater
possibility of carrying through the persistence of value, since the world of experience and its
conditions do not stand as mere slag containing no valuable ore. The advantage this second type has
over the first hangs together with a more complete carrying through of the proposition of the
persistence of value.

When there are nevertheless difficulties here too for the demonstration of this proposition, they
lie first and foremost in the fact that the work and the development in time stand indeed as real,
but yet only as means. They have only mediate, not immediate value. They are only a preparation for
what is to come. And that which is to come, the divine kingdom, does not even come as a natural and
necessary fruit of the development and the work, but is inserted into time by a supernatural road,
so that the bond between past and future is cut through. There is after all no inner relation of
equivalence between them, since human work has vanishing significance in relation to the great
thing that is to come---indeed easily becomes a hindrance. The work becomes properly negative: what
matters is not to immerse oneself in human relations, in the goods and tasks of human life, but to
watch and pray in order to be ready when the consummation comes. Life in time becomes a life in
expectation, and only that has value which can prepare for the life to come. Hereby there arises a
certain likeness between Christianity (that is, primitive Christianity) and Buddhism, for it can in
practice lead to the same result to deny the essential continuity between past and future and to
declare the temporal relation as a whole illusory.

But at other points too the demonstration of the proposition of the persistence of value meets
difficulties within this type. And these points are of special interest, because the religious
consciousness has itself been attentive to them and sought to overcome the difficulties. It is of
great philosophy-of-religion interest to see how the attempt to overcome the difficulties is made
precisely by virtue of the proposition of the persistence of value. It is the greatest theologian
within Christianity's history, Augustine, who undertakes this attempt.

The first point is the dogma of eternal damnation. It is not necessary here to discuss whether this
dogma can be derived from the New Testament writings; it is enough that the orthodox church
doctrine, both the Protestant and the Catholic, maintains it. It would undoubtedly be easier to
carry through a proof that Christianity rests on the presupposition of the persistence of value if
it could be established that the New Testament teaches that all are at last saved. But what
interests us here is that precisely the dualism of the saved and the damned has for the thinking
religious consciousness been able to stand as a necessary consequence of the persistence of
value---naturally of the persistence of that value which for this consciousness was the highest
value. However the world's course may end, the Christian consciousness holds fast that it serves
for God's honour and glorification. Augustine seeks to ground this more closely. Divine justice is
satisfied by the wicked being treated according to desert and placed where it belongs to them to be.
Indeed, Augustine defends this dualism not merely as ethically necessary, but also as aesthetically
necessary. When the world's course ends with the eternal damnation of the many, this might seem to
entail an imperfection in existence's definitive state;---but, says Augustine, what seen from the
individual part's standpoint is imperfect and causes disgust may very well be a perfection from the
whole's standpoint, since for harmony it is merely required that everything come into the right
place, and since precisely the opposition between the saved and the damned serves to increase the
beauty of the whole (\textit{omnes ita ordinantur \ldots\ in pulchritudinem universitatis, ut qvod
horremus in parte, si cum toto consideremus plurimum placeat}).\textsuperscript{80} That dualism,
then, in which an objection to Christianity has so often been found, is for Augustine precisely a
testimony to the persistence of the highest ethical and aesthetic values. There expresses itself in
his line of thought undoubtedly a striving to hold fast the fundamental proposition of the
persistence of value---against an objection that was grounded on this same proposition. He will
transform that from which the objection was drawn into a confirmation.---How far the concepts of
value that Augustine here applies are valid is a question by itself. For Augustine there was no
doubt that God's blessedness and honour were compatible with, indeed even required, that dualism.
The consequence that has in more recent times been drawn from Augustine's premises, that God cannot
be blessed when the world's course has the character he assumed, lies far from him. Just as little
does it cause him difficulty how some human beings can be blessed when they know that not all are.
He has, so far as I know, not touched on this objection. But Thomas Aquinas has later answered it.
To be sure, says the great Schoolman, Aristotle teaches that no one can be happy without having his
loved ones about him; but that holds only for the earthly life; in the eternal life the individual
stands before God, and love of the infinite object fills him so that he has no need to love finite
beings; only love of God, not love of neighbour, is necessary to the state of blessedness, so that
even if there were only a single soul that came to enjoy God's presence, it would be blessed without
having any neighbour to love!\textsuperscript{81}---In the face of a line of thought such as this
one must well remember that the significance that faith in the persistence of value has in the
individual case depends on which values are known and presupposed. Love of humanity cannot at a
standpoint such as the one here discussed belong to the highest values, and it therefore does not
come along either when faith's content is to be determined. There lies in this a confirmation of the
conception of religion's essence that I seek to make good.---For the rest, Thomas Aquinas teaches
that there is a great difference between the feeling a human being ought to have here in life, where
he is still on the way (is a \emph{viator}), and the one he will have in blessedness, where he
clearly understands the connection (as a \emph{comprehensor}). For the earthly human being compassion
with the damned is praiseworthy, but such a feeling cannot arise in the blessed, who understand God's
judgments.\textsuperscript{82} Unfortunately neither Thomas nor any other theologian has succeeded in
giving us the ``higher'' psychology and ethics that are here pointed to, any more than they have
succeeded in giving us the ``higher'' epistemology which in another connection has shown itself
necessary for their standpoint. It has not even been possible to show how a ``wayfaring'' human
being, a \textit{viator}, can connect any real meaning with the ``higher'' standpoint to which
reference is made---an ethical meaning one cannot at any rate connect with it.---But, as said, this
does not properly concern the religious axiom in its generality.

The second point where a difficulty shows itself for holding fast the theological standpoint we here
discuss lies in the dogma of creation.---While Buddhism does not enter upon any explanation of how
``the world'' comes into being, Christian theology gives an explanation which might seem to contain
the assumption of a diminution of value. For the created is inferior to the Creator: it is finite and
limited, and there presents itself for it the possibility of a fall which does not exist for the
Creator in his eternal and ideal reality. But here too Augustine finds in his concept of justice a way
to save the persistence of value: ``Who is so out of his mind that he dares demand that the work
should stand equal with him who has produced it, the founded with the founder?'' ``It would be
blasphemous impudence to place nothing on a level with God by demanding that what God has created out
of nothing should be of the same nature as that which springs from God's own
essence!''\textsuperscript{83} The diminution of value at the creation falls away for Augustine
through justice being satisfied; through it continuity is preserved. We see clearly lying at the
foundation in the religious thinker's conception the proposition that I have conjectured to be the
religious axiom. In its consequence, however, Augustine's line of thought would lead us to a result
similar to that in which Buddhism (and on the whole the first type) ended. The divinity is thought as
unchangeable, untouched by creation and fall, as well as by the opposition between the saved and the
damned. God does not need the world and is not altered because the world is altered. Not even God's
goodness entails need and dependence, for since he is a perfect being that can receive no increase in
perfection, his goodness too can persist without needing other beings. This consequence Thomas draws:
\textit{Bonitas ejus potest esse sine aliis.}\textsuperscript{84} With what right it is nevertheless
still called goodness is a matter by itself. Here it was only to be emphasized that both religious
types---in spite of their indubitable striving to hold fast the persistence of value---come up against
difficulties they have not overcome---difficulties in which they show themselves as imperfect
expressions of the ideal of religion that they themselves acknowledge.

\medskip

\noindent\textbf{70.} The result of the investigation now made of the two chief forms of
religion is that in both of them there expresses itself a decided tendency to assert the
persistence of value, a tendency that appears quite particularly in the attempts to reject
objections drawn from an apparent diminution of value.

When nevertheless neither of the two chief forms succeeds in holding fast the proposition of
the persistence of value strictly, this comes not merely from the fact that historical forms
of life never fully express their own real essence or idea, because extraneous elements and
circumstances co-operate. It comes also from a circumstance that hangs closely together with
religion's essence, as this must be if faith in the persistence of value is the religious
axiom. It is namely clear that faith in the persistence of value must lay claim to forces and
interests that are thereby drawn away from the very work of finding and producing values in
the given world. Thereby there constantly arises a greater or lesser relation of opposition
between religion and the other spiritual domains. Since the religious relation itself gives
rise to evaluation and grounds values, this relation of opposition can also be designated as a
relation of opposition between the primary and the secondary values. Quite particularly does
this appear in the relation between religion and ethics, and it will be discussed more closely
in the ethical philosophy of religion. But in order that it might be understood why the carrying
through of the religious axiom must always be imperfect, reference had already to be made here
to the fact that faith in values easily comes to stand in a certain opposition to the work of
discovering and realizing values. In the two chief forms of religion this shows itself in the
relation to development in time, which does not come to its full, positive, inner validity. Life
in time comes at last to stand either as an illusion or as a reality of vanishing significance.

And with this it also hangs together that religion has an inclination to start from too narrow a
concept of value, or to start from the assumption that the sum of values has been found once and
for all, so that it is now only a matter of preserving it, or at least of living in faith in its
preservation. So long as new experiences come, new values can come too, and these religion will
only after a certain resistance take up into the sum of that in whose persistence it believes.

In another way too the incompletability of experience hinders the carrying through of the
religious axiom. There will always be new experiences before which this axiom must stand its
test. Even if we imagine the sum of values concluded once and for all, this concluded sum will
always again be able to be subject to vicissitude, and the religious ideas that were able to
express its persistence during the change of earlier experiences do not always suffice without
more ado when the experiences acquire a different character. The religious axiom here shares the
fate of all axioms.

Finally it is here too of great influence that the religious consciousness expresses itself by the
help of more or less figurative ideas. It is only in philosophy-of-religion reflection that we set
up the proposition in the form in which we here discuss it. In the religious consciousness itself
it appears in concrete, figurative forms, drawn from analogies with more or less limited relations.
A glance at the different ideas of God, theories of incarnation, theories of atonement, and ideas
of immortality shows that they are due to combinations of images drawn from special domains of
experience (family and legal relations, for example). It is then no wonder that when these images
are raised out beyond their original connection and thereafter combined with elements from a wholly
different connection, no consistent line of thought can arise. In the attempt to erect a structure
of thought here, the images' original poetic significance and effect are lost, without any
compensation for it being reached in thought's consistency.---

The hypothesis that faith in the persistence of value is the essential thing in all religion has
then stood the test as well as could be expected.

Indirectly this hypothesis was prepared in the epistemological philosophy of religion, which showed
that religion's significance could not be that of giving a scientific knowledge of existence. If
religious ideas were to have any value, it must then be that of serving as figurative form and
expression for other sides of the life of the soul than the intellectual.

Positively the hypothesis was grounded by the psychological investigation of religious experience
and religious faith. It showed itself that the relation between value and reality was the domain
where religious experience belonged, in distinction from other experiences that either merely
concern values or merely concern reality. And it showed itself that religious faith by its nature
seeks to hold fast a constant relation between value and reality; it is itself faithfulness, and it
presupposes something analogous to faithfulness in existence. There is then an inner connection
between religious experience and faith and the assumption of the persistence of value.

This purely psychological grounding now seems to be confirmed by an analysis of the two historical
chief forms of religion.

\begin{center}---\end{center}

\subsection{General Philosophical Discussion of the Proposition
  of the Persistence of Value}
\label{ssec:value-philosophical}

\noindent\textbf{71.} If faith in the persistence of value is the essential thing in all
religion, first and foremost in all positive religion, it acquires a great interest to
investigate how far such a faith is compatible with a strictly scientific conception of the
world. A distinction must be drawn between its being \emph{compatible with} scientific inquiry
and its being able to be \emph{derived from} this. It is, as I have sought to show, not the pure
need for knowledge, but the need to obtain harmony between knowledge (understanding in the
scientific sense) and evaluation, thus between reality and value, that can lead to faith in the
persistence of value. Although the motives of knowledge and of evaluation are different, their
effects may very well be compatible. If we could not and dared not espouse any assumptions about
existence other than those that scientific inquiry can find out and prove, the proposition of the
persistence of value would have to fall. Science cannot out of itself produce religious faith.
Science works its way towards the assumption that all changes in existence are conversions from
one form of existence to another, conversions that take place in definite relations of measure.
In this direction scientific understanding advances; ever better the new is known as a form of
the old; the new is inferred from the old; and to every new phenomenon a preceding one is pointed
out, to which it corresponds as effect to cause. The eye is opened more and more to a great
continuity in existence. But even if we were to imagine this ideal of science completely realized,
absolutely nothing would thereby be decided with respect to the continuity of value. Existence's
processes could be continuous without there being continuity as regards their value. The valuable
could be a rare guest, to be found at a few stations in the great process of development, but
vanished without trace at all other places. If one thinks that existence's inmost essence is
exhausted by the content of experience being shaped according to science's fundamental
propositions, there is no place for a faith. But such an opinion cannot be proved. The possibility
stands open that the great causal connection which more and more discloses itself to science is
frame or foundation for the unfolding of a valuable content, an unfolding that takes place
precisely through the laws and forms that scientific inquiry has pointed out. The proposition of
the persistence of value need say nothing other and nothing more than this.

\medskip

\noindent\textbf{72.} If the proposition of the persistence of value thus stands as different
from the scientific fundamental propositions (without therefore being incompatible with them),
then in its relation to experience it presents a certain analogy to them. The scientific
fundamental propositions cannot be strictly proved. They stand as fundamental hypotheses, as
principles that guide our inquiry and search by indicating to us how we are to ask, and how we are
to shape our problems. Indeed, the very fact that we ask at all---that we set up problems at
all---is due to a more or less conscious acknowledgement of the scientific fundamental
propositions. The causal principle, for example, has the great significance that by virtue of it
we ask, in the face of every event, which other events form the presupposition of this event's
occurrence. If the causal principle did not lie involuntarily at the foundation in our thinking, we
should content ourselves with noting and describing everything that had happened---drawing up a
catalogue of existence's changes. When the causal principle in science appears in particular as
what we have called the principle of natural causes, this has the significance that we do not think
we have understood an event before we have found its unavoidable connection with earlier events
that are just as sure and indubitable experiences as the new event is. In a similar manner it
stands also with special propositions, as for example the proposition of the conservation of
energy. By virtue of this proposition the investigator of nature asks, at every new expression of
energy, to which other kind of energy it forms an equivalent, and what quantum of this other kind
corresponds to a definite quantum of the new kind. Our scientific fundamental propositions are
hypotheses that are at the same time \emph{anticipations:} we anticipate nature's course in our
questions and conjectures, and our knowledge advances through the testing and confirmation
(verification) of these conjectures. One philosophical school (the aprioristic, idealistic) lays
more weight on the anticipation; another (the empirical, realistic) lays more weight on the
verification; but in themselves these two sides of our knowledge do not stand in conflict with one
another.

As regards now the verification of the fundamental propositions, of the fundamental anticipations,
it will never be able to become complete. It consists in this, that all new experiences show
themselves on closer investigation to satisfy the expectations we cherish by virtue of those
fundamental propositions. But so long as experience is continued, this investigation will not be
able to be concluded. And this becomes of so much the more penetrating significance because the
firm, causally determined connection between the changes is the last criterion we have for
distinguishing between dream and reality. When we doubt the reality of something that shows itself
to us, we test whether we can bring it into unavoidable connection with other things whose reality
stands fast for us; if this shows itself impossible, we regard what showed itself to us as illusion
or hallucination, unless doubt arises about those other things we regarded as real; in this last
case we must test the other things by their connection with things about whose reality we have not
(yet) reason to doubt---and so forth. Since the criterion of reality can constantly only be carried
through approximately, the very concept of reality is properly an ideal concept. We operate so
briskly with it, as with various other concepts, in spite of their unfinished character. Existence
does not lie concluded before us---and who knows whether it is concluded at all? Perhaps it is in a
constant becoming, so that it is no accident that new content of experience constantly wells forth.
Thereby it will be understood that existence is always incommensurable with our knowledge, however
great the advances this makes.

As it stands with the fundamental propositions for our understanding of existence, so it stands also
with the principles that lie at the foundation of our ethical judgments. The ethical principles state
the rules we must follow in our evaluation of human actions if we would remain in agreement with
what is the inmost, the central thing in our nature. They are therefore altered when the central
thing in us is altered, and their application and confirmation in experience always becomes
imperfect, since the great connection in which human actions arise, and in which they intervene, is
unsurveyable for us and is placed in constant development.---From an ethical point of view
existence's unconcluded character acquires a quite distinctive significance. For if existence were
in itself finished, there would be no place for and no ground for an ethics. Of ethical striving
there can only be a question when the world's course is unconsummated, and when its continuation
goes in part through human willing and acting. One of the reasons why religion and ethics so easily
come into conflict with one another lies precisely in the fact that religion is inclined to assume a
finished, once and for all absolutely perfect principle for existence.

In analogy with the fundamental propositions of science and of ethics, the religious consciousness
will be able to set up the fundamental proposition of the persistence of value as an
anticipation---although the verification will be far more difficult and imperfect than in the
theoretical and the ethical domains. The motive for this setting up will essentially be a practical,
personal need that is called forth by experience of the relation between value and reality. The
proposition will express the expectation that amid the changes of the times and the course of
alterations a valuable kernel is preserved in existence, and this expectation will in turn give rise
to a striving to find, in the new forms and phenomena in which existence appears, new garments for
the old values. It is an expectation and a striving that need not let itself be bound by any dogmatic
authority, just as little as it comes into conflict with any scientific presupposition or needs to
lead to conflict with any scientific result.

In all domains life anticipates experience. Life begins with a surplus of strength. Life must begin
with a great investment of strength in order that the resistance to its unfolding and assertion may
be overcome. This holds for the spiritual life as well as for the purely physical. With this hangs
together the involuntary anticipation, the bold formation of ideas, that is characteristic of the
life of thought. We live in expectation before we live in memory; we look forward before we look
back. We are born in faith and begin with confident anticipation. It is then experience's business,
through confirmations and disappointments, to dam in and determine the original anticipations. When
there is conflict between the ideas we have formed and experience, it may come either from our
experience being too limited, or from our ideas being incorrect. The discovery of such a conflict
can, in all domains, lead to great crises. What matters then is whether there is a spiritual
elasticity that can either extend experience so that it fits the ideal anticipations, or transform
the ideas so that they agree with the testimony of unfailing experience. That is life's great art.

The proposition of the persistence of value will often only be able to be held fast by the values
themselves being altered, or by their appearing in other forms than before. Faith in the persistence
of value can therefore only maintain itself on the presupposition of an uninterrupted spiritual work
that is able to distinguish between kernel and shell, not merely so long as it is old, familiar
fruits that existence offers, but also when it bears new fruits. It may happen that what was at
first shell later becomes kernel, and conversely.

But if experience here cannot lead further than to limitation and closer determination, can it not
lead to a complete cessation of the need that leads to setting up the proposition of the persistence
of value? Let us imagine that the need to find events' causes were always disappointed, the causal
principle never receiving any confirmation whatever: would it not then at last fall away? All
spiritual requirements stir only under certain conditions and fall away with these. Every personal
need must struggle to persist, as all life on the whole must. This possibility cannot then be denied
as regards the religious need. And the history of religion teaches that a faith properly dies only
from lack of conditions of life, not because it is strictly theoretically disproved. As Auguste Comte
has said: Apollo and Minerva have never been refuted. It is the whole soil on which the need to hold
fast these figures of gods alone could thrive that has changed. The gods die from lack of
nourishment, and the nourishment comes from the spirit's living need. But the experience of the
history of religion has hitherto shown us only a falling away of special forms or objects of
religious faith, not a drying up of the very living spring from which faith sprang.

Whether such a drying up of this spring will happen in a distant future cannot be known. For the
present there is no reason to assume it. The religious crisis in which we are placed may be a process
of transformation; it need not be a death-struggle.

\medskip

\noindent\textbf{73.} As so often happens in similar cases, this crisis is prepared by changes that
severally do not seem to be of significance in principle. It has in the course of the last centuries
become clearer and clearer that religion's significance cannot consist in the scientific
understanding of existence that it might be able to afford. Negatively this is grounded by the fact
that science leads to understanding by other roads and in other forms than religion. Positively it is
grounded by the growing psychological insight into religion's essence, which leads to emphasizing the
elements of feeling and will as the essential thing in religion, while the significance of ideas is
comparatively subordinate. The conclusion of this whole series of changes can only be the assumption
of all religious ideas' symbolic or poetic character.

If now the need to hold fast the persistence of value maintains itself after the poetic character of
all mythical and dogmatic ideas has been recognized, then continuity in human spiritual development
can be preserved at a very essential point. It is not a philosophy that thereby comes to step into
religion's place. It is a kind of faith that steps into the place of other kinds of faith. Faith can
always only be an \emph{object} for philosophy, not a \emph{product} of philosophy, except in so far
as philosophy can point out the psychological possibility of a certain faith under certain spiritual
conditions of life. The philosophy of religion investigates the epistemological, psychological, and
ethical conditions to which every kind of faith is subject. But it can construct no faith; it can
only describe, analyse, and evaluate the faith that life itself at different standpoints leads one to
develop. The philosophy of religion is a comparative science; it investigates the standpoints and
forms of faith in their relation to the spiritual conditions of life and in their struggle with
these, just as comparative biology investigates organic beings in their relation to and their struggle
with the material conditions of life. And it finds---so far as I can see---the possibility and the
progressive development of a religious standpoint without myth, dogma, and cult, a fruit of personal
experiences that procure expression in a poetry of life which shapes itself differently according to
the individual's distinctive character. Even if images and symbols from the classical times of the
positive religions are retained, this happens at such a standpoint with historical and psychological
justification. In the first place, such reinterpretation as is here necessary is constantly undertaken
within the positive religions. And in the second place, the great poetry that is often laid down in
myths and dogmas is not the private property of the individual people, the individual church or sect,
but belongs to humanity and stands at the disposal of everyone whose experience of life goes in the
same or an analogous direction as the experience to which those images owe their first
coming-into-being or shaping. But there are those in whom the image- and symbol-forming process takes
place at first hand, whether its fruits can acquire significance for the individual alone or for
several along with him.

\medskip

\noindent\textbf{74.} The expression \emph{positive religion,} which has often been used in what
precedes, needs a closer explanation. It is probably to be understood on the analogy of the expression
positive law. What is meant by it is then a religion that has been shaped in definite traditions, in
common forms and customs, and it thus presupposes a community with fixed norms. Positive religion
stands in opposition to the ``natural'' religion at which the single individual can arrive by the road
of his own conviction, just as positive law stands in opposition to natural law, which is derived by
reflection on the conditions of communal life and on the individual's relation to the community. The
word positive then means what is laid down (\textit{quod positum est}) in definite forms.

The positive religions that history shows us all have the distinctive feature that they have their
centre of gravity in a cult, and that their content consists of myths, legends, and dogmas that are
regarded as valid expressions of absolute truth. And it is a general feature of these myths and dogmas
that personal beings (one or more) are regarded as the author of the world, or at least of the most
remarkable things that happen in the world. Purely historically regarded it would then be justified to
take up this distinctive feature into the definition of positive religion. In my \emph{Ethics} (ch.
XXXII) I have also started from this historical definition. But the possibility cannot (as I have also
remarked in the work cited) be dismissed that a religious community might be formed whose faith found
a poetic, symbolic expression, without any dogmatic conclusion. There may very well develop a common
symbolic cult without magical character and a common symbolic circle of images, or one may unite about
common great experiences of life whose special shaping is left completely free to individuals, so that
it depends on the individual personality's originality and productive capacity how far his symbols are
to acquire significance in wider circles. It would be idle to attempt a closer construction of such a
standpoint; if it comes, life produces it by its own roads, just as it has produced the old positive
religions. The concept of positive religion would then acquire a greater extent (and thereby a smaller
content) than one would give it on a purely historical consideration.

The positive religions that history knows can psychologically be conceived as concentrations, syntheses,
in which human beings' ideas about existence have been worked together, interwoven with their deepest
experiences of life. No proof can be given that the time of such spiritual concentrations is past, even
if we for the present live in the time of criticism and analysis. And neither can any proof be given
that such concentrations must necessarily have the mythological or dogmatic character that the hitherto
known positive religions have had. Spiritual freedom and strength are needed to live on the conviction
that our highest ideas can only be similitudes. Conditions have hitherto been less favourable for the
development of this capacity; but who dares set a limit to the freedom and strength that will be able
to develop in the human race?---

It appears to me therefore unjustified to narrow the concept of positive religion in such a way that
the mythical and dogmatic personifications come to stand as a necessary appurtenance. And it would be
especially unjustified if our earlier psychological investigation, which led to acknowledging the
derived significance of ideas in the religious domain, is correct. Besides, ``natural'' religion has as
a rule also had that appurtenance. But still more unjustified would it be if one wished to regard such
personifications as necessary constituents of all religion, whether it appears as ``positive'' or not.
One will not be able to dismiss the question: what value is there in the personal beings one believes
in existing, and in one's believing in them? And then the idea of the persistence of value will show
itself as the fundamental religious idea, the one that also first gives the ideas of the personal gods
a real content and a real significance. And this idea will be able to preserve its significance even if
all personifications acquire only symbolic significance, and even if other symbols could be found than
those drawn from human life and the human figure.\textsuperscript{85} There will then be preserved a
continuity in human spiritual life, even where this leaves the positively religious standpoint. Whether
the standpoint reached when positive religion is left is really a higher one is a matter by itself. It
does not become a higher one without more ado because myths, legends, and dogmas are pushed aside. As I
have once said in another place: \emph{not} to believe something is about the least that can be said
about a human being. The negative---giving up a standpoint---says nothing about the value of the new
standpoint, or about whether a new standpoint is reached at all. The standard for ``lower'' and
``higher'' at the transition from positive religion to a free religious standpoint must be the same as
is applied within the series of religions when the one religion is said to be higher than another. In
the ethical philosophy of religion we shall come back to this point.

As I have several times remarked, I do not lay great weight on the name religion. The danger in using
the word religion in a more comprehensive sense might be that dishonest accommodation would thereby be
favoured. But on the other hand religious life in the present often assumes---and not least in such
circles as think they have an exclusive right to be called religious---forms of such a character that
it truly does not require great resignation to renounce the use of the name religion.

But however one may name a standpoint that holds fast the faith in the persistence of value without
clothing it in dogmatic form, it is a task for the philosophy of religion to investigate more closely
how far such a standpoint is possible, and what difficulties it has to struggle with.

\medskip

\noindent\textbf{75.} All evaluation takes place---like all understanding and explanation---from
the human standpoint in existence and is determined by the human conditions of life. But this
standpoint and these conditions are limited, and their inner connection with the rest of existence
cannot be surveyed by us. Only if we imagined humanity and its globe as the central thing in
existence would this difficulty not exist; but the Copernican conception of the world has extended
the horizon and thereby made the problem of the persistence of value far more complicated than
before. The unsurveyable connection into which the history of our highest values is woven makes it
not merely impossible to give an objective grounding of the proposition of the persistence of value,
but also makes it impossible to give the faith in the persistence of value a definite, let alone a
vivid form. Beyond the general thought that what has real value stands in so intimate a relation to
the forces that stir in existence that it will in one form or another continue to persist, we cannot
get.---This is a difficulty that is common to all religion. The positive religions nevertheless as a
rule start from the assumption that existence is limited, and none of them has given any solution of
the doubts that arise when the limitation falls away.

The extension of the horizon must at any rate lead to the conviction that our concepts of value
cannot be conclusive. A highest value can as little be pointed out as a first cause. We are often
compelled to halt at causes for which we cannot in turn find causes, and thus we are also compelled
in our evaluation to halt at values that for us are the highest, and which must therefore be those
that directly or indirectly lie at the foundation of all other values we can find. But if in our
faith we venture to ascribe significance to the concept of value beyond human existence---which is
only possible by the road of analogy\textsuperscript{86}---then we must also familiarize ourselves
with the thought that \emph{our} highest values can themselves only be elements in a still more
comprehensive order of things, a kingdom of values whose fundamental laws and whose extent we can as
little represent to ourselves as we can represent to ourselves a comprehensive order of nature of
which the conformity to natural law that our science discloses to us should be only a fragment.
Faith in the persistence of value therefore presupposes a bold resignation; only thereby can this
faith hold itself up even there where our capacity to evaluate has reached its limit, because our
experience has reached its limit.

But even if we do not go beyond the purely human domain, it is easy to see that we cannot form
definite and vivid ideas of value's future forms. There take place, namely, as a matter of
experience, constant displacements of value, as has already been mentioned in an earlier connection.
What at first has value only as a means can later acquire value as an end, and what at first has
value as an end can later show itself to be an introduction and preparation for new, hitherto unknown
and more comprehensive values. There can thus take place subjective displacements of value, in that
the motive of evaluation is altered, and objective displacements of value, in that wholly new values
come into being. And such displacements are the only forms in which experience establishes a
persistence of values. They are the real content of human history. But they show then at the same
time at once that it is not the individual, definite value that persists. In order that value may
persist, it must be altered, just as the germ must be altered in order to become a plant. There holds
in the domain of values an analogy to what anatomy calls epigenesis, in contrast to preformation:
just as during growth such thoroughgoing changes take place that the germ is not merely a plant or an
animal on a reduced scale, so neither can we without more ado find the later values preformed in the
earlier ones. Only experience can show what new forms the values will assume, and this experience
extends far beyond what single individuals and generations will be able to survey.---Another
scientific analogy also lies near. Just as little as one can derive from the general thought of the
conservation of energy the special conversions of energy that actually take place, just as little can
we derive from the idea of the persistence of value what new forms the values will be clothed in. In
both places it is experience that alone can instruct us about the new forms' character, and experience
is never concluded.

If we were not sufficiently warned beforehand against a dogmatic formulation of the proposition of the
persistence of value, the law of displacement will contain an emphatic warning. It is not the definite
empirical values in whose persistence we can believe.

We must even go a step further. It follows from the psychological law of relation, according to which
every feeling is determined in its character (quality and degree of strength) by a certain opposition
to the preceding and simultaneous states and elements, that every value is altered when the relations
are altered, and that an unchangeable value would be no value. When the background of the feeling is
altered, the feeling too must be altered; and when the background is uniform and unchangeable, the
feeling will be weakened and at last fall away. Value can only persist through alterations and
conversions. This hangs together with the reality of the temporal relation and of differences
generally. Only by a purely mystical road---a road that consistently ends in ecstasy---will one be able
(sometimes) to attain to disregarding this order of things. It will nevertheless easily be seen that
the law of relation entails that the alterations and differences themselves acquire their great value
when they are conditions of value's arising. The religions have generally been inclined to overlook
this, even when they acknowledged the significance of the temporal relation. For the sources and
conditions of value which precisely changeability and differences contain, they have had but slight
eye. With a certain shyness they have looked away from this, and only by a flight from this so
essential side of reality have they thought to be able to assert the persistence of the valuable. It is
a point where there are still significant experiences to be made and new spiritual lands to be
discovered.

The confusion of individual definite values with eternal values is irreligious. But it is an
irreligiousness of which most religions make themselves guilty. Their religious postulate then acquires
the form: ``If the kind and form of worth that I know does not persist, then I do not care about the
persistence of value---or rather, then I do not acknowledge that what persists is or has value.'' It is
an egoistic form of religiosity, which is far from rare. Faith in personal immortality in particular is
often grounded in this way. As if there could be no meaning in existence even if I were not to be
immortal! Of the individual personality's relation to the great kingdom of values no clear idea can be
formed, and neither affirmation nor denial can here be grounded. The special idea that the individual
forms for himself on this question will always be assignable to the domain of poetry and symbolism as
soon as it goes beyond the general thought that what has real value will persist in one form or another.
Such a poetic clothing may (whether it is a positive or a negative decision of the question of personal
immortality that puts on this clothing) have its great significance and be indispensable; but the
clothes do not make the man, even if the man cannot do without clothes. Towards the life after death the
evangelical admonition ``Take no thought for the morrow!'' is to be applied with greater right than
towards our relation to the next day within earthly life. Whether it is precisely the forms of value
that we know that persist, experience will show.

\medskip

\noindent\textbf{76.} The preceding consideration has moved within the concept of value. But the
question cannot be held back whether the whole difference we make, and must make, between value and
reality holds beyond the human standpoint---whether it is not precisely a purely human difference. This
difference might indeed hang together with the fact that neither our concept of value nor our concept of
reality is concluded and concludable. Our concept of value is empirical, and our concept of reality is
ideal. And we can derive neither reality from value nor value from reality. None of the ways in which
dogmatics and speculation have tried to undertake such a derivation can stand a closer test. Our thought
is then led to the possibility of a principle from which at once value's arising and persistence and
reality's connection might spring. But it is a thought that cannot be carried through in scientific form.

Meanwhile the mere possibility of value's persistence beyond the domain of human life seems to
presuppose the constant persistence of life of the soul (whether this in turn entails the persistence of
the individual ensouled individuals or not). For value presupposes a relation to life of the soul. If we
knew a form of existence in which the difference between value and reality fell away, there would not lie
so great a problem in this; but when we hold to our experience, we must see how the matter stands when
there is a difference between value and reality, and when within reality there is a difference between
spirit (life of the soul) and matter.

In this respect it is of importance that it is not possible to derive the spiritual from the material.
The criticism of materialism in its different forms has sufficiently established this. Life of the soul
stands as a distinctive, irreducible side of existence, however one may think of its relation to other
sides of existence. In the common ground, inaccessible to us, from which both the spiritual and the
material spring (since the latter can as little be derived from the former as the former from the
latter), there may be a possibility of a continuity beyond the limited extent in which our experience
shows us life of the soul in the world, even if here too the special forms in which this continuity
appears cannot be represented by us. The main thing is that life of the soul stands as an independent
element in existence, not as a superfluous addition or without inner connection with existence's other
elements.---Yet even a materialistic conception of the world could very well hold fast a faith in the
persistence of value, namely if it assumed that matter always contained conditions for the arising of a
life of the soul within which life's earlier course of development could be continued.

But neither must we at this point forget our experience's limitation. We have no guarantee that existence
does not have other sides or attributes than the two---the spiritual and the material---which are the only
ones our experience presents. Neither an idealistic nor a materialistic conception of the world can
therefore be fully carried through. As regards the problem of value, that possibility of the incompleteness
of our knowledge of existence's fundamental forms (attributes) entails that the horizon is extended, in so
far as more roads and ways become possible for the displacements of values and for their persistence
during the displacements than is the case when one starts from the assumption that the opposition between
the spiritual and the material is contradictory, so that there is a forced choice between them. Such an
extended horizon warns both against narrow-minded doubt and against narrow-minded dogmas.

Let me illuminate this point by the help of an analogy. When sulphur changes its state, passes from solid
to liquid and thence again to gaseous state, it also changes colour. But the colour cannot be explained by
the state, any more than conversely. If we now knew nothing else about sulphur than its states and colours,
the problem of the inner connection between state and colour would be insoluble. But now we also know
sulphur's different temperature and know that it is the heating to which both the changes of state and the
changes of colour are due. It is such a third thing that we, as regards existence as a whole, do not know.
Only if we knew such a thing would the problems of spiritual life, and especially the problem of value, be
able to find their solution for us. As it is, the fate of values, as well as of life of the soul generally,
stands for us in a fragmentary form which does not make a faith in the persistence of value impossible, but
brings it about that it cannot become anything other than a faith.

\medskip

\noindent\textbf{77.} When there is talk of the scientific conception of the world, one often lays too
great weight on the general laws of nature, which are to be sure also what strikes the eye most when one
compares the modern picture of the world with the old. But the laws are nevertheless properly abstractions.
They denote the forms and rules according to which movement and development in existence take place. But
the essential thing is nevertheless that movement and development, both in the spiritual and in the
material domain, always have a certain definite direction and thereby a certain definite individual
character. Existence is like a great individuality; it stands for us as one of countless possibilities; we
must take it as we take the individual human being, who is always more than an example of a general law. An
individuality arises through an interplay of laws; it is (to use H. C. Ørsted's word) a
\emph{togetherfoldedness} (in contrast both to ``singlefoldedness'' and to manifoldness), which we must
take as experience gives it. However far the boldest hypotheses go back, we constantly find certain definite
forces operating in certain definite directions. Even if one could derive all later forces and directions
from these---for us---original ones, the original forces and directions nevertheless stand for us as mere
facts that must be taken as facts, just as the original dispositions with which a human life begins must be
taken as factually given in their individual determinateness. Existence is a \textit{unicum}, and if we have
a right to compare it with a drama, then it is a drama with a definite, special exposition and course of
scenes, which must be taken as they once for all are.

The investigator of nature and of history seeks to find his bearings in this drama's course. But since only
a single scene, perhaps even only a single speech, of this great drama falls within our experience, we
cannot construct the whole drama for ourselves. And if we attempted it, the drama we produced would
nevertheless in turn show itself to be only a part of a still greater drama. Every mythology, every
dogmatics, and every metaphysics leaves---however it may go with the inner coherence and consistency of
their constructions---always unanswered questions, partly with respect to what lies before the beginning of
the drama they depict, partly with respect to what follows after its close. Neither the foundation nor the
consequences are exhausted.

A human individuality too is inexhaustible in its distinctiveness. It is impossible for us to find and
connect all its elements so that the ``togetherfoldedness'' comes forth in its special character. But still
more does this hold with respect to existence as a whole, where we do not even, as with a limited
individuality, have the reciprocal action with a surrounding world to learn from. The problem of existence
acquires its sharp sting precisely from the fact that existence, when we try to conceive it as a whole, must
have a special, individual character, and that the individuality we try to ascribe to it cannot stand as
conditioned by the relation to anything different from it, as is otherwise the case with all individual
forms of existence.

For our general theoretical and practical fundamental propositions this acquires its great significance. If
they were to be able to be completely established as holding, existence would have to stand for us as an
individual whole whose inner law those fundamental propositions were. Thus the causal principle would only
be able to be strictly proved if all phenomena could be referred to their definite place within a great
whole; the same holds of the proposition of the conservation of energy, which can only be proved for
concluded and isolated wholes. The unprovability of the fundamental propositions hangs together with
existence's---for us and perhaps in itself---unconcluded character. The proposition of the persistence of
value shares the fate of the other fundamental propositions; only if we could grasp existence in its whole
``togetherfoldedness'' would it be able to become clear what validity it has. And in every attempt to assert
it, it must not be forgotten that every such general proposition means only an incision into or a glimpse of
reality from a single side, not an exhaustive determination. Both existence's individual character and its
infinity here constantly set us limits.

\medskip

\noindent\textbf{78.} The most serious difficulty for holding fast the religious fundamental
proposition arises rather from existence's presenting too many possibilities than from its
offering too few. There takes place in existence a constant conflict and struggle between elements
and individuals. Countless germs are cast out, but only few take root and carry life and
development further. And where development is furthered, it happens under mutual inhibition,
through crowding out and pushing aside. Development takes place from a mass of different starting
points, and the tendencies of development that these starting points possess by no means stand from
the first in harmonious relation to one another. Most frequently a result is reached at one point
only through nothing's being reached at other points. Resistance and hindrance can to be sure have
their great worth by being releasing forces which free an energy that would otherwise slumber; but
it is very far from being the case that the resistance and the hindrance that experience shows us
are limited to this.

It is all development's sporadic character that causes the running costs in existence to seem
greater than the yield. Existence shows itself prodigal and therefore ruthless. And the prodigality
effects, so at least it seems, a loss of value.

The problem of the disharmonious in existence leads back to the problem of unity and multiplicity.
It has shown itself impossible to derive multiplicity from unity, and even if it were possible, the
problem of disharmony would thereby appear only so much the more glaringly: for if multiplicity---the
many sporadic starting points---had really sprung from a principle of unity, how would it then be
possible to explain that there is not always a harmonious, but often a disharmonious relation between
the different tendencies of development in the manifold elements and individuals? There might then
well be conflicts between them, but not greater than the release of energy
required.\textsuperscript{87}

It is of no use here to refer to the myth of a fall; that is only to put several riddles in place of
the one. For how is a fall possible when every element and every individual in existence has
originally sprung from a harmonious principle? And a fall explains at most the disharmony within the
human world, not the other disharmonies in existence.

If one will say that the possibility of a fall had to be there if the human being was to be a
personality, since such a one develops only through choice between possibilities, then one comes into
a serious self-contradiction when one ascribes personality to the divinity and yet will not concede
the possibility of a fall in its case. The very possibility of a fall is moreover already the
expression of a disharmony: if I have the possibility of committing a piece of villainy tomorrow, then
I \emph{am} already a villain today. If one concedes an original possibility of a fall as posited at
the creation, then one has also conceded that the disharmony was produced by the creation itself,
which indeed from another point of view (cf.\ §69) must consistently stand as a diminution of value.

No scholasticism can avert this consequence.

There is a side of the matter that is as a rule overlooked in the discussion of this problem. One is
inclined to start from the assumption that it is existence's totality, the principle of the whole, that
has ``right'' on its side, and that the individual element or individual whose tendency of development
leads to a breach of this totality has the ``wrong.'' But it is a great abstraction when one in this
way places the principle of the whole in opposition to the individual elements or individuals. The
force that stirs in each of these is after all a part of the whole's force, and the laws that work in a
``togetherfolded'' manner in the individual being are individualized interplays of existence's general
laws. Every single element of existence has so far the same ``right'' as the whole, and it is not
rebellious because it follows its own law, its own tendency of development. It is an oriental line of
thought that makes one inclined in advance to find the ground of disharmony in the elements and not
also in the principle of the whole. It was against this line of thought that Goethe reacted in the
magnificent poem ``Prometheus'': it was precisely the inner germ of divinity in the individual finite
being that he wished to assert in opposition to the customary external concept of God. Only from a
dualistic religious standpoint could this poem be regarded as irreligious. The profound religious
conception according to which God himself suffers has (however it may arrange the matter for itself
mythologically or dogmatically) in reality given dualism its death-blow, without itself knowing or
willing it.---But the settling of the whole account naturally becomes so much the more difficult when
regard must be had not merely to the possibility of a total harmony that might arise through the
interplay of the finite harmonies and disharmonies, but also to those regions in existence where
disharmony prevails, and where one cannot hear how it fades into an all-embracing harmony. Such regions
are after all independent parts of existence and cannot be dismissed with the plea that they represent
only subordinate passages in the great world-concert. Only an oriental ``superman'' (of which the
supermen constructed in the most recent times are impotent plagiarisms) could sit in his heaven and
find harmony in his existence under these conditions. Here lies the great justification and the
beneficial effect of Voltaire's and Schopenhauer's protests.

When existence, already from a purely theoretical and intellectual standpoint that asks only about
scientific understanding, can come to stand as irrational or unfinished, this holds no less from the
standpoint of evaluation. It is confirmed here too that religion does not solve the riddles that science
cannot solve. The advantage religion perhaps has here over science is that it can more easily lead
thought away from the insoluble riddles, so that their sting is less felt; this is an advantage where
there is to be willing and acting---where the struggle is to be taken up and courage is to be kept up.---
In order to escape the sting that is in the problem we here stand at, one must---to use an image of
Bayle---prevent thought from seeing its own shadow, just as Bucephalus could only be steered when it
was not frightened by its own shadow on the ground. Consciousness must be occupied with an ideal
content, concentrated on a shining thought, in order not to be checked in its striving by the idea of
existence's dark and disharmonious sides: this is a condition for the very struggle against the dark and
disharmonious being able to be carried on. But to look away from the problems is not to solve them. And
it is then most consistent to concede the insolubility, especially when the case is such that life can
be led greatly and beautifully even if one does not get a solution to all riddles. It is not existence's
whole factual constitution that properly concerns us as striving beings, but the conditions for
continued development that it presents. Naturally the strength of the resistance and the degree of the
disharmony will determine the character of our experience concerning the relation between value and
reality. But the proposition of the persistence of value is in itself independent of how little or how
much value there is in existence; it says only that the value that is there persists. Only if one lets
reality spring from value does the problem's sting become unbearable. If there is a possibility of
asserting and perhaps increasing the value that once for all is found in existence, then the religious
axiom can be upheld.

There are two thoughts that will always be of penetrating significance for our position towards the
problem that the disharmonies in the world cause.

In the first place our experience is limited in relation to unsurveyable existence, and it must then
constantly be tested whether an extension of experience might not entail a different result. In purely
theoretical problems already this thought plays an essential role. When our experience by strict
inference leads to self-contradiction, we do not therefore give up the assumption of the validity of the
logical principles, but test whether the ground of the self-contradiction did not lie in our experience's
being too limited. Often it then shows itself that the contradiction falls away when we can extend
experience. Thus objections to the persistence of value too will often fall away when we extend our
horizon through the discovery or production of values there where hitherto there seemed to be no
possibility of them. It was the most important and just about the only valid thought in Leibniz's
well-known attempt to assert an optimistic world-view, that by reason of the idea of existence's
infinity we are more favourably placed towards the problem of the disharmonies than the thinkers of
antiquity (Plotinus, Augustine).

In the second place life is constantly rejuvenated. Every year has its spring, every generation its
youth. New germs therefore constantly come forth, and there is no trace of the world's having grown old.
A new section of existence's history constantly begins, in relation to which the earlier series of
experiences on which we build our judgments about reality's value is only a prehistory; the whole earlier
drama now becomes only an exposition to a greater drama. Existence is not merely unsurveyable, but
inexhaustible.

A faith in the persistence of value is therefore in spite of everything psychologically possible. Whether
it also really is found, and in what form and what extent, experience must show.

\medskip

\noindent\textbf{79.} To a great extent the living feeling of existence's disharmonies (in us and outside
us) rests on contrast-effect. It is the ideal that judges. The opposition between reality and the great
ends we can set ourselves conditions the disparaging judgments that are passed on reality as it is. The
young person who for the first time meets reality's sluggish and defiant resistance, without being able to
acknowledge the value that even such resistance can have, asks in astonishment: is this what it is to
exist? And the greater the demands one comes with, not on one's own behalf alone but on behalf of the
highest one knows, the greater the darkness one discovers round about one. One might very easily escape
discovering the disharmonies if one limited and dulled one's sense for harmony and ideality. Much pain is
avoided by him who does not attach his desire to great things and to comprehensive circles of interests.

It belongs then to the human being's nobility that he can grieve and suffer. It is a sign of what
oppositions the human being can contain. It testifies to the deep, powerful unity in human nature. Pain is
in itself a symptom of dissolution, and only that can be dissolved which has unity in its being; pain is
the feeling that this unity is threatened. And when spiritual life can now be held together and be led in
concentrated form in spite of the pain, it shows itself clearly what energy it possesses.

The very fact that the disharmony is felt is then a sign that there are values in existence. Pessimism has
an optimistic element as its background. The contrast-effect can to be sure become so strong that one
wholly loses the sense for the background's value; but that will only happen when one sinks away into
purely passive states.

The last position that can be taken up towards the problem we have here occupied ourselves with lies on the
boundary between religion and ethics. Value is not absolutely dependent on its own persistence. What we
value as beautiful and good can keep its value whatever its fate may be, and whatever shadows may fall on
its way. A thought does not necessarily become less true, a feeling less pure and noble, because they must
pay their tribute to time. Value is not always measured by the clock. And if it should be the lot of the
good and the beautiful to perish, would they therefore be less good and beautiful? The richer in content
life is, the more we forget time. In the very work of finding and producing values there lies a faith in
these values' persistence, in so far as this work presupposes the possibility that the values can be
released and find their place in existence, whatever and however great this may be. If faith in the
persistence of value is the kernel of all religion, then in a standpoint that is wholly absorbed in finding
and producing values there is a hidden religion present. One can \emph{be} religious, be in a state of
religion, without \emph{having} a religion; and conversely the religion one only \emph{has} stands in a
very loose relation to the personality. If one will define God as the principle of the persistence of value
in existence, then everyone who (however it may stand with his catechism) works to assert value in existence
is a ``child of God.''

\medskip

\noindent\textbf{80.} It is the task of the humanities to understand and to evaluate spiritual phenomena,
not by holding to their fully developed form alone, but by seeking back to the forces from which they
proceed, and thereupon investigating whether conditions are present for these forces to be able to continue
working under new conditions and in new forms. Just as in outer nature there is an inner connection amid the
change of the outer forms, so it will also show itself in the spiritual domain, the more unprejudicedly we
immerse ourselves in it. If the preceding investigation has led to a correct result, ideas in the religious
domain belong to the forms that change, not to the kernel itself, although they naturally, so long as they
correspond to the whole spiritual stage of life, exercise a great influence on the kernel's condition.
Religious life presents, according to the result I have reached, a greater continuity than the religious
ideas. As an attempt at understanding existence, at solving the riddles of the world, religion has lost the
battle. But it lives in human feeling and need, and from this it will always anew be able to set the will in
motion for the finding out and producing of values, just as with the imagination's help it will be able to
form images and symbols in which the highest poetry of life is expressed.

But the problem of the persistence of value here returns again in a more special form. Will there not be a
loss of value bound up with the transition from dogma to symbol? Spiritual life must after all be able to
have a quite different closedness and concentration when the highest ideas it knows---or rather: the ideas of
the highest that it knows---appear under vivid forms in whose reality and validity there is unconditional
belief, and which can at the same time be common to all individuals, or at least to wide circles. The many
reservations and critical misgivings that lead over from dogmatism's standpoint to that of the poetry of life
seem as though they must weaken the concentration that life cannot do without. So essential an alteration in
the significance of ideas cannot fail to intervene deeply in the life of feeling and will.

There is much in modern spiritual development that testifies to the justification of this question. Not for
nothing has the figure of Hamlet become a type in literature right up to the most recent time, even if the
dimensions have gradually become rather small, and even if the painful process of dissolution in which the
Shakespearean hero is placed on the boundary between instinct and reflection has gradually given place to a
self-satisfied decadence that regards the foam as the most valuable product of the surge. The foam has its
beauty, but true culture is not merely foam but wave, not merely effect but cause.

The first remark that must be made here is that if there is here a dissolution, it begins already within
positive religion itself as soon as its classical times are past. Within the Church already, religious life
has lost its unity. Successively one has had to concede that religion does not teach us astronomy, physics,
natural history. Biblical criticism (or, by a more correct name, the literary history of the biblical
writings) leads ever more clearly to seeing that the Bible does not teach us history either. And in practice
it is conceded---for example by the theological reinterpretations of the Sermon on the Mount's clear
commandments---that it does not teach us ethics either. The serpent has already here come into paradise. The
split is proclaimed, and the problem is set for everyone who will think. Religion, which in its classical
times was everything, now tends towards becoming the last consolation. In its classical form it can
now---by those who see clearly\textsuperscript{88}---only be held fast by a convulsive exertion which shows
that concentration upon it is not so natural as formerly.---The Church will have to go through a course in
the use of symbolic ideas if it does not succeed in defeating modern biblical literary history. The figures of
the patriarchs will pass over from history to myth or legend, and the task will become to hold fast their
religious significance in this new form. They are small lessons that the Church takes on at a time. It was two
hundred years in accepting natural science's independence. Time will show how long it will be before it
accepts the independence of scientific historical inquiry. And after that comes the turn of psychology and
ethics.

Immediate security is abolished as soon as independent seeking comes to play an essential role. Protestantism
already denotes so far---in its beginning---a dissolution; later it has been able to acquire a factual
security. Only where religion encloses the human being on all sides as an objective power that does not depend
on the human being's conviction, only there is there full security. Where self-activity is clearly traced, and
where one's own choice expressly makes itself felt, full security cannot be attained. There is therefore
something justified in the view that religion is something one cannot oneself acquire or produce. The Russian
peasant finds it in order that his Tartar fellow-townsmen have a different religion from himself, but he takes
offence at the sectarians within his own religion, and the motivation is this: ``The Tartar peasants have got
their religion from God, just as they have got the colour of their skin from him; but the Molokans (the Russian
Protestants), they are Russians who have themselves thought up their faith.''\textsuperscript{89}

---It is easier for human beings to find God in fixed, transmitted arrangements than in the inner seeking and
yearning that can lead out beyond the outer historical connection, and in which the individual personality
seems, to an external way of regarding it, to be quite left to itself.

A self-acquired view of life does not so easily acquire the firmness that a transmitted faith can have, in
which the most important work has long since been done, and in which centuries' accumulated experiences of
life are laid down in images whose execution and application have been accommodated to the generations'
spiritual needs. With becoming and new formations there follow unrest and oscillation on account of the new
forms' unfinishedness and untestedness. Only gradually can inward absorption in the new spiritual conditions
remedy the dizziness that easily follows on departure from the narrow but secure forms in which life was
formerly led.

In many people---by one of those displacements that play so great a role in the life of the soul---criticism
and negation become the main thing, not merely a means of liberation. The nutcracker becomes more important to
them than the kernel it helps to get out. An ironic, blasé attitude becomes the consequence and is perhaps
regarded as a sign of the most advanced standpoint. It is not from such a form---one of philistinism's many
forms---that a transition to a new stage of life becomes possible. Development has there run aground, and the
really new must come from a wholly different quarter. But then neither is it of any use to cling to the old
forms when life is departing from them. One then becomes homeless in the new world that is being built up round
about. From fear of becoming a Hamlet one becomes a Don Quixote.

\medskip

\noindent\textbf{81.} There is no lack of models of an art of life and of thought which through the purgatory
of the problems is able to press forward to the free human being's standpoint and realize this in harmonious
conduct of life. Spinoza and Goethe, Fichte and Stuart Mill are such models, and there is no reason to believe
that their number should not increase. Even in spirits who were not able to reach such an energetic harmony
there may be found references and directions to a new ideal of personality in the direction in which those
older models already point. Both from Carlyle, Eugen Dühring, and Friedrich Nietzsche there is much to learn.
Even where one's own capacity does not suffice on account of unfavourable inner and outer conditions, the need
can stir and give its testimony to life's direction.

There is no reason to doubt that new forms of harmonious concentration will be able to be reached when the
after-effects and side-effects of dogmatic narrowness and critical blaséness have disappeared. The imagination,
which formerly unfolded so vigorously within the mythical and dogmatic frames of the positive religions, will
be able with freedom to create its own forms when it is once again pressingly experienced what a wonderful
thing life is---what oppositions it contains, and what struggle it demands.

The great models from the past are not lost because their tracks can no longer guide us on our ways. It is once
for all so, that what was able wholly to fill an earlier generation's life becomes for later generations only an
element in a new whole; but it is precisely one of the most important forms by which the persistence of value
becomes possible, that what was first a whole later becomes an element. In this way the Christian view of life
has made use of Jewish, oriental, and Hellenic spiritual life, and thus a new view of life will make use of
Christianity. The bond of unity for the new view is not finished; but why not hope that it is being woven in
silence? Every serious striving can make its contribution to it.

\begin{center}---\end{center}

\section{The Principle of Personality}
\label{sec:personality}

\subsection{The Significance and Justification of the
  Principle of Personality}
\label{ssec:pers-significance}

\begin{flushright}
\textit{La seule chose universelle doit être l'entière liberté}\\
\textit{donnée aux individus de se représenter à leur manière}\\
\textit{l'éternel énigme.}\\[2pt]
\textsc{Guyau}.\\[4pt]
\textit{(The only universal thing must be the entire freedom given to}\\
\textit{individuals to represent the eternal enigma in their own way.)}
\end{flushright}

\bigskip

\noindent\textbf{82.} On the boundary between the psychological and the ethical part of the
philosophy of religion lies---as regards grounding and significance---the principle of
personality, by which we here in the philosophy of religion understand the fundamental
proposition of the justification and value of personal differences in the religious domain. The
principle of personality has also its great ethical and legal significance, and there is a certain
analogy in its grounding and significance in the different places. In ``\emph{Personlighetsprincipen
i filosofin}'' (Stockholm 1912) I have sought to point out its application in the different
domains. Its necessity in the religious domain springs partly from presuppositions that also find
application in ethics, partly from presuppositions that are distinctive of religion.

A personal being ought never to be treated merely as a means, but must always at the same time,
and first and foremost, be regarded as an end. This has its ground in the fact that personal beings
within our experience stand as existence's centres of value, that is to say, as the living
mid-points where value can be felt and acknowledged. It is personality that, in the world of our
experience, gives everything else value. An inhibition of a centre of value must necessarily entail
a diminution of value. Therefore every inhibition---every compulsion and suffering---can only have
justification as an educative means, a means whose time must one day be past. All authority
inhibits, compels, or pains. Therefore authority can always only be a means, and the principle of
authority must be subordinated to the principle of personality, just as mediate value must always
be subordinated to immediate. The burden of proof must always rest on him who would inhibit, limit,
compel, and pain.\textsuperscript{90} Authority draws its justification from being a condition of
the full carrying through of the principle of personality; but personality does not draw its
justification from any authority. Positive religion, as it has hitherto appeared, always has a
greater or lesser tendency to make authority absolute. In Catholicism especially, the principle of
authority appears with full consistency as an absolute principle. Authority here sets its own limits,
decides what belongs directly or indirectly under religion and what does not---establishes, for
example, that the proofs of God's existence are scientifically valid. Complete intellectual
submission to the Church is demanded, and it is a breach with the Church when one attempts to instruct
it instead of letting oneself be instructed by it. There arises here an irreconcilable conflict
between the principle of authority and the principle of personality.

An inhibition is exercised when the individual personality does not come to work according to its full
distinctiveness. Where religious faith is predominantly due to tradition, imitation, or habit, there
will easily be forces and elements in the individual's nature that do not come to work along with the
rest. Only where he self-actively appropriates and develops his religious conviction can he be present
in it with his whole nature. The ideal must therefore be the greatest possible self-activity, and all
tradition and authority ought only to be awakening, guiding, and educative. All more deeply penetrating
religious tendencies have acknowledged this principle in greater or lesser extent. It is the principle
of inwardization. A misapprehension of it violates the very principle of the persistence of value and
is therefore irreligious. The greatest possible scope must be demanded for individual differences in
the religious domain. It is uniformity and schematism, not differences and distinctive features, whose
justification is to be established.

In this grounding, philosophy of religion and ethics meet. But there is also a grounding that is
special to the philosophy of religion, and that builds on the religious-psychological investigation
undertaken in what precedes.

Individual personalities are not merely centres of value, but centres of experience. It is in them that
the relation between value and reality is experienced, and it is such experience that forms the
foundation of all religion. It is then of the greatest importance that this experience become as
comprehensive and unprejudiced as possible. But it can only become so when individuals are given the
right to work with all their capacities and with all sides of their being. There are many values, and
reality is unsurveyable, not merely in extent, but in the manifoldness of its qualitative relations.
Each individual is, within the world of values as well as within the world of reality, placed
differently from others. He must take his section, seize his nuance, and the more he is himself, the
more he comes to make experiences that no one else can make. The art of life must therefore be
one-sided, just as Julius Lange has shown that pictorial art must be. ``A work of art,''
says Lange,\textsuperscript{91} ``cannot convey the all-sided conception of the human figure; only
history, through its whole course in all its stages, achieves that. Art is precisely the way of seeing
that is determined by the warm relation of feeling to the object and consistently carried through, and
is therefore by its nature one-sided\ldots. The individual artist and the individual work of art can
never get further than to say something definite about their object according to the subject's relation
to it; the warmer, fresher, and more finely to the point, the better.'' The same grounding that holds
for pictorial art holds also for the art of life. All art rests on the principle of personal truth, and
personal truth is present only where the individual's whole distinctiveness stamps the work, whether the
work is an outer picture, or is the shaping of one's own life of feeling and will, the work of art that
every human being is called to produce. Afterwards the comparing and ordering thought can set to work,
draw up its catalogues, undertake its reductions and combinations, set up its types and its formulas. But
this is a work of the second rank. All dogmatics and all criticism are in the religious domain merely
secondary. When a definite religious idea is fixed, it can only be in order to obtain form, expression,
explanation for the immediate experience, and here it can always be asked with what right precisely this
idea is laid down, and with what right it stands as other and more than as the expression of the
individual's experience; criticism and dogmatics can here at once come into conflict. Only by virtue of
the absolute principle of authority would certain forms of ideas be able to be declared universally valid
and necessary expressions of all religious experience.

Religious-psychologically regarded, Schleiermacher hit the nail on the head when he said: ``Immediately,
everything in religion is true; for how else should it have come into being? But immediate is only that
which has not yet passed through the concept, but has grown forth purely in
feeling.''\textsuperscript{92}

Dogmatics is the index to life's book, not the book itself.

It is a Nemesis that very often strikes religion, that it thinks it is fighting for the life where it is
only fighting for the index.

And criticism has often imitated it in this, by thinking that it has done away with religion when it has
established the index's untenability. No index can step into the book's place, and only in the book
itself, not in any index whatever, do all the nuances and distinctive features come forth in their natural
truth.

\medskip

\noindent\textbf{83.} Historically we also find a growing acknowledgement of the principle of personality
in the religious domain.

The penal laws withdraw more and more from this domain. Religious crimes play a smaller role in Greek and
Roman legislation than in Jewish law, and likewise a smaller role in the Christian era than in
Greco-Roman antiquity, and in modern times than in the Middle Ages. According to the ancient and the
medieval conception, the domains of Church and State wholly or partly coincided, and he who separated
himself from the Church became an enemy of the State. The legal and ethical acknowledgement of the
principle of personality has exercised its influence on religion; but this would hardly have been possible
if within religion itself a development had not taken place in the direction of inwardness and personal
truth. Such a development can be traced in Indian religion by comparison between Brahmanism and Buddhism,
in the development of Greek religiosity from the time of Aeschylus and Socrates, within Christianity by
comparison between the Catholic world and the Protestant. In the most recent times there is here especially
to be noted the ever more frequent reference to the significance of the great personalities as founders and
models. The model steps more and more into the dogma's place. And a model one follows best by oneself
becoming a personality in one's own way, as the model was one in his.

% DROPPED IN THE 3RD EDITION. The 1901/1906 text had a subsection here,
% "Hovedgrupper af personlige Forskelligheder" (Main Groups of Personal
% Differences). It is absent from the 3rd ed., which runs straight from
% "Personlighedsprincipets Betydning og Berettigelse" to "Buddha og Jesus" —
% this is the "Sammenarbejdelse af de to ... Grupperinger af
% religionspsykologiske Typer" announced in the 1924 Fortale. Retained as a
% comment so the divergence from the first edition stays on the record.

\subsection{Buddha and Jesus}
\label{ssec:buddha-jesus}

\noindent\textbf{84.} It serves for the illumination and confirmation of the principle of
personality that at the source of the two highest popular religions there stand two great
personalities, each with its own distinctive character. Buddha and Jesus divide the world between
them. An attempt at a psychological characterization of them will then in this connection be
natural, however difficult it may be.

The difficulty rests especially on two circumstances.---In the first place, myth, legend, and dogma
early took possession of these figures. The very strong light that went out from their
personalities overwhelmed their contemporaries and immediate posterity and excluded a purely
objective historical account of them. It has gone with them as with the angel who in Dante's
``Purgatory'' shows the way upward, and who ``is hidden by the radiance he himself spreads.''---In
the second place, antiquity did not have the interest, which in the most recent times has come
forward so strongly, in following a personality's development step by step from its childhood to
its full maturity. The figures of antiquity stand before us fully finished, without our learning
what stages and what crises they have gone through to reach the fixed form in which they appear to
us.

I have in an earlier place (§34) given a characterization of Buddhism and Christianity. But it will
not be mere repetition when we here occupy ourselves with the founding personalities. There can
always be a difference between the source and the stream. At the source the original dispositions
appear; the stream is determined by many other circumstances.

\medskip

\noindent\textbf{85.} \emph{Buddha} was a prince's son who in the midst of life's splendour and
glory became attentive to suffering, infirmity, and death. This led him to brooding over life,
although he was not for his own person struck by life's shadow-sides. Neither by the road of
speculation nor by that of asceticism could he find peace. Consummation he reached only when he saw
through life as a great illusion begotten and nourished by thirst, drive, desire. From the need to
live spring sensuality, hatred, cruelty, besides all pain and all sorrow, and all this falls away
when the very need for life is abolished. There are, according to Buddha, five fetters that one must
free oneself from: selfishness, doubt, desire, asceticism (as an end), hatefulness. Their origin is
the same. If one can free oneself from them, there is attained the state of exalted peace called
Nirvana, which cannot be denoted by any positive predicate, because it stands in opposition to every
state that human experience knows. According to a legend Buddha speaks thus of Nirvana just before he
leaves his palace to become a hermit: ``When the fire of sensual desire is extinguished, then it is
Nirvana; when the fire of hatred and delusion is extinguished, then it is Nirvana; when pride, false
opinion, and all other passion and all anxiety are extinguished, then it is Nirvana.'' On a later
occasion he speaks thus to his disciples: ``Some ascetics and Brahmans direct false accusations
against me and say: a denier is he; he proclaims destruction, annihilation, the abolition of real
life!---What I am not and do not teach, that those dear ascetics and Brahmans falsely accuse me of
being and teaching. One thing only do I proclaim, now as before: suffering, and the extirpation of
suffering!'' Neither virtue nor knowledge was the highest for him, but ``the matterless, perfect
extinction of all illusions.'' All further metaphysical questions Buddha rejected, thinking that when
a human being has been struck by an arrow, what matters is to heal the wound, and it is a matter of
indifference who shot the arrow and what the bow and the string are made of.\textsuperscript{93}

What Buddha taught was in a certain sense nothing new. But there was accomplished in his mind and
through his experience of life a distinctive combination and condensation of tendencies that had
already stirred in the preceding religious and philosophical development in India. The concept of
Nirvana is found already in the Upanishads, on whose doctrines Buddhism is on the whole in the most
essential points dependent.\textsuperscript{94} It is an inwardization and a liberation, a practical
and concentrated living-into, that characterizes Buddha and leads him to produce a new form of life.
Neither cult nor myth, neither asceticism nor speculation now became the decisive thing, but the
personality's quiet, inward consummation through exaltedness above everything ``material,'' that is
to say, everything finite and changing. All outer differences become indifferent, including the caste
differences that played so great a role in India. Buddha did not come forward as a reformer of the
outer conditions. He did not wish to overthrow Brahmanism. But in his personality he furnished a model
of personal consummation that led out beyond everything existing. Access to the highest stood open to
everyone without regard to caste or sex. The struggle for peace of soul, for inner unity and freedom,
was for him the one thing, and for this struggle he awakened an enthusiasm that was consummated in
the joy of victory at having overcome the world.

Buddha's greatness is to have set up the ideal of personality as unity and freedom. Personality is not
absolutely tied to this or that, is no ``matter,'' but is a world of inwardness which on account of its
opposition to everything outer cannot be expressed in a positive manner.

He belongs, psychologically regarded, to the expansive, and disposition plays a greater role in him
than emotion. He is a reflective, not an intuitive nature, and in his reflection synthesis, which
abolishes all differences, plays a greater role than analysis. He is---in spite of the energy without
which Nirvana cannot be reached---a passive nature; or the activity is at any rate laid claim to by the
withdrawing process that is to make a completely passive relation to existence possible. He conceives
personality's unity in its relation of opposition to all the elements of personal life, not at the same
time as the totality within which they can find their place as organized members in a great connection.
The concept of personality is in Buddha like the lion's den, into which all tracks led, but out of which
none led. In the Buddhist collection of poems \emph{Dhammapada} (the Way of Truth) it is said, quite in
Buddha's spirit: ``He who is forbearing among the implacable, gentle among the violent, without desire
among the desirous, him I call a wise man.''\textsuperscript{95}

When Buddha had himself reached consummation, after the tempter (Mara, the Indian devil) had vainly
sought to hinder him in it, the latter tried a last means, urging him to keep the peace he had won for
himself and not to teach others how it could be reached. Characteristic is Buddha's answer: ``Whether
the Consummate One shows the disciples the truth or not, he is and remains the same. How is that
possible? Well, the Consummate One has denied and cut off at the root the illusion that defiles, and
that sows repeated existence, hatches pain, produces life, old age, and death!''\textsuperscript{96}---%
Thus he proclaims the truth because he has no ground for not proclaiming it; he stands entirely neutral
towards the possibility of proclaiming and of not proclaiming. Thus indeed must he stand who has no wish
any more. There therefore comes to be lacking a psychological foundation for the inward love of humanity
that the Buddhist movement bears the stamp of. The individual's consummation stands in decided
opposition to the wish for others' consummation. And yet fellow-feeling with human beings must have been
a motive for Buddha's break with his princely life. He had not himself been struck by the suffering he
saw around him; it was life's general lot that set his feeling in motion. It must therefore be due to an
imperfect accounting when he makes non-hatred one with love: hatefulness was indeed one of the five
fetters that were to be burst. It is not enough to burst the fetter; it also depends on the use of the
free mobility. That the barriers between human beings are broken down is not enough for positive
connection. Because the partition between two masses of ice is removed, they do not melt together; the
warmth must be added. This warmth has certainly been present already in original Buddhism; but the
anxiety about life's unrest, the critical reflection, and the bent towards passivity have stirred so
strongly that the sympathetic motives have not come forth to full consciousness. Love is once for all
akin to life and hope and is impossible unless the wish's right is acknowledged.---A motive that has not
least worked to inhibit positive and active fellow-feeling in Buddha and his first disciples is the
thought that when one attaches one's love to an object, the consequence is a greater possibility of pain
and sorrow than when the interest embraces only one's own self and its development. For a father who
lamented the loss of his son, Buddha is said to have had only this consolation: ``Yes, so it is,
householder. What one holds dear brings woe and lamentation, suffering, grief, and despair.'' Thus
feeling must be limited so that it shall not assume the form of pain. The thought that through the sorrow
and the loss that have love as their presupposition a higher spiritual development can be reached than
when love is shunned or limited seems not to have come forward in Buddha. That father went away offended
from the Indian sage, and he was right that a ``consummation'' won in this way cannot express the highest
personal value.

Yet love of humanity is motivated by Buddha and his followers in a positive manner by being seen as a
consequence of the attained consummation---and so far in a manner that comes into conflict with the pure
neutrality Buddha acknowledges towards Mara. There are, on closer view, two motivations, the one purely
psychological, the other of a more metaphysical kind.---Love is like an outflowing or radiation of the
inner life to others. This conception agrees with Buddha's expansive nature; it was properly inconsistent
when Nirvana halted the expansion. In the introduction to the Jataka book, love of humanity (in the
English translation: good-will) is named as the ninth perfection, and it is depicted thus: ``As water
cleanses all in the same way, the evil as well as the righteous, from all manner of dust and dirt, and
fills them with refreshing coolness, so shalt thou refresh friend and foe with thy love; and when this
ninth perfection is won, a Buddha's wisdom shall be thine!''---But not merely through such an involuntary
outflowing is a certain connection brought about between one's own consummation and devotion to
others---also through the fact that the individual, when all the differences and barriers called forth by
illusion have fallen, recognizes himself in all other beings. In one of Buddha's discourses it is said of
him who has won true peace: ``In his loving, compassionate, joy-filled, unmoved mind he recognizes himself
in everything and irradiates the world in all directions, out of his great, deep, unlimited heart, cleansed
of anger and rancour.'' Here, then, it is stated that a positive recognition takes place, not merely a
negation of barriers, and the liberation from hatefulness here stands merely as the negative condition for
the positive expansion that the recognition (the knowledge of all beings' unity) makes
possible.\textsuperscript{97}

How self-consummation and love---being a whole for oneself and being a member of a greater whole---are to
be united, that is a great task on whose solution humanity is constantly working. In the solution Buddha
gave in the sixth century B.C., self-consummation had decidedly the upper hand. This brought it
about---in conjunction with his passive tendency and with reflection's preponderance over intuition---that
he properly founded not a congregation but a monastic order. A popular religion Buddhism became only
through the legends that were attached to the figure of Buddha, and through the forms of cult that were
taken up from the older Indian religion, as well as through a stronger emphasis on active love of
humanity. The Mahayana sect in particular has made Buddha the god of gods, whose birth as a human being is
a world-event with which all gods and demons are occupied. Clothed in a luxuriant mythological and
liturgical form, Buddhism has spread in East Asia. It has, as a modern Buddhist has said, ``made Asia
mild!'' But it has predominantly worked in a damping, lulling, restraining way, where it was not, as among
the Japanese, crossed and altered by an actively forward-striving racial temperament, and by the influence
of an earlier religion (the Shinto religion) which had particularly developed the feeling of individuality
and nationality.\textsuperscript{98}

\medskip

\noindent\textbf{86.} \emph{Jesus of Nazareth,} the carpenter's son, lived in the midst of his people,
shared its memories and its hopes, wished and suffered with it. He did not, like the Indian prince's son,
look down on life from a higher standpoint as an observer, but stood in the midst of life, in the midst of
the yearning world. He wished to purify, to idealize the wish, not to abolish it. And even though the wish
and the struggle for the wish's realization entail pain, he will not give it up; pain is not always merely
an effect whose cause is to be abolished; it can be a cause of the testing of the powers, a means of
purification, a way to loftiness. His faith in the persistence of value is carried through by the valuable's
being won and preserved through struggle and suffering, not by the whole sphere where wish, struggle, and
suffering belong being unveiled and abolished as a great illusion. The contrasts here come into their own;
Jesus is not so much the man of expansion as the man of struggling oppositions. And he is the man of will,
does not withdraw to rest, but will determine his people's fate by leading it into higher paths. He
is---likewise in contrast to Buddha---the man of emotion. Even the Johannine Jesus must be characterized
thus in comparison with Buddha. He can be violently moved in spirit, shaken in his inmost being; his soul
can be terrified and sorrowful unto death, and he can be angered. Jesus's love is not merely mild,
involuntary radiation, but is the restlessly seeking and struggling love, is a bearing, not merely an
enduring devotion. There is here no fear that love shall bring sorrow by our weal and woe being tied to a
far greater circle of interests than our own isolated self-assertion embraces. There is an inner trust in
love's power and a feeling of its capacity to widen the mind, which does not let constricting misgivings
come forward. And love strives to be consummated in a higher righteousness through the idea of a kingdom
within which each individual can find his place.

It is another race and another culture that we here stand before than that in which Buddha appeared.
Besides, five hundred years lie between them, and the Jewish prophet built on far more agitated and
composite historical conditions than Buddha. But common to them is that their greatness rests on no single
new idea or new institution, but on the wonderful concentration with which they gather up the most
significant things in their people's development---the inwardness and depth with which the thoughts of the
past are shaped by them, the gripping mood with which their personalities are pervaded, and which streams
out from them over the world.

Jesus of Nazareth is the man of intuition, of the image, of the vision of the future. In his reflection he
is an analyst, asserts the significant differences, sets barriers and points to oppositions. He will have
us go through, not around, life's greatest contrasts.

Buddha rejects all images and thoughts as insufficient to express the highest; herein his intellectual
energy shows itself. Jesus has here no misgivings---but then neither does he show us where the limit of the
figurative lies, and thereby there were called forth among his followers the great doubts and the great
struggles. How far does the symbolic extend?---becomes the constant question. It is conditioned by the
founder's intellectual character that this question stands unresolved.

In the idea of a kingdom of God, as Jesus expressed it, there lay, as already said, the possibility that
individual personalities could reach the highest without the differences being effaced.---Yet there is a
point where Jesus's position towards history becomes analogous to Buddha's. For he does indeed lay great
weight on the development towards a future goal, but he maintains at the same time that the goal is not
reached in a positive manner through historical work, but that its realization takes place by a supernatural
breakthrough for which the human being is to hold himself ready. And since this breakthrough will come
speedily, the strained expectation, the constant watching and praying, becomes the essential attitude the
human being can take up towards it. What significance the constantly continued historical development and
the many-sided culture can then have becomes here at last just as great a problem as in Buddha.

But it was nevertheless of great importance that so great weight had been laid on life in expectation and
yearning. Jesus's visions of the future, the apocalyptic in his ideas, have through mighty images taught
human beings to look towards great ends that are to be reached through time, not by leaping over time.
Through transformations and displacements this contribution has been preserved for the race's continued
life, even after the narrow frame in which the contribution originally arose was dissolved. The struggling
human will has been able to find its symbols in Jesus's great images. Such a transformation and displacement
was, on the other hand, not so easily possible as regards Buddha's ideas; he brought quietives, not motives,
and the positive influence on spiritual life and culture had therefore to be far more limited. Buddha's
thoughts are like the grains of wheat that are found lying undissolved, but unfruitful, to this day in
Egypt's graves, as they were laid down millennia ago. Jesus's thoughts have shown their fruitfulness by
perishing in their original form and through the dissolution making their way to growing and working under
new conditions and through a series of historical displacements. Buddha has ``made Asia mild,'' but Jesus
has taught Europe a great \textit{Excelsior}.

\subsection{Is the Principle of Personality a Principle of
  Growth or of Dissolution?}
\label{ssec:pers-growth}

\noindent\textbf{87.} Growth and dissolution do not exclude one another, for the grain must be
dissolved in order to be able to sprout; but there might nevertheless be processes of dissolution
that did not bring new growth with them. It is often asserted that the principle of personality is
purely negative and dissolving. It puts, it seems, differences in the place of unity and
community. It isolates individuals. And it seems to conflict with every world-view that seeks back
to unity behind all oppositions, as all thought and all faith more or less have a tendency to do.

For the present the dissolving effect of the principle of personality will perhaps also be
predominant. But there is also a great deal that shall and must be dissolved in the religious domain
if this is to have a future at all. It is especially the hierarchical, tradition-governed character
that the religions have hitherto predominantly had. There is still both direct and indirect coercion
of conscience that must be got rid of, and this must for the present be at the expense of community
in mood and idea. Many will suffer under it, not merely those who cannot stand on their own feet,
nor merely the ``ecclesiastical natures,'' but also those who have a need, in spite of all
differences, to feel themselves one with others. Besides, the principle of personality, like all
great thoughts, will be able to be misunderstood and misused, so that one thinks to satisfy it by
withdrawing oneself from the influence of all models, from all instruction by tradition and others'
experience, and thereby only succeeds in becoming ``a fool on one's own account.''

In contrast to Catholicism, Protestantism emphasizes the principle of personality. It is both its
weakness and its strength. It is its weakness, because it cannot come forward with the certainty of
authority that the Catholic Church has. Thereby it comes to lack pedagogical authority towards
individuals, and it is short of closedness in the struggle against Catholicism's mighty, more and
more centralized organization. But it will more and more become its strength, if freedom and the
cause of personal truth have the future before them---and if it makes more earnest than hitherto of
carrying through the principle to which it owes its coming-into-being. Catholicism has in the
nineteenth century been able to do what Protestantism has never been able and never will be able to
do: to produce new dogmas. But if Protestantism could produce new personalities---would that not be
more valuable?---From the Catholic side the difference between orthodoxy and a freer religious
conception has recently been determined thus: ``Does truth exist outside believers? Does it
correspond to an objective reality? Does it impose itself from without? Has it immigrated from the
Beyond?---Or should it, on the contrary, consist in each individual's inner life as the fruit of
personal conscience, the resultant of individual religiosity, the expression and translation of
inward piety---should it in a word be subjective?\ldots\ Does religious truth come from God or is it
worked out in each of us? In the first case it \emph{is}, in the second case it \emph{comes to
be.}''\textsuperscript{99} To this it is to be said that if religious truth rests on experience, and
is won only by interpretation of this, then it comes to be, and constantly comes to be in everyone.
But it can for all that very well ``correspond to a reality,'' namely to the relation between value
and reality. What other relation does ``religious truth'' express than precisely this? And where does
this relation manifest itself except in individual personalities? If it is regarded as extended
beyond the limits of human life, that can only happen by the road of analogy and poetry. So long as
personal life persists, new relations between value and reality, and thus new religious truths, will
be able to appear. And if by ``God'' one understands something that is not merely ``outside'' us but
works in all reality, in all values, and thus also in the very relation between value and reality and
in the personalities that make experiences of this relation, then it is a false dilemma that is set
up when it is said that \emph{either} ``truth comes from God'' \emph{or else} ``it is worked out in
each of us.''

Protestantism is, as A. Sabatier has aptly said, a method. And the method can be no other than the one
that is given with the principle of personality, with the acknowledgement of individual personalities
as centres of value and of experience. The Reformers to be sure did not see this. They dissolved the
old church community, but a consistently practicable principle for a new community they did not find.
Hence the unrest, the strife, and the mobility within the Protestant world.

Catholicism stands related to Protestantism as Ptolemaic astronomy to Copernican. But just as little as
Copernicus saw that once he had moved the earth out from the world's mid-point, the consequence would
be that no world mid-point at all could be pointed out, just as little did the Reformers see that when
they had thrown the Pope down from authority's high seat, no absolute authority in the religious domain
could historically be pointed out any more. If one wants authority, it is most consistent to give up
Protestantism and seek back to the Church that has most completely carried through the principle of
authority.

\medskip

\noindent\textbf{88.} We live quite certainly in what has been called a ``critical'' age, not in an
``organic'' one. But it is not therefore merely dissolving forces that are at work. There is nothing
that prevents the joining together in smaller groups of personalities about a common direction of
thought or spirit, or about a common symbol. And such joining together can often become more inward and
freer than where traditional authority is the cohesive force. Besides, the principle of personality
itself can stand as the expression of a great truth, one of the highest spiritual values. Whatever a
human being's faith is or becomes---the fact that he lives in it with his whole soul, and that his
distinctive character unfolds through the discovery and appropriation of what he believes in, gives his
faith a value that the best guaranteed system of doctrine would not have. Here is something in which
all will be able to come to mutual understanding, however far apart they may otherwise stand from one
another as regards the content of their faith. When responsiveness to personal distinctive features
grows, one will less and less sacrifice the personal accent for agreement in positive or negative
dogmas. Here we stand before an extension of the spiritual world that is quite as significant as the
extension of the material world became in its time. All culture's finest flower arises through
sympathetic understanding of others' personal distinctiveness, from which it perhaps follows precisely
that they look at essential questions quite otherwise than we ourselves do. Only a few steps have
hitherto been taken on this road. But the principle of personality is a positive and fruitful principle,
because it proclaims that it is this road we must go, and because it thereby opens the possibility of a
deeper feeling of community than the one conditioned by joining together about common dogmas.

Philosophically regarded, the principle of personality testifies, the more it is confirmed in the
particular, to existence's richness and fullness. The history of philosophy shows us that it has mostly
been acknowledged and valued by thinkers who laid great weight on the thought of unity (Montaigne,
Spinoza, Fichte). Existence's principle of unity must be the more powerful, and its working have a deeper
and more inward character, the more and the stronger the differences that individual beings present. It
is easy enough to believe in a unity when one effaces or disdains the differences. With religious
devotion he whose faith in unity has a less abstract character discovers how differently the universal
life can shape itself in individual beings. It does not shake his faith; it confirms it. His doubt is
directed against the superficial faith in a unity that conceives this either quite abstractly or quite
externally; but this doubt springs precisely from a still deeper faith (\textit{alte dubitat, qvi altius
credit}).---Ethically regarded something similar holds. The kingdom of personalities on which all ethics
spells out or builds will become the richer, the more distinctive personalities in mutual harmony it
encloses. The standard of a human community's value consists not merely in the unity and law that
prevail, but also in the fullness of distinctive features, of different centres of value and of
experience, that it is able to connect.---Purely psychologically regarded, too, something similar holds.
The validity of the principle of personality sets tasks for a comparative psychology of greater extent
and depth than psychology's older frames gave room for. This holds particularly as regards the psychology
of religion. The history of religion has hitherto occupied itself predominantly with the great common
views in the religious domain. There is quite as much to learn from the history of individual views of
life; and since we are each of us ``individuals,'' not human beings in general, this last instruction,
when it can one day properly be given, will become the most valuable.---

A principle that opens such prospects and sets such tasks cannot be merely dissolving. It is one of the
most fruitful principles that has ever worked its way forward.


\subsection{Layman and Scholar}
\label{ssec:layman-scholar}

\noindent\textbf{89.} By virtue of the principle of personality, each individual's religious conviction
must be formed according to his experience of values, of reality, and of the relation between them. A
spiritual and universal priesthood holds. But one develops best towards independence and self-activity
in reciprocal relation not merely with actual life, but also with other personalities who take life in an
independent manner and have access to knowing both the world of values and the world of reality in a more
thorough and comprehensive manner than the individual at the beginning of his road is able to. The
principle of personality therefore does not exclude models and teachers. The difference between layman and
scholar does not fall away. On the contrary, the more the principle of personality penetrates, the better
and better will the laity understand that they need the learned, and they will understand better and
better how to use what they can learn from these in the right way. They will see the great significance of
the learned's work in enlightening us about existence's great connection and about the riches it encloses,
and in putting clear knowledge in the place of dreams and dogmas. And yet they will know that the final
conclusion of thoughts about life's riddles can only be won on the basis of one's own experience and faith,
neither from new enlightenment nor from old dogmatics. And it will be an enrichment of spiritual life when
each individual, more or less consciously, makes up his reckoning on the basis of his experiences and his
intellectual development. Scientifically regarded, personality itself is the last, perhaps insoluble riddle,
the ever distant concluding point: scientific thinking is itself a spiritual activity that can only be
exercised by a personality, and the last riddle would not be solved because science explained everything
else, if it did not explain its own last presupposition. In my \emph{Psychology} I have sought to illuminate
this more closely, and my result I can express in the proposition: ``The thought that explains everything
becomes its own last---and constant---problem.'' But in life personality is the first thing, that which bears
everything, science included, and sets its seal on everything. From life's involuntary course, in which
personality constantly reveals itself from new sides, to the last results of scientific thinking, there is a
long series of transitions; but when both life and science understand themselves rightly, there is no breach
or leap at any point.

\medskip

\noindent\textbf{90.} The matter presents itself otherwise at the standpoint of the principle of authority.
After the time in the Christian Church when revelation was immediately continued through ``the proof of the
Spirit and of power'' in individuals' enthusiasm had ceased, personal experiences were dammed in by
ecclesiastical authority. The difference between layman and scholar soon became greater in religious
questions than in scientific ones. Whether the Church referred the layman to the infallible Pope, who was
elected by the Church, or to an infallible book, which was due to the Church's choosing and rejecting, he had
at any rate to take up a waiting attitude. His faith became a charcoal-burner's faith.

And yet there came a time when the Church's own men felt themselves to be laymen. It was when there arose a
question of the validity of tradition, and especially of the Bible's authenticity. The principle of authority
leads indeed quite naturally to studying the history of the authorities. But who is now to decide which
historical conception of authority's development is the true one? Scientifically the matter is not in doubt.
The coming-into-being and development of all traditions and all books can scientifically only be known by the
help of philology and historical criticism. Since now priests and theologians are often as little philologists
and historians as laymen as a rule are, authority would then at last pass from the ecclesiastical authorities'
hands into those of the learned. That one would not hear of. At the Council of Trent it was laid down not
merely that the Bible's authority was grounded on the Church's, and that it was the Church's business to decide
what was the right tradition, but that the Latin translation of the Bible (the Vulgate), which had for centuries
been in use in the Church, should be the foundation of all preaching and all discussion. One did not wish
``grammarians and pedants'' to stand above bishops and theologians, or that only he who understood Greek or
Hebrew should be able to sit in judgment on a heretic.\textsuperscript{100}---In the Protestant world an
analogous situation has come forward and assumes especially in the most recent time very glaring forms.
Protestant orthodoxy led to a rule of theologians, since only the theologians versed in the Bible could decide
what the right doctrine was. The laymen's reaction against this lay to a great extent at the foundation of the
spread that Pietism and Rationalism obtained; one wished to overthrow the many little popes that had come in
place of the one great one. But when Rationalism had been organized into a system of doctrine, there arose
against it a new lay movement. One found in many circles a nutritive value in the old writings and traditions
which rationalist theology did not let come into its own, and besides, Rationalism demanded a reinterpretation
of Bible and tradition and a critical attitude towards them, while the laymen preferred to hold to the holy
words as they stood. Literalism now became a democratic principle. It was then again, as at Trent, maintained
(with us especially by Grundtvig) that ``the faith of the unlearned could not be bound to the testimony of the
learned.''\textsuperscript{101}

But the learned do not cease to inquire, and especially not to investigate the biblical writings' inner and
outer connection with the age in which they came into being. To this comes the fact that the religious ideas
are discussed from general philosophical points of view, which with greater or lesser consistency also make
themselves felt in theology. There has therefore developed an opposition between the Church and theology which
is one of the most remarkable signs of the times in our days. Theology, which is often regarded as science's
antipode, comes more and more to stand as one of science's outposts. Or---as a modern theologian has put
it---it stands as a buffer between the Church and scientific culture---a buffer at which blows are struck from
both sides.

The principle of authority thus makes front on two sides: partly against the layman's free, personal life, which
is to be bent into the ecclesiastical forms, partly against science, when this would test these forms' origin
and value. But the situation's distinctive character arises especially from the great adherence that the
principle of authority has won in the course of the nineteenth century in lay circles. It has shown itself to be
a democratic principle: it procures foothold, security, and support for those whose inner and outer conditions do
not permit independent seeking with its labours, dangers, and crises. Modern Catholicism's firm hold on the great
masses appears more and more clearly. And in the Protestant Church it is more and more laymen and lay preachers
who dominate religious life. The control of the priests' orthodoxy is exercised in Protestantism far more by the
laymen than by the ecclesiastical authorities. This influence extends to theology too, in that the priests, who
are themselves controlled by the laymen, demand influence on the filling of theological professorships. The
Church can, as has been said, neither endure nor do without theology.\textsuperscript{102}

\medskip

\noindent\textbf{91.} When the Church steps between individuals' personal life and science and would keep them
apart from one another, it violates two spiritual powers at once. That does not happen unpunished.

The intellectual radiance that formerly surrounded the ecclesiastical eminences, and through them ecclesiastical
life, has vanished. Really seriously carried through religious thinking becomes rarer and rarer. One asks only
anxiously whether the traditional dogmas are satisfied. In the Catholic Church one lives on the thirteenth
century, in the Protestant Church on the seventeenth. Alongside tradition's dominion there is indeed another trait
of a more praiseworthy character that is distinctive of the Church in our days, namely the great philanthropic work
it has organized. It is the type of St.~Vincent de Paul that comes back again. But the advantages philanthropy may
have from no time and strength going to thinking will in the long run hardly be able to outweigh the harm spiritual
life suffers from independence being inhibited by the anxiety about free discussion of life's most important
questions.

The two spiritual powers that the Church would keep apart from one another will nevertheless find each other. On
closer view, the modern development in the Church, characterized by the conscious acknowledgement of the principle
of authority in wider circles than before, is after all in agreement with the principle of personality. The lay
movement may work inhibitingly and inimically to science for a time; but in itself it is nevertheless a sign that
sleeping literalism has been succeeded by living literalism, and between these it is, according to the principle of
personality, not hard to choose. Through the stronger weight that is laid on getting each individual to join in,
personal life will be able to awaken in far more people than hitherto; and once it has awakened, it will not fail
to work transformingly on the old forms and to seek out beyond them. The most significant advances in the spiritual
domain often come from quarters from which one would least expect it. The new values very often come not from the
world of criticism and analysis, but from circles where one has lived with depth and inwardness in the old values.
History does not always go the straight ways.

\begin{center}---\end{center}

%% ============================================================
%%  IV. ETHICAL PHILOSOPHY OF RELIGION
%% ============================================================
\chapter{Ethical Philosophy of Religion}
\label{ch:ethics}

\section{Religion as the Foundation of Ethics}
\label{sec:foundation-ethics}

% [text to be added]

\section{Religion as a Form of Spiritual Culture}
\label{sec:spiritual-culture}

\subsection{Psychological Considerations}
\label{ssec:psych-considerations}

% [text to be added]

\subsection{Sociological Considerations}
\label{ssec:sociol-considerations}

% [text to be added]

\section{Primitive Christianity and Modern Christianity}
\label{sec:primitive-christianity}

% [text to be added]

\section{We Live on Realities}
\label{sec:realities}

% [text to be added]

%% ============================================================
%%  NOTES
%% ============================================================
\chapter*{Notes}
\addcontentsline{toc}{chapter}{Notes}
\label{ch:notes}

% [notes to be added]

\end{document}
